% v27.1 (2026-10-16): source attribution of the Goldman Sachs note removed (Section 16, references, acknowledgments); no other change.
% ai_compute_limits_v26_2026-09-30.tex
% v27 (2026-10-15 build): ORCID on title page; date line set to release.
% v27: Goldman Sachs Phase-2 ROIC note (Sheridan et al., Sept 24 2026)
% reconciled with the carry in a new Section; four first-round arguments recorded; observables,
% register anchors (c_opex, grid cap), glossary (NOPAT, Phase 2, take-or-pay) and references added.
% v26.7: byline - Retired, IBM Systems / Founder, Nvara.ai; nvara.ai email added.
% v26.6: acknowledgments - names only: Kountouris, Pedraza, M. Fitzpatrick, Mehta.
% v26.5: a16z State of Markets (Sept 2026) taken point by point in a new section; arms-race
% path relabelled as street consensus; observables (consensus capex, backlogs ~$1.7T, GPU rental
% prices, consensus FCF, OpenRouter shares); electricity-bill study; labor-budget objection named.
% v26.4: Copilot paid-vs-used qualification (E7 bundling, usage billing, active share as
% secondary variable); Big-5 vs Chinese-four alignment subsection; historical productivity-lag
% data (David 1990, BLS, Jorgenson-Ho-Stiroh 2008, Brynjolfsson-Hitt 2000, Census BTOS).
% v26.3: symbols named at first use; Q2-2026 disclosures (Jul 22-30) added to observables;
% Copilot seats 30M and M365 seat growth +6-7% into the assistant chain and incumbent case;
% Google backlog $514B. Register rebuilt (minor).
% v26.2: acknowledgments final (Pedraza, M. Fitzpatrick, Mehta); release statement.
% v26.1: citations added to the four new Section 14 subsections.
% v26 (pre-earnings, final): batch share + computing precedents; producer cost curve
% (the low-cost producers); September disclosures (Oracle Q1 FY27, Anthropic >$100B pace,
% OpenAI Q2); effective vs list token prices; register rebuilt (minor).
% v25 (pre-earnings release): reviewer feedback incorporated - economic vs accounting
% depreciation (secondary variable, unrecognized impairment), inference utilization,
% the commodity floor (marginal cash cost of inference), market context; acknowledgments;
% explicit commitment to a post-Q3 update (v26).
% v24: notation audit (K, M, W(e), T() shorthand, f_tr); 'Every Equation, Term by Term'
% appendix; change-marked PDF produced with latexdiff for review.
% v23: consistency pass; cannibalization/BofA reply; Microsoft separated as an
% illustration; Appendix G register softened + 'better guidance' framing; conflict
% disclosure; OpenAI figures treated skeptically; release formats stated.
% v22 (final candidate): glossary completed by symbol audit - derivation symbols,
% subscripted g's, ARPU/M365/DINUM/IDC/MIT NANDA/DOI/edge-silicon/HBM-supplier entries.
% v21: Appendix 'Every Input, Appraised' - each register entry with evidence, range,
% source bias, what would improve it, and what is likely.
% v20: glossary gains a complete table of model variables and parameters.
% v19: update protocol (register, rebuild script, disclosure calendar, re-issue rule).
% v18: incumbent's dilemma (Microsoft 365 worked case: price up, seats down,
% Copilot attach); reallocation share and seat deflation declared secondary
% variables with their own curves (f21, f22).
% v17: the enterprise software budget pool decomposed (seat-SaaS deflation, services,
% sovereignty displacement, reallocation to AI); Gartner discounted; Microsoft exposure.
% v16: response to author review of 9 Sep 2026.
% Horizon extended to 2035; adaptive capex path; revenue-flow diagram;
% units throughout; version narrative removed; slack-trap evidence added.
\documentclass[11pt,letterpaper]{article}
\usepackage[utf8]{inputenc}
\usepackage[margin=1.0in]{geometry}
\usepackage{amsmath,amssymb,amsfonts}
\usepackage{booktabs}
\usepackage{tabularx}
\usepackage[final,expansion=false]{microtype}
\usepackage{setspace}
\usepackage{cite}
\usepackage{array}
\usepackage{url}
\usepackage{xcolor}
\usepackage{graphicx}
\usepackage{abstract}
\usepackage[breaklinks=true,colorlinks=true,linkcolor=blue,citecolor=blue,urlcolor=blue]{hyperref}
\newcolumntype{L}[1]{>{\raggedright\arraybackslash}p{#1}}
\newcommand{\verify}{\textsuperscript{\textcolor{red}{[v]}}}
\setstretch{1.15}
\renewcommand{\abstractnamefont}{\normalfont\bfseries}
\renewcommand{\abstracttextfont}{\normalfont\small}

\title{\textbf{\Large The Macroeconomic Limits of AI Compute: \\ A Vintage Fermi Model of Capital Carry, \\ Sectoral Demand, and the Growth Rate the Build-Out Requires}}
\author{\textbf{Desmond Fitzpatrick, Ph.D.} \\
\small Retired, IBM Systems \\
\small Founder, Nvara.ai \\
\small \texttt{desmond.fitzpatrick@nvara.ai} \quad \texttt{desmond.fitzpatrick@gmail.com} \\
\small ORCID: \url{https://orcid.org/0009-0001-0534-2793}}
\date{15 October 2026 --- Version 27, release. \textit{The register is rebuilt after the Big-5 Q3 calls of 28--30 October and Nvidia's report of 19 November; a re-issue follows only if the change is material; see Appendix~\ref{app:update}.}}

\begin{document}
\maketitle

\begin{abstract}
The five largest US cloud and AI infrastructure operators spent \$448B on capital expenditure in 2025 and have guided to roughly \$780B for 2026~\cite{ref:epoch_capex,ref:mlq2026,ref:statista2026}. Whether that capital earns its carry is not a question about the size of AI demand today but about how fast it must grow. We build a vintage model: capex cohorts from 2023, net of the 2022 pre-AI baseline, each split into silicon (economic life $\sim$4 yr) and shell (life $\sim$20 yr) with an energization lag, and compute the annual carry --- depreciation, facility opex, and return on net book --- for each year 2027--2030. The carry ramps from $\sim$\$390B in 2027 to $\sim$\$725B in 2030 if capex merely plateaus, because the 2025--27 cohorts are still on the books as new ones arrive. Against calendar-2026 AI revenue of \$187B --- decomposed into seven sourced components (frontier-lab bookings, hyperscaler AI services net of reseller pass-through, AI-attributable advertising, consumer) with a P10--P90 of \$165--190B~\cite{ref:techcrunch_anthropic,ref:axios_anthropic,ref:epoch_revenue,ref:sacra_anthropic,ref:meta_q2_2026} --- we solve for the sustained growth rate $g^{*}$ at which the installed base breaks even: about 48\%/yr through 2030 on a capex plateau, 62\% under an arms race, 33\% under a 35\% capex cut. No plausible rate clears 2027--28; those years are a sunk deficit of \$100--250B/yr. A second, independent estimate --- fifteen sector Fermi chains, each with a 2026 base, a 2030 ceiling set by payroll, IT-budget or unit constraints, a ramp rate, and a cloud-capture share --- delivers end-customer revenue of $\sim$\$625B in 2030 (P10--P90 \$525--745B) with growth decaying from $\sim$57\% in 2027 to $\sim$14\% in 2030, a 33\% CAGR. Cloud receipts from those sectors are $\sim$\$490B against a 2030 carry of $\sim$\$725B: a deficit of \$255B/yr, 3.5 standard deviations of the propagated uncertainty, and \$190--540B/yr under an arms race. Nineteen inputs that carried most of the variance --- the 2026 base and its hyperscaler component, the cloud-capture share and its lab component, the advertising, assistant and consumer ceilings, facility opex, the non-AI capex baseline (constrained by 2023 filings), energization lags, silicon life by fleet, the silicon share from Nvidia's revenue, the cost of capital by CAPM, grid energization capacity, siting and consent deferral, water-driven cost and cooling uplifts, and the budget feasibility cutoff --- are each decomposed into sourced components, cutting the top-down root-sum-square from 286 to 104 and leaving the residual in two axes: an adoption/capability factor and the assistant ceiling. Over those two we map the probability that receipts exceed carry inside a feasibility band derived from Gartner enterprise software and services budgets. On a capex plateau the 50\% contour requires 2030 AI revenue of $\sim$\$950B --- inside the feasible band's upper half, requiring AI to take $\sim$15\% of all enterprise software and services spending by 2030; on a 35\% capex cut P(profit) reaches $\sim$50\% in 2030 and $\sim$95\% by 2032 at the paper's own modes. As of September 2026 the reported growth of frontier-lab revenue exceeds $g^{*}$ by a wide margin~\cite{ref:techcrunch_anthropic,ref:epoch_revenue}; the sector ceilings are the reason we expect it to decay. Five structural mechanisms --- slack absorption, verification, value leakage to integrators and on-device silicon, sovereignty carve-outs, and grid queues --- are what set those ceilings. If the ceilings are wrong by a factor of 1.5, the build-out pays. The paper is a probability assessment, not a forecast: at the register's current inputs the most likely outcome is the deficit above, and that most likely outcome is overturned if calendar-2028 AI revenue exceeds \$400B. Extending the horizon to 2035 does not rescue a plateau: the annual deficit narrows to zero only around 2035 and the cumulative hole is not recovered; a rational operator who cuts capex in step with realized receipts turns profitable in 2031 and recovers cumulatively by 2034. Two reviewer points sharpen the supply side: if an accelerator's value tracks the marginal cost of inference on the newest chip --- falling 50--90\% a year --- the silicon bought through 2026 is already worth \$100--250B less than straight-line books imply; and if token prices converge on the marginal cash cost of inference ($\sim$\$0.03--0.05 per million tokens, a tenth of realized prices), receipts converge on the opex line and the carry is never recovered. Goldman Sachs's September 2026 ROIC framework, restated per dollar of capital, requires 27 cents of annual revenue per dollar against the paper's 28; the two differ on scope --- one cohort or the stacked stock --- and on one input, facility operating cost. Model, data and code are open and driven from a register of inputs; every input is appraised --- evidence, range, source bias, what would improve it --- in Appendix~\ref{app:inputs}, with an update protocol in Appendix~\ref{app:update}; a glossary of terms and phrases is in Appendix~\ref{app:glossary}.
\end{abstract}

\vspace{0.5em}
\noindent\textbf{Keywords:} Fermi estimation, vintage capital model, hyperscale capex, break-even growth, productivity J-curve, surplus capture, ASC 360-10.
\vspace{0.5em}
\hrule
\vspace{1em}

\section*{Executive Summary}
\begin{itemize}
    \item \textbf{The question is $g$, not $D$.} A stock built ahead of demand is rational if demand compounds fast enough. Comparing a stock to a demand level lets each side pick its year. The vintage model instead asks: at what sustained growth rate does the installed base earn its carry, and by when? (Figures~\ref{fig:scissors} and~\ref{fig:breakeven}.)
    \item \textbf{The carry ramps.} Net of a non-AI baseline that the flat 2023 capex total pins at $\sim$\$130B for that year, growing 2--10\%/yr since, the AI-attributable stock is $\sim$\$1.1T entering 2027 and $\sim$\$2.2T by 2029. Its carry is $\sim$\$390B in 2027, $\sim$\$535B in 2028, $\sim$\$725B in 2030 on a plateau (Figure~\ref{fig:carry}) --- siting friction defers some of it a year and then adds cost. A 35\% capex cut flattens it near \$460B; an arms race pushes it to \$1T.
    \item \textbf{Break-even growth is 49--62\%/yr sustained to 2030.} From a calendar-2026 base of $\sim$\$179B, plateau needs $g^{*}\approx49\%$ (2030) and $\sim$64\% to recover the 2027--28 deficits cumulatively; rationalization needs $\sim$34\%; arms race $\sim$62\%. Across the decomposed base's P10--P90 (\$163--194B) the plateau figure runs 53\% to 46\%.
    \item \textbf{The sectors deliver about 33\%, not 49\%.} Fifteen independent sector chains (Section~\ref{sec:bottomup}, Figure~\ref{fig:sectors}) put 2030 end-customer AI revenue at $\sim$\$625B and infrastructure receipts at $\sim$\$480B --- office assistants and agents ($\sim$\$135B), advertising uplift ($\sim$\$118B), consumer ($\sim$\$60B) and coding tools ($\sim$\$60B) dominate; everything else is under \$50B each. Growth decays as sectors approach seat-, budget-, user- and yield-set ceilings. The 2030 shortfall against a plateau carry is $\sim$\$260B/yr, a 3.5$\sigma$ result (Figure~\ref{fig:reconcile}).
    \item \textbf{The observed growth rate is above the threshold today, unevenly.} Anthropic's reported run-rate went from \$9B (end-2025) to \$65B (July 2026) and OpenAI's from \$21B to \$40B~\cite{ref:techcrunch_anthropic,ref:axios_anthropic,ref:epoch_revenue}; Microsoft's consumer assistant and Meta's non-advertising AI products have no comparable disclosures. Two years of the labs' pace would clear the plateau path by 2029. The sector chains below are the reasons the paper expects deceleration; whether it comes is the 2027--28 test (Section~\ref{sec:falsify}).
    \item \textbf{Where the remaining doubt lives, and whether profit can overlap.} After nineteen decompositions (Section~\ref{sec:decomp}) the residual uncertainty sits in two axes: an adoption/capability factor and the assistant ceiling. Figure~\ref{fig:probmap} maps P(receipts $>$ carry) over both, with a feasibility band built from the enterprise software and services budget --- itself decomposed, because that budget is being torn up: seat-priced SaaS is deflating under agents, European and Chinese buyers are displacing US vendors, and the money is moving from seats to AI solutions. On a plateau the 50\% contour sits at a residual factor of $\sim$1.5 --- 2030 AI revenue of $\sim$\$950B, which is inside the feasible band (\$920B--\$1.31T) and requires AI to take $\sim$25\% of a pool that is not growing. Under a 35\% capex cut P(profit) is $\sim$50\% in 2030 and $\sim$95\% by 2032 at the paper's modes. That is the paper's result stated as a probability rather than a sign: not a prediction of loss, but a distribution over outcomes and the growth it takes to move it. 
    \item \textbf{Economic depreciation and the commodity floor.} Straight-line lives are the accounting view. The economic view (reviewer comment) is that an installed accelerator is worth what it can earn against the newest chip's marginal cost per token, which falls 50--90\% a year; on that schedule the 2023--26 silicon carries an unrecognized impairment of \$100--250B today, and the carry is front-loaded into years the paper already calls a sunk deficit. Doubling inference utilization (batching, quantization, open models) serves twice the tokens from the same stock but, if prices fall in step, no more dollars. And if prices converge on marginal cash cost --- $\sim$\$0.03--0.05 per million tokens against realized prices of \$0.10--1.0 --- receipts converge on the opex line, $\sim$\$130B against a \$725B carry. Section~\ref{sec:commodity}; the metric to watch is realized price per million tokens against that floor.
    \item \textbf{Readiness is a constraint in its own right.} Data Center Watch counts \$130B of US projects blocked or delayed in Q1 2026 alone --- matching all of 2025 --- with 833 opposition groups in 49 states, a New York moratorium, and Ipsos polling near 70\% local opposition~\cite{ref:dcw_q1_2026,ref:forbes_consent,ref:tomshw_dcw}. Water, ratepayer bills, noise and diesel are the stated grounds; Dublin and Amsterdam were the European template. The model now defers 10--35\% of each cohort a year for consent and siting, adds 2--9\% to shell cost for community benefits, water mitigation and relocation, and 0--6\% to cooling opex where evaporative cooling is refused. The net effect is to shift the carry later, not to remove it.
    \item \textbf{Why $g$ decays.} Nominal task savings convert to cash at roughly 0.28$\times$ (fractional time, verification, Parkinson--Jevons); realized savings leak to buyers and outcome-priced integrators at 10--50:1 against cloud tokens (Nordhaus); routine inference migrates to 1.45B NPUs a year; sovereignty removes 15--35\% of demand while duplicating capital; grid queues strand cohorts before they earn; and a material share of 2026 lab compute is paid with investor capital (\$122B and \$65B raised this year~\cite{ref:valueadd_openai,ref:techcrunch_anthropic}), not customer revenue.
    \item \textbf{Scale.} A \$780B/yr build-out and a \$255B/yr shortfall sound enormous and are not, against the economies that host them: world GDP is $\sim$\$115T, the US $\sim$\$30T, China $\sim$\$19T, the EU $\sim$\$20T\verify. The build-out is 1.5--2.5\% of US GDP per year depending on definition (gross, or lease-adjusted as in BofA's cycle comparison~\cite{ref:a16z_som}) --- at the peaks of the railway booms of the 1840s and 1870s and of telecoms in 1999--2001, each of which ended in impairment without ending the technology.
    \item \textbf{The street's path is the arms race.} September 2026 consensus has Big-5 capex at $\sim$\$1.05T in 2027 and $\sim$\$1.15T in 2028~\cite{ref:a16z_som}. The paper's plateau is therefore below consensus and its arms-race path is consensus; on that path the 2030 deficit is $\sim$\$545B and no year through 2035 is profitable. The strongest statement of the bull case, a16z's September deck, is taken point by point in Section~\ref{sec:a16z}; the two documents agree on nearly every observable and disagree on two judgments --- whether AI is paid from labor budgets as well as IT budgets, and whether today's rising GPU rental prices are a permanent state or a cycle top. The question here is who carries the write-down and when, not whether AI is affordable.
    \item \textbf{The sell-side hurdle agrees with the carry.} Goldman Sachs's September 2026 framework~\cite{ref:gs_roic2026} --- \$1.42T of 2028--30 revenue to earn 15\% on the \$1.73T 2026--27 cohort --- is 27.4 cents per dollar of capital per year; the paper's carry is 28.4 (Section~\ref{sec:gs}, Table~\ref{tab:gs_reconcile}). The note isolates one cohort and benchmarks to backlog; the paper stacks the cohorts and forecasts receipts. The one material input difference is facility operating cost (2\% versus 6.65\% of capital), now carried as two anchors in the register.
    \item \textbf{Beyond 2030.} Ending the analysis in 2030 would invite the objection that 2032 makes it all back. Appendix~\ref{app:horizon} runs the model to 2035: on a plateau the annual deficit narrows to roughly zero by 2035 and the cumulative deficit is never recovered within the horizon; under an arms race it widens every year; under a 35\% cut, or under an adaptive path in which capex tracks realized receipts, annual profit arrives in 2030--31 and cumulative recovery by 2034--35.
    \item \textbf{Consequence, either way.} The 2025--27 cohorts will not earn their carry in 2027--28 under any growth rate. ASC 360-10 recoverability~\cite{ref:fasb_asc360} is tested on projected cash flows, so impairment turns on whether management can still credibly forecast $g\ge g^{*}$ in 2028. That is the observable this paper puts its name on.
\end{itemize}
\newpage

\section{Method: What a Fermi Estimate Can and Cannot Say}
On July 16, 1945, ten miles from the Trinity tower, Enrico Fermi dropped scraps of paper as the blast wave passed and, from their $\sim$2.5\,m displacement, estimated a yield of about 10 kilotons of TNT~\cite{ref:fermi_trinity}; radiochemistry later gave $21\pm2$ kilotons~\cite{ref:katz2021}. He had the order of magnitude before the instruments were read. At Chicago he institutionalized the method with the piano-tuner question~\cite{ref:morrison1963,ref:vonbaeyer1993}: chain population, household size, ownership density, tuning frequency and tuner throughput, and the product lands within a factor of two of the directory count.

Why does this work? If the estimate is a product of $n$ independent factors, $Y=\prod X_i$, then $\ln Y=\sum\ln X_i$ and, for independent errors,
\begin{equation}
\sigma_{\ln Y}^{2}=\sum_{i=1}^{n}\sigma_{\ln X_i}^{2}.
\label{eq:var}
\end{equation}
Errors add in quadrature, and --- the point that matters in practice --- they add \emph{weighted}. Most factors in a well-built chain are known to $\pm$10--20\% ($\sigma_{\ln X}\approx0.1$--0.2); only one or two are loose. The total is then dominated by the loose ones and barely moved by the tight ones: nine factors at $\sigma=0.1$ contribute less than one factor at $\sigma=0.3$. Fermi's paper-scraps got within a factor of two because he had one loose factor and knew it. A corporate budget is not an order-of-magnitude exercise either: people know what they spend to within 5--20\%. This paper is in that regime. Every input below is bounded by a filing, a disclosure, a physical relation or a budget, and the reader should expect --- and hold the paper to --- precision of tens of percent on inputs and a factor of about 1.3 on the result, not powers of ten. ``Fermi'' here names the method of decomposition and propagation, not a claim that only the order of magnitude is known. Adding tight factors does not accumulate error; it partitions the estimate so that the loose factors are exposed and can be attacked individually. That is why this paper decomposes demand into fifteen sectors rather than one growth rate, and why Section~\ref{sec:results} reports a variance budget: the reader should see which two or three inputs carry the uncertainty, because those are the only ones worth arguing about. Every numbered equation in the paper is restated and explained term by term, with a worked number, in Appendix~\ref{app:equations}. The full propagation uses three different rules, set out and checked numerically in Appendix~\ref{app:errors}: (1)~\emph{within a chain}, where factors multiply, relative uncertainties add in quadrature ($\sigma^2_{\ln}$ add); (2)~\emph{across chains that are summed} --- sectors, or capex cohorts --- absolute uncertainties add in quadrature, which in relative terms weights each chain by the square of its share of the total, plus any factor the chains have in common, which adds once and does not average away; (3)~\emph{between supply and demand}, which are subtracted, absolute uncertainties again add in quadrature and the result is judged by the gap divided by that combined uncertainty (the $z$-score).

Two disciplines follow: report one or two significant figures, and call a conclusion robust only when the gap exceeds the propagated uncertainty. Where the gap is inside the noise, say so.

\noindent\textbf{Open model.} Parameters, distributions, the Epoch AI capex series and the Monte Carlo code are archived with this paper on Zenodo, the CERN-hosted open research repository; the deposit's permanent digital object identifier (DOI) is \url{https://doi.org/10.5281/zenodo.23406476}; later versions are reachable from the same record.

\section{Prior Literature}
\begin{enumerate}
    \item \textbf{The capex--revenue gap.} Cahn~\cite{ref:cahn2024} framed the ``\$600B question'' as a flow comparison; Covello~\cite{ref:covello2024} questioned whether generative models solve problems valuable enough to earn back the stock. We extend both from flow to stock and from level to growth rate.
    \item \textbf{Task-based macro limits.} Acemoglu~\cite{ref:acemoglu2024} bounds AI's contribution to total factor productivity --- output growth not explained by more labor or capital --- at roughly 0.66\% over ten years, by counting the share of tasks AI can do and the cost saving on each. Section~\ref{sec:slack} does the same accounting at the level of a firm's payroll: how much of the nominal saving on tasks becomes cash the firm can actually book.
    \item \textbf{The productivity J-curve.} Brynjolfsson, Rock and Syverson~\cite{ref:brynjolfsson2021} show that general-purpose technologies require complementary intangible investment before measured productivity rises. That lag is why enterprise adoption cycles (8--12 years) are longer than silicon lives (3--5), and it is also the strongest argument for building early: if the productivity is coming, capacity must be there first. Section~\ref{sec:demand} engages it directly. Solow's 1987 remark that computers were visible everywhere but in the productivity statistics~\cite{ref:solow1987} is the same phenomenon a generation earlier. Practitioners at the last mile report the same thing from experience: putting AI onto a poorly run process automates the process's defects, and the time goes into fixing the process, not the model.
    \item \textbf{Surplus capture.} Nordhaus~\cite{ref:nordhaus2004} estimates that innovators captured on the order of 2\% of the social surplus from their innovations over 1948--2001, the rest going to customers as lower prices and better products. A 2\% share is a 50:1 ratio of customer benefit to innovator revenue. Section~\ref{sec:lastmile} finds 10--35:1 for AI contact-center automation --- the same phenomenon, measured on one sector.
    \item \textbf{Scaling and spending data.} Epoch AI's capex~\cite{ref:epoch_capex} and revenue~\cite{ref:epoch_revenue} series, and Villalobos et al.\ on the human-text data wall~\cite{ref:epoch_data2024}, supply the empirical spine.
\end{enumerate}
The evidence for the build-out is strong and the paper uses it as input: Google Cloud's contracted backlog reached \$514B in Q2 2026, up from \$462B a quarter earlier, on 82\% revenue growth~\cite{ref:alphabet_q2_2026} (how much of it is AI-specific is not disclosed); the 2026 capex guidance has been revised upward all year~\cite{ref:mlq2026,ref:statista2026}; Goldman models \$5.3T of hyperscaler capex over 2025--2030~\cite{ref:yahoo_goldman2026}; and frontier-lab revenue has grown at multiples per year~\cite{ref:epoch_revenue}. Those are the observables the model must reconcile with, and the next section lists them with their provenance --- audited filing, company statement, or press report --- because they are not all of the same quality.

\section{What We Observe (2025--2026)}\label{sec:observe}
\begin{table}[htbp]
\centering
\small
\caption{Observables the model must reconcile with (\$B unless noted)}
\label{tab:observe}
\begin{tabularx}{\textwidth}{L{5.0cm}L{2.4cm}X}
\toprule
\textbf{Quantity} & \textbf{Value} & \textbf{Source / note} \\
\midrule
Big-5 capex 2022 / 2023 / 2024 / 2025 & 162 / 161 / 260 / 448 & Epoch AI from SEC XBRL, cash PP\&E plus new finance leases~\cite{ref:epoch_capex} \\
Big-4 capex guidance, 2026 (Q2 calls, Jul.~22--30) & 725--745 & Alphabet \$195--205B (raised from 180--190); Microsoft $\sim$\$175B calendar-2026 (FY26 cash capex \$115.9B, Q4 \$41B); Amazon $\sim$\$220B (raised from $\sim$200 on memory prices); Meta \$135--145B (floor raised twice)~\cite{ref:q2_scorecard,ref:alphabet_q2_2026,ref:futurex_q2}; Oracle adds $\sim$50 to make the Big-5 $\sim$780 \\
Q2 2026 quarterly capex & GOOG 44.9; MSFT 41; AMZN 54.2; META 31.1 & Company releases~\cite{ref:futurex_q2}; Alphabet free cash flow $-$\$5.9B, its first negative quarter \\
Google Cloud, Q2 2026 & \$24.8B (+82\%); backlog \$514B; op.\ income \$8.8B & Alphabet Q2 release, Jul.~22~\cite{ref:alphabet_q2_2026}; backlog up from \$462B in Q1 \\
Azure, FY26 Q4; AWS, Q2 2026 & Azure +43\%, TTM $>$\$100B; AWS +37\% & Microsoft Jul.~29; Amazon Jul.~30~\cite{ref:q2_scorecard,ref:directions_msft} \\
Microsoft 365 Copilot paid seats; M365 commercial seat growth & $>$30M (Jul.~2026); +6--7\% y/y through FY26 & Microsoft FY26 Q4 call and Form 8-K investor metrics~\cite{ref:q2_scorecard,ref:msft_8k_fy27} \\
Nvidia data-center revenue, FY2026 & 193.7 & Form 8-K, Feb.~25, 2026~\cite{ref:nvda_fy26}; Q1~FY27 \$75.2B, i.e.\ $\sim$\$300B run-rate~\cite{ref:nvda_q1fy27} \\
Anthropic run-rate: Dec.~2025 / Feb. / Apr. / May / Jul.~2026 & 9 / 14 / 30 / 47 / 65 & Company disclosure (Feb.)~\cite{ref:anthropic_seriesg}; Bloomberg via TechCrunch and Axios (Jul.)~\cite{ref:techcrunch_anthropic,ref:axios_anthropic} \\
OpenAI run-rate: Dec.~2025 / Aug.~2026 & 21 / 40 & Epoch revenue dataset, media reports~\cite{ref:epoch_revenue}; 2025 booked revenue \$13.1B against a \$20.9B operating loss (reported, unaudited)~\cite{ref:valueadd_openai}\verify \\
Microsoft AI revenue run-rate, Q3 FY26 & $>$37 & Earnings coverage~\cite{ref:yahoo2026}\verify \\
AWS AI revenue run-rate, Q1 2026 & $>$15 & Earnings coverage~\cite{ref:traderinsights2026}\verify \\
Meta Advantage+ run-rate, Q2 2026; disclosed AI relevance lifts & $>$75; +8\% clicks, +15.7\% conversions & Q2 2026 earnings call~\cite{ref:meta_q2_2026}; Q2 revenue \$60.8B, price per ad +12\%, impressions +14\%~\cite{ref:meta_q2_2026b} \\
OpenAI gross margin, Q1 2026; 2026 Azure spend & $\sim$39\%; $\sim$13 & Reported financials~\cite{ref:valueadd_openai}\verify \\
Anthropic reseller revenue recognition & Gross basis & Bedrock/Vertex end-customer spend booked as Anthropic revenue, partner payouts as expense~\cite{ref:sacra_anthropic} \\
Microsoft 365 commercial price increase, Jul.~1, 2026 & up to 43\% by plan; E3 $\to\sim$\$39/user/month & Copilot Chat and Security Copilot bundled into suites; Windows Enterprise also up~\cite{ref:m365_price_wl,ref:m365_price_uscloud} \\
Enterprise software market value lost, Jan.--Apr.~2026 & $\sim$\$2T; \$285B in 48 h & The ``SaaSpocalypse'': largest non-recessionary software drawdown in 30 years (JPMorgan); Atlassian's first seat-count decline~\cite{ref:saaspocalypse_ti,ref:saaspocalypse_fc} \\
European public-sector exits from Microsoft & France (national), Schleswig-Holstein, Denmark & DINUM Windows-to-Linux mandate, Apr.~2026; 20\%$\to$40\% of German IT leaders implementing European software in a year~\cite{ref:eu_sovereignty,ref:tuta_sovereignty} \\
US data-center projects blocked or delayed by local opposition & \$156B (2025); \$130B (Q1 2026 alone) & Data Center Watch~\cite{ref:dcw_q1_2026}; 75 projects in Q1 2026; 833 active opposition groups in 49 states; New York moratorium Jul.~2026~\cite{ref:tomshw_dcw} \\
Oracle Q1 FY2027 (Sept.~10, 2026) & RPO \$664B; OCI \$7.4B (+121\%); GPU utilization 97.9\% & RPO up \$26B q/q, over half from OpenAI; about half converting within 36 months; new contracts ``via prepay or bring-your-own-hardware''; FY26 capex \$55.7B, free cash flow $-$\$23.7B; $\sim$\$40B more financing planned FY27~\cite{ref:orcl_q1fy27,ref:orcl_q1fy27_mw} \\
OpenAI booked revenue, Q1 / Q2 2026; run-rate Aug.; 2025 loss & 5.7 / 6.7; 40; loss \$20.9B on \$13.1B & Sacra and press~\cite{ref:sacra_openai,ref:valueadd_openai}; IPO deferred to 2027; FT-reported \$278B cash-burn plan\verify \\
Anthropic expected 2026 exit run-rate & $>$100 & Bloomberg via NYT, Sept.~15, 2026; IPO shares expected to trade by November~\cite{ref:valueadd_openai}\verify \\
ChatGPT users; Claude app users & 1B MAU (Jun.); 900M WAU (Feb.); 56M MAU (Q2) & OpenAI/Sensor Tower via Reuters; Sacra~\cite{ref:sacra_openai,ref:aibw_openai} \\
Street consensus Big-5 capex, 2027 / 2028 & $\sim$1,050 / $\sim$1,150 & S\&P Capital IQ consensus as of Sept.~18, 2026, via a16z~\cite{ref:a16z_som}; J.P.~Morgan: 791 / 1,066 / 1,179 for 2026--28~\cite{ref:jpmam_gtm} \\
Cloud backlogs (MSFT RPO + Google Cloud + AMZN RPO), Q2 2026 & $\sim$1,700; +$\sim$150\% y/y & Company filings compiled by a16z~\cite{ref:a16z_som} \\
GPU rental prices, Sept.~2026 (\$/GPU-hour) & H100 2.63; A100 1.59; H200 3.29; B200 5.69 & Silicon Data index via a16z~\cite{ref:silicon_data,ref:a16z_som}; residual values rising through 2026 \\
Consensus hyperscaler free cash flow & $\sim$0 in 2026--27; $>$\$500B by 2030 & FactSet / Goldman Sachs via a16z~\cite{ref:a16z_som} \\
Token share on OpenRouter, mid-2026 & Frontier labs $\sim$69\% of spend on $\sim$35\% of tokens; tokens +600$\times$ in 2 yr & OpenRouter via a16z~\cite{ref:openrouter,ref:a16z_som} \\
Lab equity raised, 2026 YTD & $\sim$190 & OpenAI \$122B at \$852B; Anthropic \$65B at \$965B~\cite{ref:valueadd_openai,ref:techcrunch_anthropic} \\
Goldman Sachs Phase-2 ROIC framework (Sept.~24, 2026) & \$1.42T revenue 2028--30; \$11.6B/GW-yr & Six US hyperscalers' 2026--27 AI capex \$1.726T ($\sim$41 GW at \$42.4B/GW; 70/30 compute/shell at 5/15 yr; opex \$836M/GW-yr; 15\% ROIC on average annual capex); sensitivity \$0.91--1.89T over 0--30\% ROIC~\cite{ref:gs_roic2026}; Section~\ref{sec:gs} \\
Backlog conversion the three CSPs would need (GS) & 59\% of \$1.69T over 2028--30 & Alphabet 78\% of \$496B, Amazon 68\% of \$514B, Microsoft 40\% of \$678B; assumes no backlog growth~\cite{ref:gs_roic2026} \\
AI infrastructure capex from the semiconductor build (GS), 2026 / 27 / 28 & 861 / 1,338 / 2,008; 20 / 35 / 57 GW & All buyers (hyperscalers, labs, sovereigns, enterprises) from Nvidia, Broadcom, AMD, Marvell, MediaTek revenue per GW; implied \$42 / 38 / 35B per GW~\cite{ref:gs_roic2026}; 57 GW exceeds the paper's grid cap (22--40 GW/yr) \\
\midrule
\textbf{Calendar-2026 AI revenue, all channels} & \textbf{$\sim$178 (165--190)} & Seven-component decomposition, Section~\ref{sec:decomp}, Table~\ref{tab:basedecomp} \\
\bottomrule
\end{tabularx}
\end{table}
Provenance matters here. The capex series and Nvidia's revenue are audited filings; the lab run-rates are company statements or press reports citing them; the Microsoft, AWS and OpenAI figures are earnings-call remarks and leaked financials, not audited disclosures, and are marked accordingly. Meta's AI spending is monetized through advertising rather than sold as a product, so it has no ``AI revenue'' line at all; the paper treats it as an advertising-uplift channel (Section~\ref{sec:decomp}). Three things follow. Nvidia's \$300B data-center run-rate is consistent with the capex series and with a silicon share near half. Frontier-lab revenue is growing at multiples per year, which is the bull case realized so far. And roughly \$190B of lab equity was raised in 2026 alone, a sum comparable to the labs' combined revenue --- so a material share of the compute they buy from hyperscalers this year is paid with investor capital, not customer cash.

\subsection{Do the five, and the Chinese four, tell the same story?}
They do, and the consistency is itself evidence. Across Alphabet, Microsoft, Amazon and Meta the Q2 2026 pattern is identical: capital expenditure roughly doubling year on year, guidance raised during the year, cloud revenue growth accelerating (Google Cloud +82\%, Azure +43\%, AWS +37\%), and free cash flow turning negative at the two heaviest spenders~\cite{ref:alphabet_q2_2026,ref:futurex_q2}. The Chinese hyperscalers are on the same curve a year behind and at smaller scale: in Q2 2026 Alibaba's capex was RMB~67.7B (+75\%) and Tencent's RMB~52.8B (+176\%), each about a quarter of quarterly revenue; Alibaba's AI cloud and compute revenue grew 45\% while adjusted profit fell 38\%; Tencent's free cash flow went negative for the first time in its history on RMB~51B of prepayments for computing power; Baidu's capex reached 36\% of revenue with negative free cash flow~\cite{ref:china_q2_capex,ref:36kr_china}. TrendForce puts 2026 capex for the nine largest cloud operators --- the Big-5 plus ByteDance, Tencent, Alibaba and Baidu --- at \$830--890B, up 79--90\%, with the Chinese four raising AI capex $\sim$80\%~\cite{ref:trendforce_csp9,ref:taiwannews_csp9}. The same source notes Chinese cloud returns on invested capital of 13--20\% against 25--50\% for US peers~\cite{ref:china_q2_capex}. Two implications. The growth rates the paper's near-term chains assume (40--80\% into 2027) are what every large operator on both sides of the Pacific is reporting, so the base case is not built on a US anomaly. And the Chinese four add \$50--110B/yr of capex that is outside the Big-5 register and outside its demand pools; the model is a Big-5 model, and the Chinese stock competes for the same tokens with cheaper power and less efficient silicon (Section~\ref{sec:commodity}).

\section{The Supply-Side Carry: A Vintage Model}\label{sec:supply}
\subsection{Cohorts, not a snapshot}
Let $I_c$ be Big-5 capex in cohort year $c$ (Table~\ref{tab:observe}; Figure~\ref{fig:capex}) and $A_c=\max(I_c-N_c,0)$ the AI-attributable part after subtracting the non-AI baseline $N_c=N_{2023}(1+g_N)^{c-2023}$. The baseline is constrained by data: total capex was flat at \$161B in 2023 while AI purchases began~\cite{ref:epoch_capex}, so $N_{2023}=161-A_{2023}$ with $A_{2023}\sim$Triang(20, 30, 40), and it grows at $g_N\sim$Triang(2, 6, 10)\%/yr for ordinary cloud refresh and logistics (Section~\ref{sec:decomp}). AI-attributable capex above a grid energization ceiling of Triang(1{,}150, 1{,}400, 1{,}680) \$B/yr --- 22--40 GW of new IT load at \$38--55M per MW --- is deferred to the following year and carries return but not depreciation (Section~\ref{sec:decomp}); the ceiling does not bind on a plateau and binds in 5\% of arms-race draws. Separately, a fraction $f_{\text{delay}}\sim$Triang(0.10, 0.20, 0.35) of each cohort from 2026 is deferred a year by siting, consent and water constraints; shell cost carries an uplift $\kappa_{\text{site}}\sim$Triang(1.02, 1.05, 1.09); and facility opex carries $(1+u)$ with a cooling-uplift $u\sim$Triang(0, 0.03, 0.06) for water-constrained cooling (Section~\ref{sec:decomp}). Each cohort splits into a silicon share $\phi\sim$Triang(0.45, 0.52, 0.60) --- anchored on Nvidia's data-center revenue (Section~\ref{sec:decomp}) --- with an economic life $L_{\text{si}}$ that is the fleet-weighted harmonic mean of the training-cluster life $L_{\text{tr}}\sim$Triang(2.5, 3.5, 4.5) yr and the inference-fleet life $L_{\text{inf}}\sim$Triang(4, 5, 6) yr at a training share $f_{\text{tr}}\sim$Triang(0.35, 0.45, 0.55) --- effective $L_{\text{si}}\approx4.2$ yr (P10--P90 3.75--4.6) --- and a shell/power/cooling share with life $L_{\text{sh}}\sim$Triang(15, 20, 25) yr. Silicon is energized in its spend year at an energization fraction $\ell_{\text{si}}\sim$Triang(0.5, 0.7, 0.9) and shell at $\ell_{\text{sh}}\sim$Triang(0.2, 0.4, 0.6), the rest the year after (lower-case $\ell$ is a fraction; upper-case $L$ is a life in years). With straight-line depreciation and a mid-year convention, the carry $H_t$ in year $t$ --- the annual cost of owning the installed base, summed over cohorts $c$ and over the two components $j$ (silicon, shell) --- is
\begin{equation}
H_t=\sum_{c\le t}\sum_{j\in\{\text{si},\text{sh}\}}\left[\frac{G_{c,j}}{L_j}+c_{\text{opex}}\,G_{c,j}+r\,G_{c,j}\!\left(1-\frac{t-c+\tfrac12}{L_j}\right)\right]\mathbf{1}\!\left[t-c+\tfrac12<L_j\right],
\label{eq:Ht}
\end{equation}
where $G_{c,j}$ is the in-service gross value of component $j$ of cohort $c$, $L_j$ its life, the facility-opex share $c_{\text{opex}}\sim$Triang(0.059, 0.0665, 0.074) of in-service gross (decomposed in Section~\ref{sec:decomp}), the cost of capital $r\sim$Triang(0.090, 0.0945, 0.099) from CAPM components (Section~\ref{sec:decomp})~\cite{ref:damodaran2026}, and $\mathbf{1}[\cdot]$ an indicator that is 1 while the cohort is younger than its life. Forward paths: \textbf{Plateau} (780 flat), \textbf{Arms race} (1{,}000, then 1{,}100; consistent with the \$1T+ 2027 projections in~\cite{ref:mlq2026}), \textbf{Rationalize} (550, then 500).

\begin{figure}[htbp]\centering\includegraphics[width=0.85\textwidth]{f1_capex.pdf}
\caption{Big-5 capex by company, 2022--2025 (Epoch AI, SEC XBRL~\cite{ref:epoch_capex}) and 2026 guidance~\cite{ref:mlq2026,ref:statista2026}. The dashed line is the 2022 level used as the non-AI baseline.}\label{fig:capex}\end{figure}

\begin{figure}[htbp]\centering\includegraphics[width=0.85\textwidth]{f2_carry.pdf}
\caption{Annual carry $H_t$ of the AI-attributable stock under three capex paths (Monte Carlo medians; band is P10--P90 for Plateau).}\label{fig:carry}\end{figure}

\begin{table}[htbp]
\centering\small
\caption{Annual carry $H_t$ (\$B, median; Plateau P10--P90 in parentheses)}
\label{tab:carry}
\begin{tabular}{lcccc}
\toprule
\textbf{Path} & \textbf{2027} & \textbf{2028} & \textbf{2029} & \textbf{2030} \\
\midrule
Plateau (780) & 390 (365--415) & 535 (505--570) & 650 (615--685) & 725 (685--770) \\
Arms race (1000, 1100) & 425 & 650 & 850 & 1{,}015 \\
Rationalize (550, 500) & 350 & 430 & 465 & 460 \\
\bottomrule
\end{tabular}
\end{table}
In 2028 under Plateau the median carry decomposes as depreciation $\sim$\$280B, facility opex $\sim$\$130B, return on net book and on deferred capital $\sim$\$125B. The carry ramps because 2027 is the cheapest year --- the 2026 cohort is half energized and the 2027 cohort barely --- and rationalization flattens rather than cuts because the 2025--26 cohorts are on the books whatever happens next.

\subsection{Silicon life: accounting versus economics}
The operators' own accounting lives for servers are 5 years at Amazon (cut from 6 in January 2025, explicitly citing the pace of AI hardware change~\cite{ref:amzn_10k_life}), 5.5 at Meta~\cite{ref:hudson_meta}, and 6 at Microsoft and Alphabet~\cite{ref:hudson_msft,ref:hudson_policy}. We use an economic life of 3--5.5 years because power, not silicon health, is the binding constraint: where interconnection is queued~\cite{ref:pjm_queue,ref:eirgrid_cru,ref:amsterdam}, a 700\,W prior-generation part cannot hold a substation slot against a successor with several times the throughput per watt; secondary-market volume cannibalizes new cluster pricing; and 24/7 thermal duty is an accelerated-aging regime for stacked memory. Against this, 2020-vintage A100s remained rented through 2026, and Alphabet and Microsoft have not shortened their lives. So we show the result both ways: with accounting lives (4.5--6.5 yr) the plateau break-even for 2030 moves from 49\% to 51\%, and with an 8-year life it is 49\% --- longer lives lower annual depreciation but keep every cohort on the books longer, and the two effects nearly cancel by 2030 (Table~\ref{tab:breakeven}). The dual-use defense --- repurposing 40--100\,kW/rack liquid-cooled halls for databases and web hosting --- does not recover the carry; the capital is committed to dense tensor math.

\subsection{Thermal acceleration: high failure rates are the intrinsic norm}
The accounting lives of 5--6 years were set on the experience of CPU servers at ASHRAE-class inlet temperatures. AI accelerators are a different thermal regime: 700--1{,}000\,W parts at junction temperatures of 80--95$^{\circ}$C under sustained duty, with stacked HBM and through-silicon vias in the same envelope. Three standard reliability relations bear on their economic life:
\begin{itemize}
    \item \textbf{Arrhenius acceleration.} For temperature-driven wear-out mechanisms the acceleration factor is $AF=\exp\!\left[\frac{E_a}{k}\left(\frac{1}{T_1}-\frac{1}{T_2}\right)\right]$; with the JEDEC default activation energy $E_a\approx0.7$\,eV~\cite{ref:jedec_jep122} this is close to the industry rule of thumb that failure rates double for every 10$^{\circ}$C~\cite{ref:arrhenius_rule}. A part held 20$^{\circ}$C hotter than the population its accounting life was calibrated on has roughly a quarter of that life.
    \item \textbf{Electromigration.} Black's equation, MTTF $\propto J^{-n}e^{E_a/kT}$~\cite{ref:black1969}, couples the same exponential to current density; accelerators run both high.
    \item \textbf{Thermal cycling.} HBM stacks and TSV interconnects fail by solder and micro-bump fatigue under thermal cycling (Coffin--Manson), which is driven by load swings rather than steady temperature; inference fleets with bursty duty see both mechanisms at once.
\end{itemize}
Field data agree. Meta's Llama~3 training run on 16{,}384 H100s recorded 419 unexpected interruptions in 54 days, 58\% attributed to GPU faults including HBM3 errors --- about one accelerator-related interruption every three hours at a scale that is now routine~\cite{ref:llama3}. High failure is not an anomaly of early hardware; it is the expected behavior of silicon run at the top of its thermal envelope.

Three consequences for the model. First, the economic-life distribution Triang(3, 4, 5.5) is, if anything, generous; a thermally motivated Triang(2.5, 3.5, 5) raises the 2027--28 carry by \$10--15B/yr and lowers 2029--30 slightly as cohorts roll off faster (2030 break-even $g^{*}$ moves from 49\% to 48\%). Second, and more important, shorter life converts a growing share of a flat \$780B/yr into \emph{replacement} rather than \emph{growth}: at a 3.5-year life the plateau path is buying back roughly a third of its own stock every year, so capacity --- and therefore revenue potential --- rises much more slowly than cumulative spend. Third, the salvage and second-life defenses that follow are weaker than their proponents assume, because the parts being down-cycled are the ones that have already consumed most of their Arrhenius budget.

\subsection{The salvage fallacy in detail}
Corporate disclosures have extended server useful lives to six years~\cite{ref:hudson_policy,ref:amzn_10k_life}, and a common defense holds that frontier accelerators can be down-cycled into batch inference for 6--8 years at meaningful residual value. Three constraints bound this:
\begin{itemize}
    \item \textbf{Power is the binding input.} Where interconnection is queued (PJM, ERCOT, Dublin, Amsterdam)~\cite{ref:pjm_queue,ref:eirgrid_cru,ref:amsterdam}, a 700W prior-generation part cannot hold a substation allocation against a successor delivering several times the throughput per watt. Older parts are de-racked for the grid slot, not for the silicon.
    \item \textbf{Competing with one's own ghost.} Secondary-market volume of prior-generation parts cannibalizes new cluster pricing, as it did for mainframes in IBM's own history.
    \item \textbf{Thermal duty.} 24/7 operation at peak thermal load is a known accelerated-aging regime for stacked memory and TSV interconnect; a multi-year second life at full utilization is not the default assumption an engineer would make. We have not found a public, peer-reviewed failure-rate study specific to HBM3e/HBM4 under hyperscale duty and do not cite one\verify.
\end{itemize}
Against this, the empirical record is mixed: 2020-vintage A100s remained fully utilized in rental markets through 2026. We therefore use $L_{\text{si}}\sim$Triang(3, 4, 5.5) rather than the 2.5--3.5 years of harsher analyses, and note that even $L_{\text{si}}=8$ does not by itself close the gap by 2030 (Table~\ref{tab:breakeven}).

\subsection{The dual-use fallacy in detail}
AI halls are engineered for 40--100 kW/rack with direct-to-chip liquid cooling and non-blocking fabrics, at several times the cost per megawatt of 8--12 kW enterprise halls. Repurposing them for relational databases and web hosting does not recover the carry; the capital is physically committed to dense tensor math.



\section{Demand: Base, Growth, Quality}\label{sec:demand}
\subsection{Anchor to what is observed}
Terms, first, because three flows are easy to confuse. \emph{End-customer revenue} $D_t$ is what buyers --- enterprises, consumers, advertisers --- pay for AI products and services in year $t$; $R_s(t)$ is the same thing for one sector $s$, and $D_t=\sum_s R_s(t)$. \emph{Infrastructure receipts} are the part of that money that reaches the owners of the data centers, after the frontier labs keep their gross margin and after the part that runs on devices, on customers' own premises or on integrators' small models --- collectively, \emph{leakage} --- never reaches a cloud meter at all. The \emph{cloud-capture share} $\theta_s$ is the fraction of a sector's end-customer revenue that becomes infrastructure receipts. Figure~\ref{fig:flow} draws the flows; the carry test compares infrastructure receipts, not end-customer revenue, with $H_t$.

\begin{figure}[htbp]\centering\includegraphics[width=\textwidth]{f18_flow.pdf}
\caption{Where an AI customer's dollar goes. End-customer revenue enters through three channels; the frontier labs keep a margin and spend the rest on compute; part of every channel leaks to devices, premises and integrators; what remains is the infrastructure receipts that must cover the carry.}\label{fig:flow}\end{figure}

Calendar-2026 end-customer AI revenue across all channels is $D_{2026}\sim$Triang(163, 179, 194) \$B, the sum of seven sourced components, the hyperscaler component itself decomposed from Microsoft's disclosures (Section~\ref{sec:decomp}). Not all of it reaches the infrastructure layer: frontier labs keep a margin and spend the majority of their costs on compute~\cite{ref:epoch_spending}, so we apply an infrastructure share $\theta\sim$Triang(0.77, 0.815, 0.86), a channel-weighted decomposition with the lab channel itself decomposed by lab (Section~\ref{sec:decomp}). Demand then grows at a sustained rate $g$:
\begin{equation}
D_t=\theta\,D_{t_0}(1+g)^{t-t_0},\qquad t_0=2026,
\end{equation}
where $\theta$ is the revenue-weighted cloud-capture share, $D_{t_0}$ the 2026 base, and $g$ the sustained growth rate in percent per year (a $g$ of 0.49 means 49\%/yr). Does $g$ itself follow a curve? In the bottom-up estimate it does: the sectors' implied growth falls from $\sim$58\% in 2027 to $\sim$14\% in 2030 (Table~\ref{tab:bottomup}), the shape of a logistic. The top-down sweep deliberately holds $g$ constant so that a single number answers ``how fast must it grow'' --- a sustained rate is easier to test against the labs' disclosures than a decaying one. This makes $g$ the single parameter the debate turns on, and it answers the J-curve objection directly: if productivity arrives late but arrives~\cite{ref:brynjolfsson2021}, then building early is right --- and the bull must name the $g$ at which it pays.

\subsection{Revenue quality}
Two caveats on the base. First, run-rates overstate booked revenue: OpenAI's reported 2025 booked revenue was \$13.1B against an end-2025 run-rate of \$21B~\cite{ref:valueadd_openai,ref:epoch_revenue}; we use calendar-year figures. Second, a material share of 2026 lab compute purchases is funded by the $\sim$\$190B of equity raised this year rather than by customer revenue~\cite{ref:valueadd_openai,ref:techcrunch_anthropic}. OpenAI's 2025 operating loss is reported to have exceeded its revenue~\cite{ref:valueadd_openai}; the figure is leaked and unaudited, and the paper treats it with corresponding skepticism, using only its direction. That skepticism is itself a cost of the disclosure regime: a company at this scale, with these commitments, that publishes run-rates in the months of its choosing and lets its losses reach the market by leak, generates churn --- in its counterparties' recoverability forecasts, in its lenders' covenants, and in analyses like this one --- that a single audited statement would remove. The same applies to every lab and to the hyperscalers' ``AI revenue'' lines: partial disclosure does not protect the discloser, it adds noise to everyone's estimate of it, including the register this paper maintains. Those receipts are real cash to hyperscalers and are counted in $\theta D$; but they are bridging, not durable, and the model's $g$ is a growth rate of customer revenue, not of investor capital. Related-party round trips --- Microsoft--OpenAI, Nvidia--CoreWeave, Amazon/Alphabet--Anthropic --- are of order \$10B/yr\verify and are small against this; the investor-funded channel is the circularity that matters in 2026.

\section{Bottom-Up Demand: Fifteen Sector Fermi Chains}\label{sec:bottomup}
The top-down path in Section~\ref{sec:demand} has one free parameter, $g$. A Fermi estimate is only as good as the number of independent factors in it, so we build a second estimate that does not use $g$ at all. Each sector $s$ is a chain of four estimates: a calendar-2026 base $R_s$ (calibrated so the sectors sum to the observed \$179B --- the base is not independent evidence; the ceilings and ramps are), a 2030 ceiling $C_s$ derived from a physical or budgetary constraint (Appendix~\ref{app:sectors}), a logistic ramp rate $k_s$, and a cloud-capture share $\theta_s$ (the fraction of end-customer revenue that reaches the infrastructure layer after lab margins, on-device execution, on-premise deployment and integrator capture). End-customer revenue follows
\begin{equation}
R_s(t)=\frac{C_s}{1+\left(C_s/R_{s,0}-1\right)e^{-k_s(t-t_0)}},\qquad t_0=2026,\qquad D^{\text{BU}}_t=\sum_s R_s(t),\qquad \text{Receipts}_t=\sum_s\theta_s R_s(t),
\end{equation}
where $R_{s,0}$ is the sector's 2026 revenue (\$B), $C_s$ its ceiling (\$B), and $k_s$ its ramp rate in units of per year: at $k=0.8$ the sector covers half the distance to its ceiling in about $\ln 2/0.8\approx0.9$ years. The curve starts at $R_{s,0}$, rises fastest early, and flattens as it approaches $C_s$; the sector's growth rate is not an input but falls out of the curve.
Two common factors --- $M_C$ on ceilings ($\sim$Triang(0.7, 1.0, 1.4)) and $M_k$ on ramp rates ($\sim$Triang(0.7, 1.0, 1.3)) --- are multipliers drawn once per trial and applied to every sector. They represent the macroeconomic background the sectors share --- a recession, an interest-rate shock, a regulatory turn, or a capability step that makes agents trustworthy everywhere at once --- and they keep fifteen independent draws from averaging to an implausibly narrow band. Sectors are listed in order of their 2030 ceiling. Sectors exposed to sovereign procurement carry the $\gamma_{\text{sov}}$ carve-out in the sovereignty scenario, and in every scenario a European displacement ramp $g_{\text{EU}}\sim$Triang(3, 5, 8)\% by 2030 for the movement of European buyers away from US vendors documented in Section~\ref{sec:decomp}.

\begin{table}[htbp]
\centering\scriptsize
\caption{Sector parameters (\$B; triangular min/mode/max) and the constraint that sets each ceiling}
\label{tab:sectors}
\begin{tabularx}{\textwidth}{L{2.9cm}L{1.6cm}L{1.9cm}L{1.5cm}L{1.6cm}X}
\toprule
\textbf{Sector} & \textbf{2026 base $R_{s,0}$ (\$B)} & \textbf{2030 ceiling $C_s$ (\$B)} & \textbf{Ramp $k_s$ (1/yr)} & \textbf{Cloud share $\theta_s$ (0--1)} & \textbf{Ceiling constraint (Appendix~\ref{app:sectors})} \\
\midrule
Office assistants / agents & 20/30/45 & 100/150/230 & 0.5/0.8/1.1 & 0.45/0.6/0.75 & Disclosed paid seats (20M M365 Copilot, Apr.~2026) grown with decaying rate, $\times$ price, $\times$ (1 + agentic multiplier), hard-capped at 3--5.5\% of 2030 IT spend; Section~\ref{sec:decomp} \\
Advertising uplift & 25/40/55 & 105/136/170 & 0.4/0.7/1.0 & 0.9/1.0/1.0 & 8--18\% relevance/pricing lift on 2030 Meta+Google ads (\$700--950B) plus AI-native formats; decomposed in Section~\ref{sec:decomp} \\
Consumer subscriptions & 18/25/35 & 50/64/80 & 0.5/0.8/1.2 & 0.75/0.85/0.95 & 1.3--2.0B unique users $\times$ 4--8\% paid $\times$ \$20--26/month, plus free-tier ads and premium tiers; Section~\ref{sec:decomp} \\
Coding tools & 15/20/28 & 40/60/90 & 0.6/0.9/1.3 & 0.7/0.8/0.9 & 29M developers $\times$ paid share $\times$ seat+token spend; non-recurring refactor impulse \\
Search \& internal use & 12/20/30 & 30/50/80 & 0.4/0.7/1.0 & 0.9/1.0/1.0 & Defends rather than adds to \$250B search; internal productivity \\
Financial services & 6/10/15 & 25/40/60 & 0.4/0.7/1.0 & 0.6/0.75/0.9 & AI share of $\sim$\$1.5T banking/insurance IT and ops spend; model-risk governance \\
Customer service / CX & 5/8/12 & 25/40/60 & 0.5/0.8/1.2 & 0.35/0.5/0.65 & 10--15\% of \$350B contact-center labor displaced; vendors capture a share \\
Public sector / defense & 5/8/12 & 25/40/60 & 0.3/0.5/0.8 & 0.7/0.85/1.0 & $\sim$5\% of $\sim$\$700B government IT; sovereign carve-out applies \\
Healthcare / clinical & 4/6/9 & 20/35/50 & 0.3/0.5/0.8 & 0.5/0.65/0.8 & Ambient scribes, prior-auth, coding; liability and procurement pace \\
Autonomous fleets & 5/8/12 & 16/27/43 & 0.3/0.5/0.8 & 0.3/0.45/0.6 & Fleet software market forecasts~\cite{ref:av_fbi}; NHTSA and chassis cycles; on-vehicle compute \\
Generative video / gaming & 5/8/12 & 17/26/37 & 0.4/0.7/1.0 & 0.8/0.9/1.0 & Cloud GPU COGS vs consumer budgets~\cite{ref:newzoo} \\
Legal \& compliance & 3/5/8 & 12/20/30 & 0.3/0.5/0.8 & 0.5/0.65/0.8 & $\sim$2--3\% software share of \$950B legal services; verification rules \\
Science / R\&D & 2/4/7 & 10/25/45 & 0.4/0.7/1.0 & 0.7/0.85/1.0 & $\sim$5\% of \$300B pharma R\&D plus materials, climate, engineering \\
Factory / robotics software & 2/4/6 & 8/15/25 & 0.3/0.5/0.8 & 0.3/0.45/0.6 & Robot stock $\times$ software fee~\cite{ref:ifr}; Asian wage floor; edge inference \\
Eldercare / assistive & 0.5/1/2 & 3/6/12 & 0.3/0.5/0.8 & 0.1/0.15/0.25 & 160M OECD seniors $\times$ penetration $\times$ fee; 85\% on-device \\
\midrule
\textbf{Sum (medians)} & \textbf{$\sim$205} & \textbf{$\sim$740} & & \textbf{$\sim$0.76} & \\
\bottomrule
\end{tabularx}
\end{table}

\begin{figure}[htbp]\centering\includegraphics[width=0.9\textwidth]{f9_sectors.pdf}
\caption{Bottom-up end-customer AI revenue by sector, 2026--2030 (medians). Four sectors carry three-quarters of the total.}\label{fig:sectors}\end{figure}

\begin{table}[htbp]
\centering\small
\caption{Bottom-up demand (\$B; median with P10--P90) and implied year-on-year growth}
\label{tab:bottomup}
\begin{tabular}{lccccc}
\toprule
 & \textbf{2026} & \textbf{2027} & \textbf{2028} & \textbf{2029} & \textbf{2030} \\
\midrule
End-customer revenue & 204 (191--218) & 321 (291--355) & 445 (388--513) & 551 (469--648) & 626 (526--744) \\
Infrastructure receipts ($\sum\theta_s R_s$) & 160 (148--172) & 250 (225--277) & 344 (299--398) & 423 (360--499) & 479 (402--570) \\
Implied growth, end-customer & & 57\% & 39\% & 24\% & 14\% \\
\bottomrule
\end{tabular}
\end{table}

Growth here is a consequence, not an input: it is high while the large sectors are far from their ceilings and falls as they approach them. Three readings. First, the sectors are consistent with the observed 2026--27 pace: 61\% growth into 2027 is in line with the labs' guided deceleration~\cite{ref:epoch_revenue_x}, so the bottom-up model is not a bear case in the near term. Second, growth decays fast because the four large sectors approach ceilings set by seat counts, IT budgets and advertising yields, not by model capability. Third, the cloud-capture share matters as much as the ceilings: a quarter of end-customer revenue never reaches the infrastructure layer --- it stays with labs as margin, runs on-device, or accrues to integrators.

\subsection{Structural limits on the ceilings}

\textbf{Startup mortality.} Of the 2024--25 funded cohort, roughly 40\% had failed within 24 months and the trajectory points to $\sim$80\% by end-2026\verify~\cite{ref:cbinsights}. When labs ship native features, single-feature wrappers collapse; their API spend is the low tail of the enterprise pool.

\textbf{The local-silicon siphon.} About 1.2B smartphones and 250M PCs ship annually with 40--80 TOPS NPUs. If 350M knowledge workers issue $\sim$15 routine text tasks per day at $\sim$500 tokens each, routine volume is
\begin{equation*}
350\text{M}\times 15 \times 250 \times 500 \approx 6.6\times 10^{14}\ \text{tokens/yr},
\end{equation*}
worth of order \$70--330B/yr at routine-model prices of \$0.10--0.50 per million tokens. If 60--85\% of that runs on-device or on private small models, the cloud loses \$40--280B/yr of potential base-load. On-device inference is not free (device cost, power, governance), but its marginal price to the cloud is zero.

\textbf{Frontier versus routine.} Multi-step reasoning and complex synthesis are under $\sim$5\% of enterprise token volume; the 90\%+ that is summarization, drafting and formatting is exactly the share that distilled open-weight models and NPUs absorb.


\section{Decomposing the Dominant Uncertainties}\label{sec:decomp}
Left as single judgment ranges, three inputs would carry most of the model's uncertainty: the 2026 revenue base (64\% of top-down variance at a Triang(160, 210, 300) range), the cloud-capture share $\theta$ (18\%), and the advertising ceiling (51\% of bottom-up sector variance). An estimate improves in exactly one way --- by decomposing its loosest factor into tighter, separately sourced ones --- so this section does that for those three and for a further sixteen, and then draws what-if curves over what is left. Each decomposition is reported with the undecomposed range it replaces, and with the direction it moved the central estimate, because a decomposition that only ever tightens in one direction would be a sign of motivated choices. Appendix~\ref{app:inputs} appraises every input in turn --- what is known, the range chosen, the likely direction of bias in its source, and what would improve it --- so that any reader can substitute a better-informed range.

\subsection{The 2026 revenue base}
\begin{table}[htbp]
\centering\small
\caption{Calendar-2026 end-customer AI revenue, decomposed (\$B; triangular min/mode/max)}
\label{tab:basedecomp}
\begin{tabularx}{\textwidth}{L{5.2cm}L{2.2cm}X}
\toprule
\textbf{Component} & \textbf{Range} & \textbf{Basis} \\
\midrule
Anthropic, booked 2026 (gross basis) & 50 / 57 / 65 & Q1 \$4.7B, Q2 $>$\$11.5B reported~\cite{ref:axios_anthropic}; end-July run-rate \$65B~\cite{ref:techcrunch_anthropic} implies Q3 $\sim$\$16--19B and Q4 $\sim$\$20--28B \\
OpenAI, booked 2026 & 26 / 29 / 33 & Run-rate \$21B (Dec.~2025), $\sim$\$33B (May), $>$\$40B (Aug.)~\cite{ref:epoch_revenue,ref:sacra_anthropic,ref:valueadd_openai} \\
Other labs (xAI, Mistral, Cohere, Perplexity, non-US) & 4 / 7 / 12 & Sparse disclosure\verify \\
Hyperscaler AI services excluding lab hosting and reseller pass-through & 25 / 35 / 45 & Microsoft AI run-rate $>$\$37B less Azure-OpenAI hosting; AWS AI $>$\$15B less Bedrock pass-through; Google Cloud AI less Vertex pass-through~\cite{ref:yahoo2026,ref:traderinsights2026}\verify \\
AI-attributable advertising uplift (Meta + Google) & 25 / 40 / 55 & Meta: 8--12\% relevance/pricing lift disclosed~\cite{ref:meta_q2_2026} on $\sim$\$230B ads $=$ \$18--28B; Google assumed comparable magnitude with less disclosure, \$10--25B \\
Consumer, search and internal use not already in lab revenue & 8 / 15 / 25 & Gemini, Microsoft Copilot consumer, AI search monetization, internal productivity\verify \\
Overlap correction & $-15$ / $-8$ / 0 & Anthropic books Bedrock and Vertex end-customer spend as gross revenue~\cite{ref:sacra_anthropic}, so a portion of ``AWS AI'' and ``Google Cloud AI'' is the same dollar counted twice \\
\midrule
\textbf{Sum} & \textbf{178 (P10 165, P90 190)} & $\sigma_{\ln}=0.053$, versus 0.13 for the undecomposed Triang(160, 210, 300) \\
\bottomrule
\end{tabularx}
\end{table}
Two things happen relative to a single judgment range of Triang(160, 210, 300). The uncertainty falls by a factor of 2.5, because the two largest components are public quarterly figures rather than a guess. And the central value falls from \$210B to \$178B, because run-rates overstate calendar bookings and because gross-basis reseller accounting double-counts. Decomposition does not only tighten an estimate; it can move it, and here it moves against the build-out.

\subsection{The cloud-capture share}
$\theta$ is a channel-weighted average of four shares that are each better known than the aggregate: the lab channel ($\sim$50\% of the base), where infrastructure receipts are the labs' compute cost as a fraction of revenue --- OpenAI at $\sim$39\% gross margin and $\sim$\$13B of 2026 Azure spend on $\sim$\$29B of bookings~\cite{ref:valueadd_openai}, Anthropic's margin undisclosed\verify --- giving $\theta_{\text{lab}}\sim$Triang(0.50, 0.65, 0.85); hyperscaler direct services ($\sim$20\%), $\theta\sim$Triang(0.9, 1.0, 1.0); advertising ($\sim$22\%), $\theta=1$; and consumer/other ($\sim$8\%), Triang(0.75, 0.85, 0.95). The weighted result is $\theta\sim$0.82 (P10--P90 0.77--0.87), $\sigma_{\ln}=0.045$ versus 0.072 before. What remains uncertain is one number --- the labs' compute-cost share --- and Figure~\ref{fig:whatif3} sweeps it.

\subsection{The advertising ceiling}
The 2030 ceiling decomposes as (Meta ads 2030 $\sim$Triang(330, 380, 430)) $+$ (Google ads 2030 $\sim$Triang(380, 450, 520)), both from 2026 levels~\cite{ref:meta_q2_2026b} at 10--15\%/yr, times an AI-attributable lift $\sim$Triang(0.08, 0.12, 0.18) anchored on Meta's disclosed 8\% click and 15.7\% conversion lifts and 12\% price-per-ad growth~\cite{ref:meta_q2_2026}, plus AI-native formats $\sim$Triang(10, 25, 60). Result: \$136B (P10--P90 \$111--164B), $\sigma_{\ln}=0.15$ versus 0.22 before. Meta's Advantage+ run-rate of \$75B~\cite{ref:meta_q2_2026} is not an uplift figure --- it is revenue flowing through an AI-managed product --- but it bounds the base on which any lift can act, and it is consistent with the range.

\subsection{The common adoption factor}
The factor that correlates the sectors ($\sigma_{\ln}=0.14$) is not one thing. It splits into an IT-budget cycle component --- enterprise IT spend growth of 8--10\%/yr with a well-forecast variance, $\sigma\approx0.05$~\cite{ref:gartner_it}\verify; an inference price-deflation component --- revenue is volume times price, and price per token for fixed capability has fallen $\sim$10$\times$/yr~\cite{ref:epoch_spending} while volume has risen faster, net $\sigma\approx0.08$; and a residual --- whether agentic systems work well enough to be trusted with whole workflows, and how fast enterprises reorganize around them --- $\sigma\approx0.10$. The first two are improvable with data and are largely observable by 2027. The third is the irreducible core of the whole debate, and it is the axis of Figure~\ref{fig:whatif1}.

\subsection{Facility opex}
The supply side's loosest input after the base was facility opex, at Triang(0.05, 0.08, 0.11) of in-service gross. It decomposes into power --- PUE $\sim$Triang(1.12, 1.22, 1.35) $\times$ 8{,}760\,h $\times$ utilization Triang(0.65, 0.80, 0.92) $\times$ tariff Triang(\$0.075, 0.10, 0.13)/kWh~\cite{ref:eia_epm}, divided by all-in capex of Triang(38, 45, 55) \$M per MW of IT load --- giving $\sim$1.9\% of gross per year; staffing and maintenance Triang(1.2, 2.0, 3.0)\%; property tax and insurance Triang(0.7, 1.0, 1.4)\%; network and bandwidth Triang(0.5, 1.0, 1.6)\%; other Triang(0.3, 0.6, 1.0)\%. Sum: 6.65\% (P10--P90 5.9--7.4\%), $\sigma_{\ln}=0.085$ versus 0.157. The central value falls as well as the spread --- a decomposition that moves in the build-out's favor, as roughly half of them should if the choices are not motivated.

\subsection{The assistant ceiling}
Microsoft disclosed 15M paid M365 Copilot seats in January 2026, 20M in April and more than 30M by July, with the largest single deal at 740{,}000 seats~\cite{ref:msft_copilot_seats,ref:q2_scorecard}. Three qualifications keep that number honest. A paid seat is a commercial commitment, not evidence of use: the 10M added in one quarter coincided with the May launch of the \$99 E7 suite that bundles Copilot into E5, the July base-price increases, and the introduction of usage-based billing alongside per-seat licensing~\cite{ref:ai2_copilot,ref:windowsforum_copilot}; Microsoft reports weekly engagement matching Outlook and Teams and satisfaction doubling over three quarters, but independent estimates put weekly active use at 20--30\% of assigned seats in some large deployments~\cite{ref:ai2_copilot}. The commercial base itself grew 6\% while Microsoft 365 commercial cloud revenue grew 17\%, so the growth is price and mix, not seats~\cite{ref:msft_8k_fy27,ref:fool_copilot}. The chain below therefore treats paid seats as the upper bound on adoption and carries an \emph{active share} of Triang(0.25, 0.45, 0.70) of paid seats as a secondary variable: the number a usage disclosure would settle, and the one that separates a bundle from a product. The 2030 ceiling chain is: paid seats across vendors mid-2026 Triang(35, 42, 50)M; annual seat growth Triang(1.3, 1.6, 2.0)$\times$ in 2027--28 decaying thereafter, hard-capped at 250M (about 55\% of the M365 and Workspace commercial base); price Triang(\$25, 30, 36)/seat/month; agentic and workflow consumption at Triang(1.2, 2.0, 3.0)$\times$ seat revenue; the whole capped at an AI share of Triang(3, 4, 5.5)\% of $\sim$\$6.5T 2030 IT spend, $\sim$\$270B. Result: \$150B (P10--P90 \$100--229B). This decomposition does not tighten --- $\sigma_{\ln}$ is 0.31 against 0.20 for a single judgment range --- because the seat-growth term is genuinely wide and a judgment range would understate it. That is the correct outcome: the assistant ceiling is the loose axis, and it now has a data-anchored floor and a budget-anchored hard cap. What the ceiling means: the most that enterprises could spend on AI assistants and the agentic workloads built on them in 2030, given seats, prices and IT budgets.

\subsection{Silicon life: the cutoffs}
Silicon life does not decompose further, but it is bounded on both sides by evidence: below by the thermal-acceleration physics of Section~\ref{sec:supply} (an economic life under 2.5 years would imply failure rates no operator has reported), above by the operators' own accounting lives of 5--6 years (an economic life beyond the accounting life would mean the companies are over-depreciating, which none has claimed). The training/inference fleet mix (Section~\ref{sec:supply}) gives an effective 4.2 years, inside those cutoffs.

\subsection{The feasibility cutoff: a budget that is itself being torn up}
A first-pass cutoff --- 8\% of total IT spend plus the advertising and consumer caps, $\sim$\$810B, drawn as a line --- would rest on an assumed share. A second pass would take Gartner's 2026 enterprise software (\$1.44T) and IT services (\$1.87T)~\cite{ref:gartner_it_jul2026}, grow them at 7--13\%/yr, and apply an AI share; that treats the pool as stable and Gartner as reliable, and neither holds. Gartner revised its 2026 IT-growth forecast from 10.8\% in February to 14.2\% in July~\cite{ref:gartner_it_feb2026,ref:gartner_it_jul2026} --- a 3.4-point swing in five months on the same year --- which is the precision such forecasts deserve, and the pool itself is being reorganized by the very technology the paper is pricing. Three things are happening to it at once.

\textbf{Seat-priced software is deflating.} Between January and April 2026 the enterprise software sector lost roughly \$2T of market value --- \$285B in one 48-hour stretch after Anthropic's Claude Cowork launch --- in what JPMorgan called the largest non-recessionary software drawdown in over thirty years~\cite{ref:saaspocalypse_ti,ref:saaspocalypse_nx}. Atlassian reported its first-ever decline in enterprise seat counts; Salesforce fell $\sim$30\% and Adobe $\sim$36\% year-to-date~\cite{ref:saaspocalypse_fc}. The market's reading is the paper's: if an agent does the work of ten people, the customer buys one seat, not ten. Gartner itself projects 40\% of enterprise applications carrying AI agents by end-2026, from under 5\% in 2025~\cite{ref:saaspocalypse_arbi}. Per-seat application software is $\sim$55\% of the software pool; the model gives it 2026--30 growth of Triang($-6$, 0, $+5$)\%/yr, against Triang(4, 8, 12)\% for infrastructure software, databases and security. Customers paying Microsoft for extended security updates to \emph{avoid} Windows 11 are the small-scale version of the same fact: incumbents are collecting rent on products their customers would rather not use, and rent is the first thing a budget cut removes.

\textbf{Buyers outside the US are displacing US vendors.} In April 2026 France's DINUM announced a national migration of government workstations from Windows to Linux; Schleswig-Holstein is mid-way through a full exit from Microsoft Office, Exchange and Windows; Denmark is following; France's national health fund is moving 80{,}000 staff to sovereign tools~\cite{ref:eu_sovereignty}. The trigger was not price but law: Microsoft's disabling of the International Criminal Court prosecutor's account on a US order made the CLOUD Act's reach concrete, and a 2026 survey of 1{,}818 IT decision-makers in Germany, France and the Nordics found the share actively implementing European software doubled from 20\% to 40\% in a year, with 87.5\% of German respondents favoring European products in critical infrastructure~\cite{ref:tuta_sovereignty}. Europe is $\sim$23\% of the global pool; the model moves Triang(10, 20, 35)\% of it away from US vendors by 2030. China is $\sim$10\% and is already 70--95\% domestic. The IT services pool, the largest single component, grows at Triang($-2$, 3, 7)\%/yr as AI removes consulting hours faster than implementation work adds them.

\textbf{The money moves from seats to solutions.} What the seat and services deflation does not do is disappear. The dollars are reallocated: to outcome-priced AI solutions, to agents that replace the seats, to the labs and integrators that sell them. The model puts the reallocated share of the accessible pool at Triang(12, 20, 32)\% by 2030, plus Triang(2, 4, 7)\% of genuinely new AI spend.

\textbf{Adding AI while raising the bill: the incumbent's dilemma.} The seat-priced vendor's arithmetic is general and comes first; one company is then used as an illustration, and only as that. On July 1, 2026 Microsoft raised Microsoft 365 commercial prices globally --- up to 43\% on some plans, E3 to about \$39 per user per month, Windows Enterprise up --- bundling Copilot Chat and Security Copilot into the suites in the name of ``value expansion''~\cite{ref:m365_price_wl,ref:m365_price_uscloud,ref:m365_price_ow}; standalone Copilot seats, which customers had shown little appetite for, were left unpriced and folded into the base. That is the seat-priced incumbent's only move, and Figure~\ref{fig:incumbent} shows its arithmetic. For a vendor whose seats are falling to agents and to sovereignty exits, revenue is seats times price; a price rise of $p$ per year with a seat-loss elasticity $\epsilon$ (seats lost per unit of cumulative price increase) recovers ground only while $(1+p)^4\,[1-\epsilon((1+p)^4-1)]$ exceeds the seat decline. At an elasticity of 0.35 and seat losses of 3\%/yr, an 8\%/yr price rise roughly holds nominal revenue flat and no more. \emph{Illustration --- Microsoft 365 commercial.} Taking widely quoted but unverified inputs\verify --- Triang(420, 450, 480)M paid seats at Triang(\$19, 21, 23)/month, $\sim$\$113B in 2026 --- with a seat trend that is still positive today --- Microsoft's own investor metrics show M365 commercial seats growing 6--7\% a year through fiscal 2026~\cite{ref:msft_8k_fy27}, so the model takes Triang($-4$, 0, $+4$)\%/yr over 2026--30, a deceleration from that pace as agents and sovereignty exits bite rather than an immediate decline --- cumulative sovereignty exit Triang(3, 8, 15)\% by 2030, price growth Triang(5, 8, 12)\%/yr and elasticity Triang(0.15, 0.35, 0.60), the base suite revenue in 2030 comes out around \$110--140B on $\sim$350--420M seats: near or modestly above 2026 on these inputs, with the disclosed seat growth doing the work. The observed +6--7\% is the number the paper's seat-deflation argument has to overcome, and it has not yet. Copilot attach is what would change it: if paid attach on the remaining seats reaches 12--35\% by 2030 (from $\sim$6.5\% today, 30M of $\sim$460M) with agentic consumption behind it, total Microsoft 365 revenue is roughly \$130--170B --- above 2026, but short of an 8\%/yr trend in most draws. Microsoft is used here because it discloses more than its peers, not because it is uniquely exposed; Salesforce, Adobe, ServiceNow and SAP face the same equation with less public data. The lesson generalizes to every seat-priced vendor: raising the bill for AI features does not offset seat loss; only converting seats to paid AI consumption does, and that conversion is the same adoption variable this paper cannot pin from 2026. It is why the assistant-sector ceiling and the residual factor, not the aggregate budget, are where the argument lives.

\begin{figure}[htbp]\centering\includegraphics[width=\textwidth]{f21_incumbent.pdf}
\caption{Left: seat revenue in 2030 relative to 2026 for a seat-priced vendor, by seat decline and price growth, at a seat-loss elasticity of 0.35. Right: Microsoft 365 commercial as an illustration on unverified inputs --- the base suite is more likely than not below 2026 revenue by 2030; the total holds up only if Copilot attach rises fivefold.}\label{fig:incumbent}\end{figure}

\textbf{Result.} The 2030 pool is \$3.75T (P10--P90 \$3.5--4.0T), not the \$4.85T that 7--13\% growth would give; the part accessible to US vendors is \$3.25T; the AI share of it is $\sim$25\% (20--32\%); the enterprise AI cap is \$825B (\$650B--\$1.03T); with advertising (\$140--200B) and consumer (\$80--130B) caps the total feasibility band is \$1.10T (P10--P90 \$920B--\$1.31T). That is almost exactly the band a stable, growing pool with a 15\% AI share would give --- the two errors cancel --- but the composition is entirely different, and the composition is the point. The feasibility of the aggregate is unchanged; what changes is who receives the dollar. Seat-priced incumbents lose it; European and Chinese providers keep more of it at home (the displacement enters the demand chains as a ramp on the sovereign-exposed sectors, Section~\ref{sec:bottomup}); and the largest single loser is the company that is both the largest seat-priced vendor and the most exposed to the sovereignty exit --- Microsoft, whose AI receipts grow while its conventional Office, Windows and Azure base is being cut from both sides. The band is drawn as a band, not a line, in Figure~\ref{fig:probmap}.

\textbf{Two secondary variables.} The reallocation share (12--32\% of the accessible pool moving to AI solutions) and the seat-priced SaaS growth rate ($-6$ to $+5$\%/yr) are the loosest terms in the band and neither can be pinned from 2026 data: the first will be visible in Gartner's and the vendors' 2027 category reporting, the second in seat-count disclosures over the next four quarters. Rather than let them hide inside the band they are treated as \emph{secondary variables} --- held at their central ranges in the base case and swept explicitly in Figure~\ref{fig:secondary}. The reading is that the plateau's 50\% contour ($\sim$\$950B) stays inside the band unless the reallocation share falls below $\sim$12\% \emph{and} seat software contracts at 6\%/yr with heavy European displacement; the composition changes far more than the total.

\begin{figure}[htbp]\centering\includegraphics[width=\textwidth]{f22_secondary.pdf}
\caption{The two secondary variables in the feasibility band. Left: the cutoff versus the share of the accessible pool reallocated to AI, for three seat-software growth rates. Right: the cutoff versus seat-software growth, for three levels of European displacement of US vendors. The dash-dot line is the plateau's 50\%-probability contour.}\label{fig:secondary}\end{figure}

\subsection{The non-AI capex baseline}
The \$162B the five companies spent in 2022~\cite{ref:epoch_capex} decomposes by company --- Amazon fulfillment and transport $\sim$\$36B, AWS non-AI $\sim$\$28B, Microsoft $\sim$\$24B, Alphabet $\sim$\$31B, Meta $\sim$\$32B, Oracle $\sim$\$4B\verify. Growing that figure forward at an assumed 0--8\%/yr would give $\sim$\$190B for 2026. The filings constrain it lower (below); the by-company split is what the constraint is applied to.

\subsection{Energization lag, by component}
A single lag Triang(0.3, 0.5, 0.7) hid two different physical processes. Accelerators ship and rack within one to three quarters of purchase ($\ell_{\text{si}}\sim$Triang(0.5, 0.7, 0.9) in the spend year); shell, power and cooling take twelve to twenty-four months ($\ell_{\text{sh}}\sim$Triang(0.2, 0.4, 0.6)). The effective lag is 0.56 with $\sigma_{\ln}=0.11$ versus 0.17, and its variance share in Table~\ref{tab:varbudget} fell from 8\% to under 1\%.

\subsection{Silicon life, by fleet}
A single economic life hid the training/inference split. Training clusters are retired when a successor's throughput per watt makes the power slot worth more than the silicon (Section~\ref{sec:supply}): Triang(2.5, 3.5, 4.5) years. Inference fleets serve fixed models and run longer: Triang(4, 5, 6). At a training share of Triang(0.35, 0.45, 0.55) the fleet's effective life (harmonic mean, since depreciation is $1/L$) is 4.2 years (P10--P90 3.75--4.6), $\sigma_{\ln}=0.077$ versus 0.124 for the undecomposed Triang(3, 4, 5.5).

\subsection{The consumer ceiling}
Gemini's app reached 950M monthly users and Meta AI 1.2B by mid-2026~\cite{ref:consumer_mau}; with ChatGPT and others, unique consumer users are Triang(1.3, 1.6, 2.0)B. Paid conversion Triang(4, 6, 8)\% at Triang(\$20, 23, 26)/month gives \$27B of subscriptions; free-tier advertising on Triang(0.8, 1.0, 1.3)B ad-eligible users at Triang(\$15, 25, 40)/user-year gives \$27B; premium tiers Triang(5, 10, 15). Ceiling \$64B (P10--P90 \$55--75B), $\sigma_{\ln}=0.12$ versus 0.18, and lower than the earlier \$80B mode.

\subsection{Six further decompositions, and a grid variable}
\textbf{Hyperscaler AI services excluding labs.} Microsoft's $>$\$37B AI run-rate~\cite{ref:yahoo2026} decomposes into Azure hosting of OpenAI ($\sim$\$13B, already in the OpenAI channel~\cite{ref:valueadd_openai}), M365 Copilot seats (20M $\times$ \$25--36 $\times$ 12 $\approx$ \$7B~\cite{ref:msft_copilot_seats}), GitHub Copilot ($\sim$\$2B), and other Azure AI services (remainder). AWS AI ($>$\$15B) less Bedrock pass-through, and Google Cloud AI less Vertex pass-through, add \$8--12B each. Calendar-2026 hyperscaler ex-lab revenue: \$39B (P10--P90 \$35--42B), $\sigma_{\ln}=0.067$ versus 0.118. The 2026 base becomes Triang(163, 179, 194).

\textbf{The lab channel's infrastructure share, by lab.} OpenAI's compute purchases --- $\sim$\$13B on Azure plus \$6--12B with Oracle and CoreWeave --- against $\sim$\$29B of bookings give 0.7--1.0; Anthropic's margin is undisclosed but bracketed by reported 2024--25 figures\verify at a compute share of Triang(0.45, 0.55, 0.70); other labs Triang(0.6, 0.7, 0.85). Revenue-weighted: $\theta_{\text{lab}}=0.635$ (P10--P90 0.59--0.69), $\sigma_{\ln}=0.058$ versus 0.108. Overall $\theta$ becomes Triang(0.77, 0.815, 0.86).

\textbf{The non-AI baseline, constrained by 2023.} Total capex was flat at \$161B in 2023 while AI purchases of \$20--40B began~\cite{ref:epoch_capex}, so the non-AI baseline must have \emph{fallen} to $\sim$\$130B in 2023 (Amazon cut fulfillment capex that year). Growing from there at Triang(2, 6, 10)\%/yr gives \$156B for 2026 (P10--P90 \$145--168B) and $\sim$\$197B by 2030 --- lower than the $\sim$\$190B an assumed growth from 2022 would give. The AI-attributable stock is correspondingly larger. A constraint from a filing beats a growth assumption, whichever way it cuts.

\textbf{Silicon share, from Nvidia.} Nvidia's \$300B data-center run-rate~\cite{ref:nvda_q1fy27} at a Big-5 share of Triang(0.45, 0.55, 0.65) is \$165B; server integration, custom accelerators (TPU, Trainium, MTIA), memory, networking and storage add roughly as much again. Against AI-attributable capex of $\sim$\$620B --- the share must be taken against AI capex, not the \$780B total, or it comes out implausibly low --- the silicon share is Triang(0.45, 0.52, 0.60).

\textbf{Cost of capital, by CAPM.} Risk-free Triang(4.0, 4.3, 4.6)\%, equity premium Triang(4.5, 5.0, 5.5)\%, levered beta Triang(1.1, 1.2, 1.35), pre-tax debt cost Triang(5.5, 6.5, 7.5)\%, equity weight Triang(0.75, 0.82, 0.90), tax 21\%~\cite{ref:damodaran2026}: $r=9.45\%$ (P10--P90 9.0--9.9\%), $\sigma_{\ln}=0.038$ versus 0.058. The mode fell a point; the carry fell $\sim$\$10B.

\textbf{Grid energization capacity.} Capital can only earn once it is \emph{energized}, and the grid limits how much can be connected in a year. Global AI data-center interconnection capacity of Triang(22, 30, 40) GW/yr of IT load~\cite{ref:pjm_queue,ref:eirgrid_cru}\verify at Triang(\$38, 45, 55)M per MW is \$1.15--1.68T/yr of energizable capex. Capex above it is deferred a year and carries return but no depreciation or opex. It does not bind on a plateau; it binds in $\sim$5\% of arms-race draws, where it \emph{lowers} the near-term carry by delaying depreciation while the capital sits idle. It is in the model because a physical ceiling should be, and because it becomes decisive if capex exceeds $\sim$\$1.3T/yr.

\subsection{Readiness: consent, water, and siting}
A funded, grid-connected project must also be \emph{allowed} to be built, and the record says that constraint now binds. Data Center Watch counts \$156B of US data-center projects blocked or delayed by local opposition in 2025 and \$130B in the first quarter of 2026 alone --- 75 projects, roughly the whole of 2025 in three months~\cite{ref:dcw_q1_2026,ref:tomshw_dcw}. Active opposition groups more than doubled to 833 across 49 states; more than 300 state bills were filed in six weeks; statewide moratoria were proposed in 14 states and New York enacted a halt on large new facilities in July 2026~\cite{ref:dcw_q1_2026,ref:kucoin_dcw}. Cancellations rose from 6 in 2024 to 25 in 2025 and more than 20 in early 2026, including a \$14B Arizona site withdrawn and a \$24.7B Virginia site in litigation~\cite{ref:kucoin_dcw}. Ipsos polling moved from about half of Americans opposing a data center near them in late 2025 to roughly 70\% by mid-2026~\cite{ref:tomshw_dcw,ref:forbes_consent}. The stated grounds are the ones this paper's supply section already prices in other forms: electricity and ratepayer bills, water, noise, diesel backup, and opaque permitting. On bills specifically the record cuts the other way so far: Watten, Bistline and Blanford find that data-center load \emph{lowered} average residential rates by about 6\% over 2019--24 by spreading the grid's fixed costs~\cite{ref:watten2026}; the opposition is nonetheless real and growing, and the deferral variable is calibrated on what communities do, not on what rates did. Europe reached this point first: Dublin~\cite{ref:eirgrid_cru} and Amsterdam~\cite{ref:amsterdam} imposed connection and zoning limits in 2021--22, Frankfurt and Paris face grid and land constraints, and the Netherlands has restricted hyperscale sites nationally. China is the exception: siting is a state decision, and the \emph{East Data, West Computing} program directs capacity to the western provinces where power and land are cheap and local objection is not a factor. The deferral variable below is calibrated on US data and applied to the whole Big-5 stock, which is conservative for the Chinese portion (small for these five companies) and roughly right for Europe.

\textbf{Water, as a Fermi chain.} A 2030 AI fleet of Triang(60, 80, 110) GW at Triang(0.65, 0.80, 0.92) utilization draws $\sim$570 TWh/yr; at a water-usage effectiveness of Triang(0.3, 1.0, 1.8) L/kWh and an evaporative share of Triang(0.4, 0.6, 0.8), direct consumptive water is $\sim$90 billion gallons/yr (P10--P90 53--141), before the 3--4$\times$ larger indirect draw of the power generation behind it. Why is the water not simply recycled? Evaporative cooling works \emph{by} evaporating it: the heat leaves as vapor, and that water is consumed, not returned. Closed-loop and dry cooling return the water but reject the heat to air instead, which takes 5--10\% more power. For scale: US freshwater withdrawals are about 320 billion gallons per \emph{day}, so 90 billion gallons per \emph{year} is under 0.1\% of the national total, and far less than agriculture or thermoelectric power plants. The issue is local: a 500\,MW campus draws like a town of 50{,}000 people, from one aquifer, in the dry regions where land and power are cheapest. It is why Arizona, Aragon and the Dutch Randstad refuse evaporative cooling. Operators are not exaggerating their efficiency work --- WUE has fallen steadily --- but the concentration, not the total, is what communities object to.

\textbf{Three variables.} (i) A siting and consent deferral $f_{\text{delay}}\sim$Triang(0.10, 0.20, 0.35): the fraction of each cohort from 2026 pushed a year by permitting, litigation or relocation, treated like the grid deferral --- it carries return but not depreciation. The mode is set from the ratio of disrupted project value to annual pipeline (\$130B/quarter against a US pipeline of order \$1T), discounted because most delayed projects eventually proceed elsewhere. (ii) A shell cost uplift $\kappa_{\text{site}}\sim$Triang(1.02, 1.05, 1.09) for community-benefit agreements, water mitigation, dry cooling, private generation and relocation to friendlier jurisdictions~\cite{ref:forbes_consent}. (iii) A cooling opex uplift $u\sim$Triang(0, 0.03, 0.06) where evaporative cooling is refused.

\textbf{Effect.} The three variables shift the carry rather than remove it: 2027 falls by $\sim$\$20B (deferred cohorts), 2030 rises by $\sim$\$15B (deferred cohorts arriving with higher shell cost and opex). Break-even $g^{*}$ is unchanged at 49\%; the bottom-up 2030 deficit is $\sim$\$15B larger and $z$ about 0.2 more negative than without the friction; rationalization's 2030 P(profit) is roughly ten points lower, because the cut path's late cohorts are the ones deferred into 2029--30. Readiness friction is, on net, a bear variable for the carry test and a bear variable for the revenue test, since the capacity it delays cannot serve demand in the year it was meant to.

\subsection{What remains, and the what-if curves}
After nineteen decompositions the top-down root-sum-square falls from 286 to 104 (Table~\ref{tab:varbudget}), and the residual uncertainty sits in three places: the residual adoption/capability factor, the assistant ceiling, and the labs' compute-cost share. Each has a cutoff of its own --- the IT-budget share for the first two, the labs' reported gross margins for the third --- and Section~\ref{sec:overlap} maps the probability of profit over the two that matter most.

\begin{figure}[htbp]\centering\includegraphics[width=\textwidth]{f13_whatif_mcap.pdf}
\caption{What-if 1. The 2030 $z$-score and P(deficit) as the residual adoption/capability factor scales bottom-up receipts from 0.6 to 1.6, by capex path. The plateau path crosses $z=0$ only above $\sim$1.5 --- receipts 50\% above the sector chains' central estimate.}\label{fig:whatif1}\end{figure}

\begin{figure}[htbp]\centering\includegraphics[width=0.8\textwidth]{f14_whatif_ceilings.pdf}
\caption{What-if 2. Median 2030 deficit (Plateau) across the two ceilings that could close the gap. The dot is the paper's modes. No combination in the plotted range reaches break-even; even a \$300B advertising ceiling with a \$600B assistant ceiling leaves a median deficit near $-\$60$B, because by 2030 the logistic ramps have reached only 60--80\% of ceiling and assistant receipts are $\sim$60\% of assistant revenue.}\label{fig:whatif2}\end{figure}

\begin{figure}[htbp]\centering\includegraphics[width=\textwidth]{f15_whatif_theta_base.pdf}
\caption{What-if 3. Left: 2030 deficit versus the labs' compute-cost share (the only loose part of $\theta$). Right: top-down 2030 deficit at $g=49\%$ versus the 2026 base across and beyond its decomposed range.}\label{fig:whatif3}\end{figure}

The sign survives each residual variable moved alone across its plausible range. What it does not survive is the residual factor above $\sim$1.5 on a plateau --- the sector chains being wrong by half in the same direction at once. No pair of ceilings does it alone (Figure~\ref{fig:whatif2}); it takes ceilings, ramp rates and cloud-capture shares moving together, which is a single statement about agentic adoption in the enterprise. That is where this paper's argument and the bull case actually disagree.

\section{Digital Sovereignty and Grid Friction}
\textbf{Capital duplication ($\kappa_{\text{sov}}$).} Data-residency regimes (GDPR Chapter~V, the EU AI Act, BSI C5, India's DPDP Act)~\cite{ref:gdpr,ref:aiact,ref:dpdp} require in-jurisdiction processing for large classes of data, forcing smaller, under-utilized facilities in Frankfurt, Paris, Dublin, London, Tokyo and Riyadh in place of a few gigawatt sites in cheap-power regions. We model the resulting carry multiplier as $\kappa_{\text{sov}}\sim$Triang(1.1, 1.2, 1.3). This is an engineering judgment about utilization and scale, not a published measurement.

\textbf{Demand carve-out ($\gamma_{\text{sov}}$).} Procurement preferences for domestic champions (Mistral, Falcon, Bhashini) and air-gapped defense enclaves remove a fraction $\gamma_{\text{sov}}\sim$Triang(0.15, 0.25, 0.35) of the sovereign-exposed sectors (office assistants, CX, financial, healthcare, legal, public sector, science) from US hyperscalers.

\textbf{Grid queues.} Multi-year interconnection backlogs in PJM~\cite{ref:pjm_queue}, Dublin's large-user connection constraint through 2028~\cite{ref:eirgrid_cru}, and Amsterdam's establishment restrictions~\cite{ref:amsterdam} cap the physical velocity of deployment. Capital committed to sites that cannot energize is cash drag: it is in the installed base but produces no revenue. We do not model this separately; it cuts against the bull case.

\textbf{Effect.} Applying $\kappa=1.2$, $\gamma=0.25$ to the bottom-up model roughly doubles the 2030 plateau deficit (Table~\ref{tab:bu_results}); in the top-down sweep it raises the plateau break-even for 2030 from 49\% to $\sim$66\%; the siting friction of Section~\ref{sec:decomp} is the domestic form of the same constraint.

\section{Circular Capital Flows}\label{sec:circ}
Reported cloud revenue is not all exogenous demand. Where a hyperscaler's equity investment in a lab returns as that lab's cloud spend, the revenue is a round trip. Modeling the ecosystem as a directed graph, the circular revenue on a cycle is bounded by its minimum edge weight $\omega(e)$ (the dollars flowing along that edge), $F(C)=\min_{e\in C}\omega(e)$.
\begin{table}[htbp]
\centering
\small
\caption{Circular versus merely non-recurring flows (annual, order of magnitude)}
\label{tab:circular}
\begin{tabularx}{\textwidth}{L{4.3cm}L{1.8cm}L{2.2cm}X}
\toprule
\textbf{Flow} & \textbf{Class} & \textbf{Order} & \textbf{Basis} \\
\midrule
Microsoft $\leftrightarrow$ OpenAI & Circular & $\sim$\$4B & Equity/credits returning as Azure consumption~\cite{ref:hudson_policy}\verify \\
Nvidia $\leftrightarrow$ CoreWeave & Circular & $\sim$\$3B & Equity and capacity backstops returning as GPU orders~\cite{ref:coreweave_s1}\verify \\
Amazon/Alphabet $\leftrightarrow$ Anthropic & Circular & $\sim$\$3B & Convertible notes returning as Trainium/TPU consumption~\cite{ref:techcrunch_anthropic}\verify \\
Multi-party (Oracle, CoreWeave) & Partly circular & $\sim$\$2--5B & Colocation and capacity swaps; edge weights not disclosed \\
\midrule
\textbf{Genuinely circular} & & \textbf{$\sim$\$1.2$\times$10$^{10}$} & Triang(8, 12, 18) in the model \\
\midrule
VC-funded startup API spend & Non-recurring & $\sim$\$1.5--2$\times$10$^{10}$ & Real cash, but $\sim$80\% of the cohort fails within 2--3 years\verify~\cite{ref:cbinsights}; modeled as the low tail of enterprise demand, not as circular \\
Oracle RPO backlog & Backlog & \$3$\times$10$^{11}$+ & A contracted future obligation to non-investment-grade counterparties~\cite{ref:oracle_10q}; a collectability risk, not a loop \\
\bottomrule
\end{tabularx}
\end{table}
Only the first three rows are round trips in the strict sense; the rest are real but fragile. The circularity that matters in 2026 is larger and sits outside this table: the $\sim$\$190B of lab equity raised this year, a large share of which is spent on hyperscaler compute (Section~\ref{sec:demand}).

\textbf{Off-balance-sheet structures.} As in the 1999--2001 telecom buildout, capacity is increasingly financed through SPVs with infrastructure funds (Blackstone, Brookfield, Blue Owl) and vendor-backstopped debt. This does not reduce the carry; it converts equity capex into debt service with DSCR tripwires that accelerate distressed liquidation when revenue undershoots.


\section{Enterprise ROI and the Slack Trap}\label{sec:slack}
\subsection{What enterprises must realize}
Take 2028 under Plateau: carry $\sim$\$570B. Enterprise buyers supply roughly 70\% of AI revenue, so they must pay $\sim$\$400B to hyperscalers and labs. Every dollar of software or token spend carries $\mu_{\text{integ}}\sim$Triang(1.0, 1.5, 2.5) dollars of integration, so total enterprise outlay is $\sim$\$1.0T, and with a required benefit margin $m\approx0.2$ over cost for an unproven technology~\cite{ref:koller2020,ref:brealey2023} --- every firm applies a hurdle to new spending, and the hurdle is higher, not lower, for a technology whose payback is uncertain --- enterprises must realize $B_{\text{req}}\approx\$1.2\times10^{12}$/yr of auditable benefit by 2028.

\subsection{The labor-substitution floor}
Against a global knowledge-worker payroll of $\sim$\$30T (1.05B workers at $\sim$\$28.5k fully loaded) and a US pool of $\sim$\$9T (62M at $\sim$\$145k)~\cite{ref:bls_oes,ref:bls_ecec}, $B_{\text{req}}$ is $\sim$4\% of global payroll or $\sim$13\% of US payroll if the US bore it alone. Benefits need not be layoffs --- revenue growth, avoided hiring through attrition, faster cycle times all count --- but they must be auditable cash flow.

\subsection{Four reasons realized yield is a fraction of nominal savings}
\begin{enumerate}
    \item \textbf{Indivisibility.} Tools save fragments of tasks; thirty minutes a day across six workflows does not aggregate across ten employees into one dismissal. Avoided hiring through attrition is the realistic path, and it is slow.
    \item \textbf{Verification.} While liability (SEC, FINRA, HIPAA, professional conduct) rests with human officers, human review is an irreducible fixed cost; labor shifts from generation to verification.
    \item \textbf{Parkinson--Jevons rebound.} Cheaper output inflates the demand for output: three proposal variants, sixty-page reports, more meetings to digest synthetic overflow. Gross volume rises; cash flow does not.
    \item \textbf{The manager's dilemma.} Compensation, prestige and political capital are indexed to headcount. Departments report ``bandwidth expansion,'' not surplus payroll.
\end{enumerate}
The lag has precedent, and the precedent is measured. Paul David's study of the electric dynamo found that factories took roughly forty years, from the 1880s to the 1920s, to reorganize around electric power before productivity responded, because the gain came from rebuilding the factory floor, not from the motor~\cite{ref:david1990}. The personal computer arrived in 1981; US labor productivity growth ran at 1.4\%/yr from 1973 to 1995 --- Solow's 1987 paradox --- and accelerated to 2.9\%/yr only in 1995--2004, a fifteen-year lag that Jorgenson, Ho and Stiroh attribute to information technology once complementary investment had been made~\cite{ref:bls_productivity,ref:jorgenson2008}; Brynjolfsson and Hitt's firm-level work found returns to computers appearing over five to seven years and only where organization changed with them~\cite{ref:brynjolfsson_hitt2000}. The internet's contribution followed the same shape. The present wave is, by the aggregate data, at the start of that curve: the Census Bureau's business survey put AI use in production at roughly a tenth of US firms in 2025\verify~\cite{ref:census_btos}, and no acceleration is yet visible in the productivity statistics. What the evidence says at the firm level so far: an MIT study of enterprise generative-AI deployments in 2025 found that about 95\% of pilots produced no measurable profit-and-loss impact, the gains concentrating in a small minority that reorganized workflows around the tool~\cite{ref:mit_nanda}; McKinsey's 2025 survey found most firms using AI somewhere but only $\sim$1\% describing their deployments as mature~\cite{ref:mckinsey_ai2025}. Where controlled studies exist the task-level gains are real --- 14\% more issues resolved per hour in a call center, larger for novices~\cite{ref:brynjolfsson_callcenter}; 26\% more completed tasks with an AI coding assistant in a randomized trial~\cite{ref:cui_copilot}; 12\% more tasks and 25\% faster in a consulting-task experiment~\cite{ref:dellacqua_bcg} --- and none of those studies reports the headcount that was removed as a result. The gap between task-level gains and firm-level cash is the slack trap's footprint, and it is the same footprint the ERP and office-software waves left: fractional time savings on the desktop, no matching line in the payroll. Companies that do cut do it wholesale, in restructurings attributed to AI in the press release, and those show up in the labor data as ordinary restructurings; the paper's 1.5\% displacement marker (Section~\ref{sec:falsify}) is what would distinguish the two. We formalize realized benefit as
\begin{equation}
B_{\text{realized}}=W\,\tau\,(1-\sigma_{\text{slack}})(1-\nu_{\text{verify}})-C_{\text{severance}},
\end{equation}
where $W$ is the payroll pool, $\tau\sim$Triang(0.10, 0.15, 0.25) the nominal share of work time AI saves, $\sigma_{\text{slack}}\approx0.65$ the share of that saving absorbed as slack, $\nu_{\text{verify}}\approx0.20$ the share consumed by verification, and $C_{\text{severance}}$ one-time separation costs. These two coefficients are the paper's most consequential and least measured assumptions; they are judgments consistent with the J-curve literature~\cite{ref:brynjolfsson2021}, not measurements. At the modes, realized yield is $0.28\,\tau$: a 15\% time saving converts to $\sim$4\% of payroll --- the global requirement, but only over a full J-curve cycle of 8--12 years, against silicon lives of 3--5.

\section{Last-Mile Services and Value Capture}\label{sec:lastmile}
\subsection{Where the surplus goes}
The \$2.65T professional-services pool (consulting $\sim$\$1.05T, accounting/audit $\sim$\$0.65T, legal $\sim$\$0.95T)\verify faces compression of billable hours (tax prep, audit workpapers, contract review) that puts of order 30\% of billings ($\sim$\$0.8T) under deflation. Outcome-priced last-mile operators capturing 10--18\% of realized savings under gain-share terms address a \$80--140B/yr market, at high gross margins, using small fine-tuned open-weight models on modest hardware. Upstream token intensity of 1.5--3\% leaves the hyperscalers \$1--4B of that pool.

\subsection{The 10--50:1 disconnect}
Contact centers illustrate the ratio. Global contact-center labor is $\sim$\$350B. A plausible 2026--27 displacement is 5--10\%, or \$20--35B/yr (a figure of \$80B, sometimes quoted, would imply a quarter of the global pool already gone, which the labor data do not show). The RAG and speech inference behind that displacement consumes \$1--2B/yr of tokens. The ratio is 10--35:1; Nordhaus's economy-wide 2.2\% innovator share~\cite{ref:nordhaus2004} corresponds to $\sim$45:1. Freight routing and travel administration show the same pattern but are mostly classical optimization, not AI tokens, and we no longer count them.

\begin{table}[htbp]
\centering\small
\caption{The value-capture disconnect: realized savings versus upstream tokens (order of magnitude)}
\label{tab:disconnect}
\begin{tabularx}{\textwidth}{lXll}
\toprule
\textbf{Workflow} & \textbf{Billing base} & \textbf{Realized savings, 2026--27} & \textbf{Upstream cloud tokens} \\
\midrule
Contact centers \& support & $\sim$\$350B global labor & \$20--35B/yr & \$1--2B/yr (RAG, speech) \\
Corporate travel \& admin & $\sim$\$1.4T business travel & \$10--30B/yr overhead & Fractional (database triggers) \\
Tax prep, audit workpapers, contract review & $\sim$\$0.8T of billings under deflation & \$40--80B/yr by 2028 & \$1--4B/yr (small fine-tuned models) \\
Freight \& logistics routing & $\sim$\$2.6T US logistics & \$50--100B/yr & Near zero (classical optimization, not LLM tokens; excluded from AI demand) \\
\bottomrule
\end{tabularx}
\end{table}

\subsection{The talent trap and the timing chasm}
An in-house MLOps and agentic pod costs \$3.5--5M/yr, and most internal proofs-of-concept do not reach production (surveys put the figure above 70--80\%\verify). CFOs are shifting to outcome-aligned external operators. The cycle --- failed internal experiment, procurement, compliance validation --- delays material revenue by 3--5 years, against a silicon life of 3--5 years. That is the timing chasm in one sentence.

\subsection{Reconciling the slack trap with the leakage story}
These are not contradictory; they dominate in different places. In high-liability knowledge work (legal, finance, medicine, engineering) the slack trap dominates: savings are real but do not reach cash flow. In high-volume, low-liability process work (contact centers, claims intake, document routing) savings do reach cash flow, and there the leakage story dominates: the buyer and the integrator keep it. Neither regime routes much to the cloud meter.


\section{Hardware Cycle and Commoditization}
\begin{itemize}
    \item \textbf{Precision shifts.} FP4 microscaling delivers several times the throughput per watt of FP8/FP16 generations, and ternary-weight architectures remove floating-point multiply from the critical path. In power-constrained facilities this strands the prior generation regardless of its silicon health.
    \item \textbf{HBM bullwhip.} Lower-precision inference cuts memory-bandwidth demand per token. Capacity additions by SK Hynix, Samsung and Micron sized to 2025 demand risk oversupply by 2027--28.
    \item \textbf{Open weights as a deflationary wedge.} Llama, Mistral, DeepSeek and Nvidia's own open models commoditize the model layer. When a 70B-class model runs on a workstation, API pricing converges toward the electricity floor, removing the platform margin the hurdle needs. The low-cost invasion is already under way: Chinese open-weight models trained at a fraction of frontier cost are run on-premise by Western enterprises for routine work even where cloud access to Chinese providers is restricted or prohibited, because the weights, once published, are not a service that can be blocked. Distilled and quantized derivatives of those models are what the last-mile operators of Section~\ref{sec:lastmile} deploy. \emph{Microscaling} --- a block floating-point scheme in which groups of values share one exponent, standardized by the Open Compute Project in 2023 --- is what makes FP4 usable for inference and is expected to be adopted fastest where hardware is scarcest.
    \item \textbf{Edge and on-premise.} Two different things reduce cloud receipts: inference on the user's own device (phones and PCs with NPUs), and inference on a customer's own servers (on-premise clusters running open-weight models). The paper's cloud-capture share $\theta_s$ combines both; neither reaches a hyperscaler meter.
    \item \textbf{Quantum computing} is not modeled. Gate rates, error correction and the absence of scalable quantum memory keep it outside the 2026--2035 window for commercial inference or training; it is mentioned only to say so.
\end{itemize}

\section{Geopolitical Divergence}
\begin{itemize}
    \item \textbf{United States:} carries most of the global hyperscale debt and capex; industrial power \$0.10--0.14/kWh in PJM and rising; 5--7 year interconnection queues~\cite{ref:eia_epm,ref:pjm_queue}. No federal carbon price; California's cap-and-trade and the northeastern RGGI scheme cover a minority of load.
    \item \textbf{Europe (FLAPD):} power \$0.145--0.195/kWh, ETS carbon at $\sim$EUR~70/t, Dublin frozen to 2028~\cite{ref:eirgrid_cru}, Amsterdam capped~\cite{ref:amsterdam}.
    \item \textbf{China:} interconnection is a planning decision rather than a queue, and siting objection is not a constraint; curtailed western renewables at \$0.02--0.035/kWh~\verify under \emph{East Data, West Computing}; state-subsidized training for industrial and humanoid robotics; open-weight export. Embodied AI is bandwidth-heavy: a four-camera robot at 30 fps and $\sim$60 spatial tokens per frame produces $\sim$2.7$\times$10$^{7}$ tokens per camera-hour, which favors cheap power over cheap chips.
    \item \textbf{Middle powers:} India (DPI), the Gulf, Korea --- sovereign models on stranded power and open weights, bypassing US hyperscalers for domestic demand.
    \item \textbf{Atoms.} Transformers on 36--48 month lead times, electronic-grade polysilicon, copper and neon set the deployment velocity; integrated supply chains with cheap state power hold a structural advantage over permitting-bound Western sites.
\end{itemize}


\section{Economic Depreciation, Utilization, and the Commodity Floor}\label{sec:commodity}
Three points raised by a reviewer with a commodities background~\cite{ref:mehta_2026} change how the supply side should be read, and each is now in the model as a secondary variable.

\subsection{Economic versus accounting depreciation}
The carry in Section~\ref{sec:supply} depreciates silicon straight-line over an effective 4.2 years --- 20--30\%/yr, close to the operators' accounting. The economic view is different: an installed accelerator is worth what it can earn against the newest chip, and the marginal cost of inference per token on the newest chip falls 50--90\% with each generation. On that view the value of a cohort follows a declining balance, $\text{NBV}_t=G(1-\delta)^{\text{age}}$, with $\delta\sim$Triang(0.5, 0.7, 0.9) per year~\cite{ref:mehta_2026,ref:epoch_spending}, and most of a cohort's value is gone within two years whether or not the silicon still runs --- the direction Amazon's own 2025 shortening of server lives points~\cite{ref:amzn_10k_life}. Two consequences. First, the carry is front-loaded: at $\delta=0.7$ the 2027 carry is $\sim$\$415B rather than \$390B and the 2030 carry $\sim$\$650B rather than \$725B, because value written off early is not written off later (Figure~\ref{fig:econdep}, left). The annual deficits from 2027 are similar in size but flatter; what changes is timing, and the timing runs against the operators, since the losses fall in the years the paper already calls a sunk deficit. Second, and more immediately, the silicon bought in 2023--26 is already impaired on the economic schedule: its straight-line net book value at end-2026 is $\sim$\$330B, its economic value \$75--230B, an unrecognized gap of \$100B ($\delta=0.5$), \$175B ($\delta=0.7$) or \$250B ($\delta=0.9$). That gap is what ASC 360-10's fair-value step would recognize if the recoverability test failed; it is the mechanism, not the useful-life note, by which the impairment wave of Section~\ref{sec:falsify} would arrive. $\delta$ is a secondary variable, and the evidence on it currently points both ways: the cost-per-token series that Epoch AI maintains falls at the rates assumed, while Silicon Data's rental index has H100 and A100 rental prices \emph{rising} through 2026 (\$2.63 and \$1.59 per GPU-hour in September), with residual values inferred from them rising too~\cite{ref:silicon_data,ref:a16z_som}. Rising rentals are what scarcity looks like before the 2027--28 cohorts energize; a decline toward \$1 per H100-hour would be the floor arriving. The variable stays swept until that series moves.

\subsection{Utilization}
Inference throughput is modeled by operators at 40--50\% utilization of nominal capacity, the rest held as latency and peak buffer~\cite{ref:mehta_2026}. Batching, quantization and the proliferation of open-weight models could double useful utilization for the heaviest workloads (register entry \texttt{inference\_utilization}, Triang(0.40, 0.45, 0.50) today; doubling case as a sweep). The same stock then serves twice the tokens. But receipts in this paper are dollars, not tokens, and doubling tokens raises dollars only if the price per token holds --- which is precisely what the next subsection says it will not. Higher utilization is a Jevons effect: more inference, cheaper, from the same capital, with no relief for the carry.

\subsection{Batch share, and what four decades of computing already taught}
Utilization does not rise on its own; it rises when work is moved into batch. The labs now sell asynchronous batch processing at 50--75\% below interactive rates with 24-hour windows, and prompt caching at $\sim$90\% off repeated context~\cite{ref:batch_pricing}; the register carries a \emph{batch share} $b\sim$Triang(0.3, 0.5, 0.7) of 2030 tokens processed this way (a secondary variable). At the mode, effective fleet utilization rises from 45\% to $\sim$64\% and the volume-weighted effective price falls to $\sim$70\% of list; for heavy pipelines that stack batch and caching the effective price is 5--10 times below list, i.e.\ within one to three times the marginal cash cost. This is not new. Mainframe operators learned in the 1970s that a machine paid for by the shift is only economic if the night is filled with batch, and ran them at 90\%-plus; the grid- and distributed-computing decades made the same trade at the price of latency; and Bitcoin mining is the closest precedent of all --- a commodity computation on special-purpose silicon, where within two years of each ASIC generation the hashprice converged to the marginal cost of the cheapest power, the equipment was obsolete in 12--24 months, and capacity migrated to wherever electricity was cheapest. Inference has every one of those properties. The lesson of the precedents is that utilization gains are always taken by the buyer, not the producer: they lower the floor rather than lift the receipts. One caution on the numbers: Oracle reports GPU \emph{allocation} utilization of 97.9\%~\cite{ref:orcl_q1fy27} --- the share of its fleet rented out --- which is a different quantity from the 40--50\% \emph{throughput} utilization within a rented job that the batch argument concerns; both can be true at once.

\subsection{Who is the low-cost producer}
In a commodity the survivors at the floor are the producers with the lowest full cost --- the BYDs, not the incumbents --- and the paper's carry should be allocated to them by tier rather than spread across the industry. Figure~\ref{fig:coststack} builds a cost curve for five producer types from the same chain as the marginal cash cost (power, PUE, tariff~\cite{ref:eia_epm}, opex, tokens per hour) plus each tier's capital cost per token (capex per accelerator slot~\cite{ref:nvda_q1fy27}, life~\cite{ref:amzn_10k_life,ref:hudson_policy}, cost of capital~\cite{ref:damodaran2026}, utilization~\cite{ref:orcl_q1fy27}). At the modes: on-premise open-weight small models for routine work cost \$0.03 per million tokens all-in; in-house-silicon hyperscalers (TPU, Trainium, MTIA) \$0.14; Chinese western-corridor capacity on domestic silicon \$0.27 (cheap power offset by less efficient chips); Nvidia-based US hyperscalers \$0.29; debt-funded neo-clouds \$0.40. Against a blended list price near \$0.30 the two Nvidia-based tiers are at or above break-even on full cost; against the effective prices heavy users already pay (\$0.03--0.06), only the on-premise and in-house-silicon tiers cover their carry. That ordering is the exposure table of Section~\ref{sec:falsify} derived rather than asserted: the neo-cloud is the marginal producer the floor eliminates first; the Nvidia-based hyperscaler is next unless it moves to its own silicon; Google's TPU fleet is the industry's lowest-cost large producer; and the Chinese corridor competes on power while its chip efficiency lags. The ``BYD of AI'' is not one company but a tier --- in-house silicon on cheap power running open-weight models --- and the carry the paper computes falls, in the end, on whoever is not in it.

\begin{figure}[htbp]\centering\includegraphics[width=\textwidth]{f24_coststack.pdf}
\caption{The inference cost curve. Producers ordered by full cost per million tokens (bar width is share of 2030 capacity); dark is marginal cash cost, light adds the carry (depreciation and return). Dotted: blended list price; dashed: effective price paid by batch-and-cache heavy users. At effective prices only the two lowest tiers cover their carry.}\label{fig:coststack}\end{figure}

\subsection{The commodity floor}
In a commodity, price converges to marginal cash cost and the return on fixed capital goes to whoever built cheapest; uneconomic capital is written off or replaced by lighter structures~\cite{ref:mehta_2026}. Inference has the shape of a commodity: an undifferentiated unit (the token), many producers, falling costs, open-weight models removing the product moat. The marginal cash cost of a token is the cash it takes to produce one more --- power and facility opex, not depreciation or return. A chain: an accelerator drawing Triang(0.7, 0.9, 1.2) kW at PUE Triang(1.12, 1.22, 1.35) and a tariff of Triang(\$0.075, 0.10, 0.13)/kWh, plus non-power opex of Triang(\$0.05, 0.10, 0.18) per accelerator-hour, costs Triang-derived $\sim$\$0.20--0.35 per hour in cash; at Triang(3, 6, 10) million tokens per hour at 45\% utilization~\cite{ref:eia_epm,ref:batch_pricing} it produces tokens at a marginal cash cost of \$0.036 per million (P10--P90 \$0.026--0.053), or half that if utilization doubles (Figure~\ref{fig:econdep}, right). Realized \emph{list} prices are \$0.10--0.30 per million for small and mid models and \$1--5 for frontier models --- roughly ten times the floor, blended --- but effective prices for batch-and-cache pipelines are \$0.03--0.06, within one to three times the floor (register entry \texttt{token\_price\_index}, effective tier). If prices converge to the floor while volume rises to meet it, receipts converge on the fleet's cash cost --- the opex line, $\sim$\$130B in 2028 --- against a carry of \$540B then and \$725B in 2030. On the commodity view the depreciation and return terms are never recovered by anyone but the lowest-cost producer, and the paper's deficit is not a shortfall to be grown out of but the ordinary fate of fixed capital in a commodity whose price found its floor. The metric to monitor is therefore not revenue but the ratio of realized price per million tokens to marginal cash cost: the paper's register now carries both (\texttt{token\_price\_index}, \texttt{marginal\_cash\_cost}), and a ratio approaching one from above is the signal that the commodity view has won.

\begin{figure}[htbp]\centering\includegraphics[width=\textwidth]{f23_econdep_floor.pdf}
\caption{Left: the plateau carry under straight-line depreciation (base) and under declining-balance economic obsolescence at 50, 70 and 90\%/yr, against bottom-up receipts; economic depreciation front-loads the carry. Right: the marginal cash cost of inference per million tokens at current and doubled utilization, against realized price points; prices are about ten times the floor.}\label{fig:econdep}\end{figure}

\subsection{Market context}
The same reviewer notes that the pure-play GPU renters --- the neo-clouds of the exposure table --- are among the most heavily shorted positions in public markets in 2026, and that the uneconomic nature of the projects is argued publicly by short sellers including Jim Chanos~\cite{ref:mehta_2026}. This is corroboration of the paper's neo-cloud row, not evidence for it: a crowded short is a market opinion, and market opinions are inputs to nothing in the register. It is recorded because it locates the paper's argument in a live debate and because the short interest itself is one more observable to watch.

\section{The Bull Case at Its Strongest: a16z's \textit{State of Markets}, September 2026}\label{sec:a16z}
Andreessen Horowitz's September 2026 \textit{State of Markets} deck~\cite{ref:a16z_som} is the most complete statement of the case for the build-out, assembled by an investor with every reason to make it well. It is also the best test of this paper, because most of its charts bear directly on a register entry. Each of its claims is taken in turn: the evidence it offers, and what the register does with it.

\begin{enumerate}
\item \textbf{``We keep underestimating the capex.''} Street consensus for Big-5 capex has been revised upward every year; the September 2026 consensus (S\&P Capital IQ) is about \$1.05T for 2027 and \$1.15T for 2028, and J.P.~Morgan's figures are \$791B, \$1,066B and \$1,179B for 2026--28~\cite{ref:a16z_som,ref:jpmam_gtm}. \emph{What the register does with it:} the paper's ``arms race'' path (\$1.0--1.1T) is not a stress case --- it is the street's central expectation --- and the ``plateau'' is below consensus. Readers who take consensus at face value should read the arms-race rows of every table: a 2030 deficit of $\sim$\$545B and no year of profit through 2035. The paper labels the path accordingly from this version.

\item \textbf{``Capex is eating cash flow, for now.''} Hyperscaler free cash flow is near zero in 2026--27; the consensus has it rebounding past \$500B by 2030~\cite{ref:a16z_som,ref:jpmam_gtm}. \emph{Register:} a rebound of that size on a capex path that keeps rising to \$1.18T requires revenue above \$1T by 2029--30 --- precisely the top-down break-even test. The two consensus series (capex up, free cash flow up) are consistent with each other only on the paper's bull case; one of them will be revised.

\item \textbf{``Reports of GPU obsolescence are exaggerated.''} Silicon Data's September 2026 rental index puts the H100 at \$2.63 per GPU-hour, the A100 at \$1.59, the H200 at \$3.29 and the B200 at \$5.69, with residual values inferred from rentals rising through 2026~\cite{ref:a16z_som,ref:silicon_data}. \emph{Register:} this is the strongest evidence on the other side of the economic-depreciation variable $\delta$ (Section~\ref{sec:commodity}), and it is why $\delta$ is a secondary variable rather than a range. Two readings coexist. Rental prices rising while token prices fall means compute is scarce today; the paper's claim is about what happens when the 2027--28 cohorts energize. And \$2.63 per hour against a marginal cash cost of \$0.20--0.35 per hour is an 8--13$\times$ margin for the GPU owner --- the same ratio of price to floor the commodity-floor section finds at the token level, from the other end. The variable to watch is the H100 rental price through 2027; a decline toward \$1 is the floor arriving.

\item \textbf{``Backlogs have doubled year on year.''} Microsoft RPO, Google Cloud backlog and Amazon RPO together approach \$1.7T, up $\sim$150\% year on year~\cite{ref:a16z_som}. \emph{Register:} added as an observable. A backlog is a contract, not a payment: Oracle's \$664B converts over years and is over half one counterparty (Section~\ref{sec:observe}); Google's \$514B is not split between AI and the rest. Backlog growth is evidence of demand for capacity, which the paper does not dispute; the carry test is about the price that capacity ultimately earns.

\item \textbf{``Hyperscaler ROIC is above the cost of borrowing.''} Bank of America puts hyperscaler return on invested capital near 20\% against a cost of debt of 4--5\%~\cite{ref:a16z_som}. \emph{Register:} that ROIC is earned on the existing businesses --- search, advertising, the installed cloud --- not on the AI cohorts, whose return is what this paper computes. Borrowing at 4.5\% to fund assets whose incremental return is in question is the ordinary structure of a capital cycle, and the deck's own page on the railroads and telecoms (capex at 1.5--2\% of GDP at each cycle's peak) makes the comparison.

\item \textbf{``Data centers may be lowering your electricity bill.''} Watten, Bistline and Blanford find that each 10\% increase in state data-center capacity lowered average residential rates by about 0.4\% over 2019--24, by spreading the grid's fixed costs~\cite{ref:watten2026}. \emph{Register:} taken as stated, and it is an important correction to the ratepayer narrative. It does not change the siting variable, which is calibrated on what communities \emph{do} (Data Center Watch's blocked and delayed projects), not on what rates \emph{did}; and the study predates the 2025--26 load growth and the Northern Virginia and Texas tariff increases. Both facts stand: rates fell on average through 2024, and opposition has risen since.

\item \textbf{``Enterprise AI spending has plenty of room to run.''} Only 20\% of organizations cite cost as a constraint; cloud reached $\sim$4\% of IT budgets in its second year and AI is below that; and the deck argues the addressable market is ``considerably larger than IT budgets''~\cite{ref:a16z_som}. \emph{Register:} the last clause is the real disagreement, stated clearly. The paper's feasibility band is an IT-budget band. The deck's position is that AI is paid from labor budgets --- the \$30T global payroll --- not from the \$3.3T software-and-services pool. Whether labor budgets actually convert to AI spend is Section~\ref{sec:slack}'s question, and the paper's answer (about a quarter of nominal savings reach cash, over a J-curve) is why it keeps the IT-budget band. If the deck is right, the band is wrong by a factor of several and the plateau pays easily; if the slack trap is right, the band holds. This is the single most consequential disagreement between the two documents and the register cannot settle it from 2026 data.

\item \textbf{``Adoption is wide but shallow; power users are pulling away.''} The top 1\% of users consume orders of magnitude more AI than the field, and the gap is widening~\cite{ref:a16z_som}. \emph{Register:} the same observation as MIT NANDA's 95\% and the Census Bureau's tenth of firms (Section~\ref{sec:slack}); the deck reads it as upside, the paper as the early J-curve. Both are consistent with the data; they differ on timing, which is the residual adoption factor.

\item \textbf{``Tokens up 600$\times$ in two years; frontier labs capture 70\% of spend on 35\% of tokens.''} On OpenRouter, Anthropic, Google and OpenAI take $\sim$69\% of spending on $\sim$35\% of tokens; all other labs 65\% of tokens for 31\% of spending; open models on China-based hyperscaler endpoints are under 0.3\% of model spend~\cite{ref:a16z_som,ref:openrouter}. \emph{Register:} this is the cloud-capture share and the commodity floor measured from the demand side: two-thirds of tokens already trade at the cheap end, a frontier premium persists for frontier problems, and the premium tier is three labs. It raises confidence in $\theta_{\text{lab}}$ and in the cost-curve ordering, and it is the deck's best evidence for a durable frontier margin.

\item \textbf{``Jevons: compute prices rise as token prices fall.''} The deck's reading of rising H100 rentals alongside falling token prices is that efficiency gains are absorbed by demand~\cite{ref:a16z_som}. \emph{Register:} the paper agrees that this is what Jevons looks like while supply is short; the Bitcoin precedent (Section~\ref{sec:commodity}) is what it looks like after supply catches up. Which regime 2028 is in is the question the arms-race capex path answers one way and the adaptive path the other.

\item \textbf{``SaaS: prove-it, not pocalypse.''} Software revenue growth has slowed but stabilized; operating leverage is stable or improving; the sell-off has been selective (cyber and vertical over horizontal)~\cite{ref:a16z_som}. \emph{Register:} the seat-software growth variable $g_{\text{seat}}$ is a secondary variable for exactly this reason --- the fundamentals have not yet confirmed the repricing --- and Microsoft's +6--7\% seat growth is already in the incumbent case. The deck's own newspaper chart, in which stocks sold off five years before earnings collapsed, is the caution it offers against its own point.

\item \textbf{``Capex is at or near the peaks of cycles past.''} Lease-adjusted hyperscaler capex at $\sim$1.5\% of US GDP, against $\sim$2\% for the railroads and $\sim$1.5\% for telecoms at their peaks~\cite{ref:a16z_som}. \emph{Register:} the paper's own scale statement (2.5\% of GDP on a gross basis) is brought into line as 1.5--2.5\% depending on definition; the comparison to the two prior cycles that ended in impairment is the deck's, and the paper adopts it.
\end{enumerate}

Taken together the deck and the paper agree on nearly every observable and disagree on two judgments: whether AI will be paid from labor budgets as well as IT budgets, and whether rising compute prices are the permanent state of a supply-constrained industry or the top of a cycle. The register names both as the places to look.

\section{The Sell-Side Hurdle: Goldman Sachs's Phase-2 ROIC Framework, September 2026}\label{sec:gs}
On September 24, 2026, Goldman Sachs's Americas Technology team (Sheridan, Borges, Schneider and colleagues) published a 29-page note, \textit{Sizing the AI Economy Needed to Justify ROIC on Hyperscaler AI Capex}~\cite{ref:gs_roic2026}, whose question is this paper's question asked from the other side of the table: not ``what is the carry'' but ``what revenue would make the carry worth paying.'' The note is taken here the way Section~\ref{sec:a16z} took a16z's deck --- claim by claim, against the register. It is the closest thing to a sell-side replication of the supply-side model that exists, and the two agree on the hurdle to within a few percent; where they differ is scope, one input, and what is assumed about demand.

\subsection{What the note computes}
The framework isolates one cohort: ``Phase 2,'' the AI-compute capex of six US hyperscalers (Alphabet, Microsoft, Amazon, Meta, Oracle and SpaceX) over calendar 2026--27, which the note puts at \$1.726T, taking 85--90\% of each company's total capex as AI-related (AWS segment capex for Amazon, the AI segment for SpaceX). That cohort is split 70\% ``compute'' (servers, chips, networking) depreciated straight-line over 5 years and 30\% ``shells'' (land, buildings, infrastructure) over 15 years. An all-in cost of \$42.4B per gigawatt (Epoch AI's year-end 2025 inputs inflated 10\%; servers \$24.1B, network \$12.6B, facilities \$5.4B) implies $\sim$41 GW added. Operating cost excluding depreciation is \$836M per GW-year (also from Epoch, plus 10\%), scaled by an operational ramp of 60\%, 80\% and 85\% in 2028, 2029 and 2030. The return threshold is a 15\% return on invested capital, defined as net operating profit after tax (NOPAT --- operating profit less a 21\% tax) in each of 2028--30 divided by the \emph{average annual} capex of the phase (\$863B). Grossing NOPAT up for tax and adding depreciation and operating cost gives the revenue the cohort must earn: \$465B, \$474B and \$476B in 2028, 2029 and 2030 --- \textbf{\$1.42T cumulative, or \$11.6B per GW-year} --- at an operating margin of $\sim$35\%. A sensitivity table runs the ROIC threshold from 0\% to 30\% and the cost per GW from \$34B to \$51B: the three-year revenue requirement spans \$908B to \$1.89T, i.e.\ \$6.2--18.6B per GW-year.

The note then benchmarks the requirement against contracted backlog rather than against a demand forecast. For the three public cloud providers alone (Amazon, Microsoft, Alphabet; Phase-2 capex \$1.22T, $\sim$29 GW) the same framework requires \$1.00T of 2028--30 revenue, which is 59\% of their combined \$1.69T backlog at Q2 2026 (up 152\% year on year): 78\% of Alphabet's \$496B, 68\% of Amazon's \$514B and 40\% of Microsoft's \$678B. The note's conclusion is that near-term return compression is ``a natural consequence of the massive upfront investment cycle rather than evidence of structurally poor AI economics,'' and that the 15\% threshold is conservative against what the companies themselves describe: Amazon's $\sim$3-year breakeven on servers with a 5--6 year life, Oracle's ``high-20s'' steady-state ROIC on large projects, and SpaceX's stated one-year payback on terrestrial AI capacity. A separate semiconductor-side build in the same note --- Nvidia, Broadcom, AMD, Marvell and MediaTek revenue per GW --- implies total AI infrastructure capex of \$861B, \$1.34T and \$2.0T in 2026, 2027 and 2028 (20, 35 and 57 GW), a wider set than the hyperscalers because it includes labs, sovereigns and enterprises.

\subsection{The same hurdle in the paper's units}
The note's requirement can be restated per dollar of cohort capital and set beside the carry of Section~\ref{sec:supply}, which is also a cost per dollar of in-service gross. Table~\ref{tab:gs_reconcile} does this for the note's 2029 column and for the paper's 2028 plateau at the modes (Appendix~\ref{app:equations}, worked example: gross \$1.9T, carry \$540B).

\begin{table}[htbp]
\centering\small
\caption{The Goldman Sachs Phase-2 requirement and the paper's carry, per dollar of capital per year}
\label{tab:gs_reconcile}
\begin{tabularx}{\textwidth}{L{3.0cm}L{1.6cm}L{1.6cm}X}
\toprule
\textbf{Component} & \textbf{GS, 2029 (\$/\$ capex)} & \textbf{Paper, 2028 (\$/\$ gross)} & \textbf{Why they differ} \\
\midrule
Depreciation & 0.160 & 0.153 & GS: 70\% compute at 5 yr $+$ 30\% shell at 15 yr. Paper: 55\% silicon at 4.2 yr (harmonic across fleets) $+$ 45\% shell at 20 yr. The shorter life and the smaller silicon share offset. \\
Operating cost (excl.\ depreciation) & 0.020 & 0.067 & GS: Epoch's \$836M/GW-yr, essentially power and maintenance, at 60--85\% ramp. Paper: power 1.9\% $+$ staffing/maintenance 2.0\% $+$ property tax and insurance 1.0\% $+$ network 1.0\% $+$ other 0.6\% (Section~\ref{sec:decomp}). The one material gap; see below. \\
Return & 0.095 & 0.064 & GS: 15\% NOPAT on average annual capex (half the cohort), grossed up to pre-tax EBIT. Paper: CAPM cost of capital (9.45\%) on undepreciated net book ($\sim$68\% of gross). GS's return is a target; the paper's is the floor below which capital is destroyed. \\
\midrule
\textbf{Required revenue} & \textbf{0.274} & \textbf{0.284} & Within 4\%. On the same \$1.9T stock the GS framework would ask \$520B/yr; the paper's carry is \$540B. \\
\bottomrule
\end{tabularx}
\end{table}

Three things follow. First, the hurdle is not in dispute. A sell-side desk with every incentive to set it low and this paper with its ranges decomposed arrive at 27--28 cents of annual revenue per dollar of capital; the note's own sensitivity (\$6.2--18.6B per GW-year) brackets the paper's modes. Readers who distrust the paper's supply side can substitute Goldman's and change nothing in the result. Second, the note's \$1.42T is the requirement of \emph{one} cohort and ignores, by construction, both the 2023--25 cohorts still on the books and the capex of 2028 and beyond; its authors say so. The paper's carry stacks the cohorts (that is what a vintage model is for), which is why its 2028 figure applies to \$1.9T of gross and its 2030 figure to $\sim$\$2.5T. Applying the note's own per-dollar requirement to the stacked stock gives the paper's result. Third, the cohort itself is larger in the note (\$1.73T) than in the paper (\$1.23T AI-attributable for 2026--27) for two scope reasons: the note counts 85--90\% of all capex as AI, where the paper nets a non-AI baseline that filings pin at $\sim$\$130B in 2023 growing 2--10\%/yr (Section~\ref{sec:decomp}); and the note adds SpaceX. Both are defensible; the register records the note's convention as the upper bound on the AI-attributable share.

\subsection{Where the paper is harsher, and where it is not}
\textbf{Operating cost.} The one input on which the two frameworks differ materially is non-depreciation operating cost: 2\% of capital per year in the note (power and maintenance, from Epoch) against 6.65\% in the paper, whose decomposition adds staffing, property tax and insurance, network and ``other'' at 4.6\% of gross. On the 2028 plateau stock that is $\sim$\$85B/yr of carry; on the 2030 stock $\sim$\$115B. If the note's figure is right and the paper's is a double count --- property tax is often abated for data centers, and network is partly in the capex --- the 2030 plateau deficit falls from $\sim$\$255B to $\sim$\$140B and its z-score from $\sim$3.5 to $\sim$2; the sign and the 2027--28 sunk deficit are unchanged. The register now carries both anchors on $c_{\text{opex}}$ (Appendix~\ref{app:inputs}) and the input is flagged for a sourced operator cost per MW-year; by the update protocol's rule (Appendix~\ref{app:update}) adopting the lower anchor would be a material re-issue, so it is not done silently here.

\textbf{Demand.} The note does not forecast the revenue; it shows what revenue is needed and observes that backlog covers it. The paper's two demand estimates put total infrastructure receipts --- from \emph{all} cohorts --- at $\sim$\$370B in 2028 and $\sim$\$480B in 2030 (Table~\ref{tab:bu_results}), i.e.\ 2030 receipts from the whole stock roughly equal to what the note requires from the 2026--27 cohort alone. The backlog benchmark is addressed in Section~\ref{sec:a16z} (item 4) and Section~\ref{sec:observe}: a backlog is a contract, not a payment; the three CSP backlogs include non-AI commitments; Microsoft's includes OpenAI's multi-year Azure commitment and Oracle's is over half OpenAI; and the note itself assumes no further backlog growth. Converting 59\% of today's backlog over 2028--30 requires the counterparties to pay, which is the revenue-quality question of Section~\ref{sec:demand}, not a different answer to it.

\textbf{Ramp and utilization.} The note's 60/80/85\% is the share of built capacity that is operational, a cohort ramp; the paper's energization lag and grid cap play the same role, and its inference utilization (40--50\% of nominal, Section~\ref{sec:commodity}) is a different quantity --- how hard operational capacity is run. The two are multiplicative, and the note's revenue per GW is stated at 100\% utilization of operational capacity.

\textbf{Gigawatts.} The note's semiconductor build implies 57 GW of AI deployments in 2028 worldwide. The paper's grid energization cap is 22--40 GW/yr. One of the two is wrong, or 2028's silicon will be bought before it can be powered; the register records the note's figure as the high anchor on the cap and the 2028 deployment number as an observable to test in the Q4 2027 and 2028 utility filings.

\textbf{Fungibility.} The note reports the operators' framing: Microsoft's ``fungible fleet,'' Amazon's modular retrofits, Alphabet shifting capacity between internal and external workloads. The paper does not dispute fungibility \emph{within} AI workloads --- training to inference, between geographies --- which is exactly what those statements claim. The dual-use fallacy of Section~\ref{sec:supply} is about AI halls hosting non-AI cloud, which none of the statements claim.

\textbf{What the register takes.} Five observables (Table~\ref{tab:observe}): the Phase-2 requirement of \$1.42T and \$11.6B/GW-yr at 15\% ROIC; the \$42.4B/GW cost and its \$24.1/12.6/5.4B split; Epoch's \$836M/GW-yr operating cost; the 59\%-of-backlog conversion the three CSPs would need; and the 20/35/57 GW deployment path. One flagged input ($c_{\text{opex}}$, two anchors) and one high anchor (the grid cap). The note's management claims --- Amazon's 3-year server breakeven, Oracle's high-20s ROIC, SpaceX's one-year payback --- are recorded as grade-J statements by interested parties, to be tested when segment disclosures permit. The note and the paper disagree on nothing arithmetical; they disagree on whether the cohorts after this one, and the ones before it, belong in the sum. The paper's answer is that a balance sheet does not get to choose.

\subsection{Four arguments from the first review round, not yet in the model}\label{sec:firstround}
The first circulated draft was built in a long exchange whose working notes contained several mechanisms that did not survive into the paper's text, mostly because they lacked a source at the time. They are listed here so they are not lost, each with what it would change and what would put it into the register.
\begin{enumerate}
\item \textbf{Take-or-pay power.} The paper's facility opex scales with load. A growing share of hyperscale power is contracted under long-dated take-or-pay power purchase agreements --- Microsoft's 20-year contract for the restarted Three Mile Island unit, Amazon's Susquehanna arrangement, Meta's 20-year Clinton contract~\cite{ref:ppa_press}\verify{} --- under which the payment is owed whether or not the capacity runs. For the contracted share, power moves from opex into carry: a fixed liability that does not fall when utilization does. At 2\% of gross for power, a 30--50\% contracted share converts \$10--25B/yr of 2028 opex into a fixed charge. \emph{Register:} add a contracted-power share when a sourced figure exists; the effect is on the \emph{shape} of the downside (the carry no longer flexes down in a demand shortfall), not on the base case.
\item \textbf{The sovereign backstop.} If the 2030 plateau deficit is $\sim$\$255B against a carry of $\sim$\$725B, a state that wished to keep the stock running --- for defense, sovereignty or employment reasons --- would need to absorb roughly 35\% of the carry. That is the scale of the implicit guarantee behind ``AI is strategic,'' and it is not small: comparable to a year of US federal R\&D. \emph{Register:} a what-if, not a variable; it is the size of the public-money question if the private one is answered the way this paper expects.
\item \textbf{Agentic loops as the Jevons mechanism.} The 600$\times$ growth in tokens on OpenRouter (Section~\ref{sec:a16z}) is mostly agents calling models many times per task. Tokens per task rising 10--100$\times$ as work moves from chat to agents is the mechanism by which Jevons could keep hyperscaler receipts growing after the per-token price collapses to the floor. \emph{Register:} already inside the residual adoption/capability factor; the observable is tokens per paid seat, which no operator yet discloses.
\item \textbf{Outcome-based procurement.} a16z's labor-budget argument (Section~\ref{sec:a16z}, item 7) has arithmetic: 2\% of the \$30T global knowledge payroll is \$600B/yr, three times the paper's 2030 IT-budget feasibility band. The paper's slack and leakage sections say why that 2\% is hard to extract and why, when it is, 10--50:1 of it stays with buyers and integrators. \emph{Register:} the feasibility band's upper edge is where this argument lives; it is already the axis of Figure~\ref{fig:probmap}.
\end{enumerate}
Two further notes from the same exchange --- a three-bloc comparison of five-year cost per megawatt (roughly \$42--50M in Europe, \$32--38M in the US, \$14--18M in China's western corridors) and a sketch of rest-of-world demand (India, ASEAN, Africa, Latin America) --- were derived without sources and are recorded here only as questions for the register\verify.

\section{Results: Break-Even Growth}\label{sec:results}
Each run draws $N=100{,}000$ samples of the parameters in Sections~\ref{sec:supply}--\ref{sec:demand} and computes the annual surplus or deficit $\Delta_t=D_t-H_t$ (positive is a surplus) for 2027--2030. Parameters are independent; since capex, energization and demand expectations are positively correlated in reality, the spread is understated and tail probabilities are indicative.

\begin{figure}[htbp]\centering\includegraphics[width=0.9\textwidth]{f3_scissors.pdf}
\caption{Carry (Plateau) against infrastructure receipts at four sustained growth rates. The crossing year is the break-even year for that $g$.}\label{fig:scissors}\end{figure}

\begin{figure}[htbp]\centering\includegraphics[width=0.85\textwidth]{f4_delta.pdf}
\caption{Median annual surplus or deficit $\Delta_t$ under Plateau for six growth rates. No rate clears 2027.}\label{fig:delta}\end{figure}

\begin{table}[htbp]
\centering\small
\caption{Plateau path: median $\Delta_t$ (\$B) and P($\Delta_t<0$) by sustained growth $g$}
\label{tab:delta}
\begin{tabular}{lcccc}
\toprule
$g$ \textbf{(\%/yr)} & \textbf{2027} & \textbf{2028} & \textbf{2029} & \textbf{2030} \\
\midrule
20 & $-215$ (100\%) & $-325$ (100\%) & $-400$ (100\%) & $-425$ (100\%) \\
30 & $-200$ (100\%) & $-290$ (100\%) & $-330$ (100\%) & $-310$ (100\%) \\
40 & $-185$ (100\%) & $-250$ (100\%) & $-250$ (100\%) & $-165$ (100\%) \\
50 & $-170$ (100\%) & $-210$ (100\%) & $-160$ (100\%) & $+10$ (41\%) \\
60 & $-155$ (100\%) & $-165$ (100\%) & $-55$ (92\%) & $+225$ (0\%) \\
80 & $-130$ (100\%) & $-65$ (98\%) & $+200$ (0\%) & $+800$ (0\%) \\
\bottomrule
\end{tabular}
\end{table}

\begin{figure}[htbp]\centering\includegraphics[width=0.85\textwidth]{f5_breakeven.pdf}
\caption{Break-even sustained growth $g^{*}$ by year and capex path, with the 2026 revenue base at its low and high ends.}\label{fig:breakeven}\end{figure}

\begin{table}[htbp]
\centering\small
\caption{Break-even sustained growth $g^{*}$ (median $\Delta_t=0$)}
\label{tab:breakeven}
\begin{tabular}{lcccc}
\toprule
\textbf{Path / variant} & \textbf{2027} & \textbf{2028} & \textbf{2029} & \textbf{2030} \\
\midrule
Plateau & $>$150\% & 92\% & 65\% & 49\% \\
Arms race & $>$150\% & 111\% & 80\% & 62\% \\
Rationalize & 142\% & 72\% & 47\% & 34\% \\
\midrule
Plateau, $D_{2026}=\$163$B (P10) & & 101\% & 70\% & 53\% \\
Plateau, $D_{2026}=\$194$B (P90) & & 84\% & 60\% & 46\% \\
Plateau, accounting lives (4.5--6.5 yr) & & 87\% & 64\% & 51\% \\
Plateau, 8-year silicon life & & 79\% & 60\% & 49\% \\
Plateau, thermal life Triang(2.5, 3.5, 5) & & 93\% & 64\% & 49\% \\
Plateau, sovereignty overlay ($\kappa=1.2$, $\gamma=0.25$) & & & & $\sim$66\% \\
\bottomrule
\end{tabular}
\end{table}

\begin{figure}[htbp]\centering\includegraphics[width=0.85\textwidth]{f6_cum.pdf}
\caption{Median cumulative $\Delta$ over 2027--30 versus $g$ (band is P10--P90 for Plateau).}\label{fig:cum}\end{figure}

\begin{table}[htbp]
\centering\small
\caption{Cumulative $\Delta$ 2027--30 (\$B, median; P10--P90; P(cumulative $<0$))}
\label{tab:cum}
\begin{tabular}{lccc}
\toprule
\textbf{Path} & $g=30\%$ & $g=50\%$ & $g=70\%$ \\
\midrule
Plateau & $-1{,}135$ ($-1{,}275$, $-995$; 100\%) & $-530$ ($-690$, $-370$; 100\%) & $+295$ ($+105$, $+485$; 2\%) \\
Arms race & $-1{,}770$ (100\%) & $-1{,}165$ ($-1{,}355$, $-980$; 100\%) & $-340$ ($-560$, $-130$; 98\%) \\
Rationalize & $-540$ ($-660$, $-420$; 100\%) & $+65$ ($-80$, $+205$; 28\%) & $+890$ ($+715$, $+1{,}065$; 0\%) \\
\bottomrule
\end{tabular}
\end{table}

\begin{figure}[htbp]\centering\includegraphics[width=0.85\textwidth]{f7_tornado.pdf}
\caption{One-at-a-time sensitivity of $\Delta_{2030}$ (Plateau, $g=40\%$) to each parameter at its minimum and maximum, others at mode.}\label{fig:tornado}\end{figure}

\begin{figure}[htbp]\centering\includegraphics[width=0.8\textwidth]{f8_hist.pdf}
\caption{Distribution of $\Delta_{2030}$ under Plateau at $g=40\%$, the neighborhood of break-even.}\label{fig:hist}\end{figure}


\subsection{Bottom-up result}
Table~\ref{tab:bu_results} and Figure~\ref{fig:reconcile} compare the carry with infrastructure receipts from the fifteen sectors, net of $\sim$\$12B/yr of genuinely circular flows $F_{\text{circ}}$ (Section~\ref{sec:circ}); $\kappa$ in the table is the sovereignty capital multiplier, 1 in the base case.

\begin{table}[htbp]
\centering\small
\caption{Bottom-up: median $\Delta_t=\text{Receipts}_t-F_{\text{circ}}-\kappa H_t$ (\$B) and P($\Delta_t<0$); cumulative 2027--30}
\label{tab:bu_results}
\begin{tabular}{llccccc}
\toprule
\textbf{Path} & \textbf{Sovereignty} & \textbf{2027} & \textbf{2028} & \textbf{2029} & \textbf{2030} & \textbf{Cum.\ (P$<$0)} \\
\midrule
Plateau & none & $-155$ (100\%) & $-205$ (100\%) & $-240$ (100\%) & $-260$ (100\%) & $-855$ (100\%) \\
Plateau & $\kappa$=1.2, $\gamma$=0.25 & $-250$ & $-340$ & $-405$ & $-450$ & $-1{,}450$ (100\%) \\
Arms race & none & $-190$ & $-315$ & $-440$ & $-545$ & $-1{,}490$ (100\%) \\
Arms race & $\kappa$=1.2, $\gamma$=0.25 & $-295$ & $-475$ & $-645$ & $-795$ & $-2{,}210$ (100\%) \\
Rationalize & none & $-115$ (100\%) & $-95$ (98\%) & $-55$ (80\%) & $+5$ (48\%) & $-260$ (89\%) \\
Rationalize & $\kappa$=1.2, $\gamma$=0.25 & $-205$ & $-210$ & $-185$ & $-135$ & $-735$ (100\%) \\
\bottomrule
\end{tabular}
\end{table}

\begin{figure}[htbp]\centering\includegraphics[width=0.9\textwidth]{f10_reconcile.pdf}
\caption{Reconciliation of the two demand estimates against the plateau carry. The top-down path at $g^{*}=49\%$ is what the stock needs; the bottom-up receipts are what the sectors deliver.}\label{fig:reconcile}\end{figure}

\begin{figure}[htbp]\centering\includegraphics[width=0.85\textwidth]{f12_growth.pdf}
\caption{Growth delivered by the sectors (year on year) versus sustained growth required by each capex path. The curves cross in the wrong direction: the sectors decelerate as the requirement falls, and never meet it on a plateau.}\label{fig:growth}\end{figure}

\begin{figure}[htbp]\centering\includegraphics[width=0.85\textwidth]{f11_sector_tornado.pdf}
\caption{Sensitivity of the 2030 bottom-up deficit to each sector's ceiling at its minimum and maximum. No single sector ceiling closes the gap; advertising and assistants are the ones that could move it most.}\label{fig:sector_tornado}\end{figure}

\subsection{Variance budget, undecomposed and decomposed}
Equation~\ref{eq:var} in practice: Table~\ref{tab:varbudget} apportions the one-at-a-time swings into shares of total variance (half-swing squared over the sum), with the inputs treated first as single judgment ranges and then as the decompositions of Section~\ref{sec:decomp}. Undecomposed, two inputs --- the 2026 base and the cloud share --- would carry 82\% of the top-down uncertainty; decomposed, the root-sum-square is 104. What is left is the fleet-mix silicon life on the supply side and, on the demand side, the base and the cloud share, both already reduced to public components; their remaining spread is the labs' second-half 2026 bookings and Anthropic's undisclosed margin, which the next two quarters of filings will resolve. On the bottom-up side, decomposing the advertising ceiling takes its contribution from half the sector variance to a tenth, leaving assistants and the common factor. The root-sum-square is less than half the linear sum throughout --- the quadrature discount --- and the tight factors contribute almost nothing even though there are many of them.

\begin{table}[htbp]
\centering\small
\caption{Variance budget of $\Delta_{2030}$ (Plateau). Half-swing is half the min-to-max change in median $\Delta_{2030}$; share is its squared contribution to the total}
\label{tab:varbudget}
\begin{tabular}{lrrrr}
\toprule
 & \multicolumn{2}{c}{\textbf{Single judgment ranges}} & \multicolumn{2}{c}{\textbf{Decomposed (Section~\ref{sec:decomp})}} \\
\textbf{Top-down input ($\Delta_{2030}$, Plateau, $g=g^{*}$)} & \textbf{Half-swing} & \textbf{Share} & \textbf{Half-swing} & \textbf{Share} \\
\midrule
2026 revenue base (7 components, hyperscaler part decomposed) & 229 & 64\% & 64 & 38\% \\
Cloud share $\theta$ (by channel and by lab) & 121 & 18\% & 52 & 25\% \\
Facility opex (5 components, plus water uplift) & 78 & 7\% & 20 & 4\% \\
Silicon life (training/inference fleet mix) & 66 & 5\% & 48 & 22\% \\
Non-AI baseline (constrained by 2023) & 42 & 2\% & $<$10 & $<$1\% \\
Energization lag (silicon/shell) & 41 & 2\% & $<$10 & $<$1\% \\
Cost of capital (CAPM) & 26 & 1\% & 29 & 8\% \\
Shell life; silicon share (Nvidia-anchored) & 13, 12 & $<$1\% & 12, 17 & $\sim$4\% \\
Grid energization cap; siting deferral; site cost and cooling uplifts (new) & --- & --- & $<$12 each & $\sim$2\% \\
\midrule
Root-sum-square & 286 & & 104 & \\
Linear sum & 628 & & 280 & \\
\midrule
\textbf{Bottom-up sector ceiling} & & & & \\
Advertising uplift & 54 & 51\% & 18 & 11\% \\
Office assistants / agents & 36 & 23\% & 36 & 45\% \\
Consumer (now decomposed) & 24 & 10\% & 12 & 6\% \\
Search \& internal; coding & 18, 18 & 11\% & 18, 18 & 26\% \\
Remaining ten sectors & $<$9 each & 5\% & $<$9 each & 6\% \\
Root-sum-square & 76 & & 50 & \\
\bottomrule
\end{tabular}
\end{table}

\subsection{Reconciling the two estimates}
The top-down sweep says the plateau stock needs 49\%/yr sustained through 2030 (about \$880B of end-customer revenue; the siting friction defers carry but adds cost, and the two nearly cancel by 2030). The bottom-up chains deliver 33\% CAGR (about \$625B), with year-on-year growth falling from 57\% to 14\% as the large sectors saturate. The two methods share only the 2026 base; the ceilings, ramps and capture shares are independent judgments, and the gap between them --- roughly \$255B of end-customer revenue or \$260B of infrastructure receipts in 2030 --- is about five times the root-sum-square half-swing of the bottom-up estimate (Table~\ref{tab:varbudget}); expressed as a $z$-score of the difference, $\Delta_{2030}/\sigma_{\Delta}=-259/73\approx-3.5$ (Appendix~\ref{app:errors}). That is the Fermi standard for a robust sign: not certainty, but a gap larger than the propagated noise. Figure~\ref{fig:sector_tornado} shows that no single sector ceiling closes it; advertising uplift at its decomposed maximum (\$170B) no longer can even with assistants at theirs (\$320B); Figure~\ref{fig:whatif2} maps the pair: no combination of the two ceilings within \$300B and \$600B closes the 2030 gap on its own, because ramps and capture shares bind as well. The concrete bull case is therefore not ``bigger ceilings'' but ``faster ramps and higher cloud capture at once'' --- the residual factor of Figure~\ref{fig:whatif1}.

\subsection{Can both be true? Cannibalization and the ``internally inconsistent'' objection}
Bank of America's note on the 2026 software repricing put the objection plainly: investors cannot believe both that AI capital expenditure fails to earn its return and that AI disrupts the whole software industry~\cite{ref:saaspocalypse_nx}. The paper's answer is that the two are one claim, and the mechanism is cannibalization. AI does not add a new spending pool to the enterprise budget; it enters an existing one. Three cases span the possibilities. If AI spend is \emph{additive}, total software-plus-AI spending equals the old software pool plus the AI pool, and the incumbents keep theirs. If it is \emph{substitutive}, every AI dollar replaces a software dollar and the total is unchanged. The evidence points between: clients do pay more for software with AI in it --- Microsoft's July price increases are the proof --- but they buy fewer seats, and the seats they drop were the incumbents' highest-margin revenue. The total rises, but by \emph{less than the sum}: the feasibility band of Section~\ref{sec:decomp} is built exactly this way, with a reallocation share (12--32\% of the pool moving from seats and services to AI solutions) plus a small genuinely new component (2--7\%). So AI revenue can reach $\sim$\$1T by 2030 while the software industry shrinks, because most of that trillion is the software industry's own money changing hands --- from seat-priced incumbents to labs, integrators and on-device silicon --- and only a minority of it reaches the infrastructure meter. Both halves of the objection are consequences of the same flow: the buyer's dollar moves, the incumbent loses it, the infrastructure owner sees a fraction. There is no inconsistency; there is a leak, and the leak is the paper.

\subsection{The overlap gauge}\label{sec:overlap}
The question a Fermi analysis should end on is not ``what is the sign'' but ``with what probability can the two distributions overlap, and inside what feasible region.'' Figure~\ref{fig:overlap} overlays the bottom-up receipts and carry distributions year by year; Table~\ref{tab:pprofit} states P(receipts $>$ carry) for every path and year; and Figure~\ref{fig:probmap} maps that probability over the two residual axes with their feasibility cutoffs drawn in.

\begin{figure}[htbp]\centering\includegraphics[width=\textwidth]{f16_overlap.pdf}
\caption{Overlap of the bottom-up receipts and carry distributions by year (Plateau). The distributions do not touch in any year.}\label{fig:overlap}\end{figure}

\begin{table}[htbp]
\centering\small
\caption{P(profit) $=$ P(receipts $-$ circular $>$ carry), bottom-up demand, by path and year}
\label{tab:pprofit}
\begin{tabular}{llcccc}
\toprule
\textbf{Path} & \textbf{Sovereignty} & \textbf{2027} & \textbf{2028} & \textbf{2029} & \textbf{2030} \\
\midrule
Plateau & none / overlay & 0\% / 0\% & 0\% / 0\% & 0\% / 0\% & 0\% / 0\% \\
Arms race & none / overlay & 0\% / 0\% & 0\% / 0\% & 0\% / 0\% & 0\% / 0\% \\
Rationalize & none & 0\% & 2\% & 20\% & 52\% \\
Rationalize & overlay & 0\% & 0\% & 0\% & 3\% \\
\bottomrule
\end{tabular}
\end{table}

\begin{figure}[htbp]\centering\includegraphics[width=\textwidth]{f17_probmap.pdf}
\caption{Probability of profit in 2030 over the two residual axes, by capex path. Solid contour: 50\%. Dashed line: the IT-budget hard cap on the assistant ceiling ($\sim$\$270B). Dotted curves bound the feasibility band: total 2030 AI revenue between \$920B and \$1.31T, the P10--P90 of the decomposed budget cutoff (Section~\ref{sec:decomp}); to the right of the band the demand cannot be financed from enterprise, advertising and consumer budgets. The dot is the paper's modes.}\label{fig:probmap}\end{figure}

The map is the paper's result in one picture, and the feasibility band is what keeps it honest. A single line at 8\% of total IT spend ($\sim$\$810B) would put the plateau's 50\% contour outside the feasible region; decomposing the cutoff into Gartner's software and services budgets, their growth, and an AI share (Section~\ref{sec:decomp}) gives a band of \$920B--\$1.31T --- almost unchanged in total even though the pool it comes from is being torn up, because the money moves from seats to solutions rather than disappearing (Section~\ref{sec:decomp}) --- and the plateau's 50\% contour --- residual factor $\sim$1.5, 2030 AI revenue $\sim$\$950B --- lies in its upper half. So the honest statement is: on a plateau, the installed base is more likely than not to earn its 2030 carry only if the sector chains are $\sim$50\% too conservative in aggregate, which is financeable in roughly 85\% of budget draws but requires AI to reach $\sim$25\% of an enterprise software and services pool that is no longer growing, up from $\sim$5\% in 2026 (Gartner's \$64B of AI platform and model spend against \$3.3T~\cite{ref:gartner_ai_platforms,ref:gartner_it_jul2026}) --- money that comes out of seat-priced software and consulting hours. Under an arms race the contour is off the chart. Under rationalization the feasible and profitable regions overlap: P(profit) is $\sim$52\% at the paper's modes in 2030; siting friction pulls it down because deferred cohorts and site costs land in 2029--30. The zeros in Table~\ref{tab:pprofit} for 2027--28 are not a modeling artifact: the 2025--27 cohorts are on the books and demand cannot double fast enough to meet them under any path.

Three questions the feasibility band raises. The first is the deepest: a16z's \textit{State of Markets} argues that AI's addressable market is ``considerably larger than IT budgets'' because AI is paid from labor budgets~\cite{ref:a16z_som} --- the \$30T payroll, not the \$3.3T software-and-services pool. If so, the band is wrong by a factor of several and the plateau pays easily. The paper keeps the IT-budget band because Section~\ref{sec:slack} finds that only about a quarter of nominal labor savings reach cash, and over a J-curve; but this is the disagreement on which the bull case turns, and the register cannot settle it from 2026 data. Two further questions the band raises. First, is AI spending \emph{additional} to IT budgets or does it \emph{displace} other software and services? Gartner's 2026 forecast has total IT spending growing 14\%, most of it AI-driven, so the near-term answer is mostly additional; but the band's AI share is a share of the whole software-and-services pool, so if AI displaces seat-based software and consulting hours one-for-one the cap on AI is the same and the hyperscalers' non-AI cloud revenue shrinks instead --- Microsoft can gain AI revenue while losing on conventional Office. Second, AI's own efficiency will shrink the services budget it draws on: if AI removes \$200--400B of the \$1.87T services pool by 2030, the band's upper edge falls by roughly that share of the AI cut. Both effects narrow the feasible region; neither is in the base case, which is therefore generous to the build-out on this point.

 No plausible $g$ clears 2027 or 2028: the 2025--27 cohorts are a sunk deficit of \$100--250B/yr. Under Plateau, $g\approx49\%$ sustained for four years breaks even in 2030 alone and $\sim$64\% recovers the 2027--28 hole cumulatively. The tornado (Figure~\ref{fig:tornado}) says what a Fermi analyst wants to know: the 2026 revenue base and the infrastructure share dominate everything on the supply side; silicon life, opex and the non-AI baseline are second-order; cost of capital and shell life barely register. The distribution at $g=49\%$ (Figure~\ref{fig:hist}) is centered near zero --- the coin-toss neighborhood by construction. The bottom-up estimate says the sectors do not reach that neighborhood.

\noindent\textbf{Where the data stand.} Frontier-lab run-rates rose roughly 3.5$\times$ (Anthropic $\sim$7$\times$, OpenAI $\sim$2$\times$) between end-2025 and August 2026~\cite{ref:techcrunch_anthropic,ref:axios_anthropic,ref:epoch_revenue}. Both companies have told investors to expect slower growth in 2026 than trend --- roughly 2.2$\times$ for OpenAI and up to 4$\times$ for Anthropic~\cite{ref:epoch_revenue_x}. Even those guided rates exceed $g^{*}$. The labs' figures are the only high-growth disclosures available; Microsoft's consumer assistant has by most accounts disappointed, and Amazon's and Meta's AI revenue is not separately reported, so ``ahead of the threshold'' rests on two companies. The paper's position is therefore: the build-out pays if the labs' growth halves and then halves again but stays above $\sim$50\% through 2030. The fifteen sector chains put it nearer 14\% by 2030 because seats, IT budgets and advertising yields run out before compute does; Sections~\ref{sec:slack}--\ref{sec:lastmile} are the mechanisms behind those ceilings; and Section~\ref{sec:falsify} says what would move the distribution.

\section{Winners and Losers}
\begin{table}[htbp]
\centering
\small
\caption{Exposure hierarchy}
\label{tab:losers}
\begin{tabularx}{\textwidth}{L{2.6cm}L{3.0cm}Xl}
\toprule
\textbf{Tier} & \textbf{Entities} & \textbf{Mechanism} & \textbf{Severity} \\
\midrule
Existential & Oracle Cloud Infrastructure & \$300B+ RPO to non-investment-grade labs; receivables risk if API growth slows~\cite{ref:oracle_10q}. & Critical \\
Existential & Neo-clouds, GPU debt funds & Balance sheets levered on merchant silicon with 3--5 year life; no software ecosystem to buffer hardware deflation~\cite{ref:coreweave_s1}. & Critical \\
High impairment & Microsoft & Cut from both sides: \$190B/yr capex and the OpenAI round trip on one side; on the other, the largest seat-priced software base in the world (Office, Windows, Dynamics) deflating under agents, a consumer assistant that has lagged, customers paying to avoid Windows 11, and the sharpest exposure to the European sovereignty exit (France, Schleswig-Holstein, Denmark). AI receipts grow while the base that funds the capex shrinks. & High \\
High impairment & Amazon & \$200B/yr capex with the Anthropic round trip; balance sheet survives, cloud margins absorb any ASC 360-10 charges on the 2025--27 cohorts. & High \\
Defensive & Alphabet, Meta & In-house silicon (TPU, MTIA) avoids merchant margins; ad cash flows absorb overcapacity, and the ad channel is a demand pool the others lack. & Moderate \\
Structural & FLAPD hubs & Tariff squeeze and grid queues versus cheap-power corridors~\cite{ref:eirgrid_cru,ref:amsterdam}. & High \\
Structural & Seat-priced SaaS, hourly consultancies & Billable-unit collapse under outcome contracts. & High \\
\bottomrule
\end{tabularx}
\end{table}
\begin{table}[htbp]
\centering
\small
\caption{Beneficiaries}
\label{tab:winners}
\begin{tabularx}{\textwidth}{L{3.0cm}L{2.8cm}Xl}
\toprule
\textbf{Segment} & \textbf{Base} & \textbf{Engine} & \textbf{Horizon} \\
\midrule
Last-mile outcome operators & US, UK, G20 & Gain-share pricing; \$80--140B addressable at high gross margin on distilled models. & 2026--32 \\
Frontier labs, differentiated & US, Europe, China & Anthropic's enterprise and coding revenue is compounding fastest on reported figures and runs on partners' silicon rather than its own balance sheet; OpenAI carries the largest compute commitments against the largest reported losses; Mistral, DeepSeek and the open-weight labs win on cost, not revenue. If growth holds they are the largest winners; if it decays they are the counterparties on whom the hyperscalers' recoverability forecasts rest. & 2026--30 \\
Consumer and edge silicon & Apple, Qualcomm, AMD & 1.45B NPUs/yr absorbing routine inference locally. & 2026--30 \\
Hardware hedgers & Nvidia (edge, physical AI, open models) & Diversification away from cluster saturation; open-model releases commoditize the layer above the chip. & 2026--28 \\
China western corridors & Qinghai, Xinjiang, Gansu & \$0.02--0.035/kWh subsidizing embodied AI~\verify. & 2026--35 \\
Sovereign pragmatists & India, Gulf, Korea & DPI and sovereign models on stranded power and open weights. & 2026--32 \\
Grid and power industrials & Transformers, cooling, transmission & 36--48 month backlogs capture hardware spend first. & 2026--30 \\
\bottomrule
\end{tabularx}
\end{table}


\section{Timeline}
\begin{itemize}
    \item \textbf{2026--27:} startup-cohort failure crosses $\sim$80\%; SEC and auditor scrutiny of related-party cloud revenue; first \$50B+ impairment wave across hyperscalers and neo-clouds; secondary GPU spot prices fall by half or more.
    \item \textbf{2028--29:} human-text data wall binds for frontier training~\cite{ref:epoch_data2024}; CFOs cut experimental AI budgets; Big-5 capex retreats 30\%+ from the 2026--27 peak; neo-cloud restructurings; OCI renegotiates unfulfilled contracts.
    \item \textbf{2030:} 100+ TOPS local silicon standard; open-weight small models the enterprise default; frontier intelligence a low-margin utility; software margin concentrated in domain-specific integrators.
\end{itemize}

\section{What Would Change the Distribution}\label{sec:falsify}
The paper's result is a distribution, not a yes or no, so the tests below are stated as observations that would move the distribution decisively rather than as binary refutations. Two are sufficient on their own:
\begin{enumerate}
    \item \textbf{Revenue.} If calendar-2028 end-customer AI revenue across all channels exceeds \$400B --- i.e.\ $g\ge50\%$ has held for two more years from the 2026 base, measured on the same seven-component basis as Table~\ref{tab:basedecomp} --- the plateau path is on track and the paper's most likely outcome was wrong. Interim markers: calendar-2027 revenue near \$330B is the bottom-up path; above \$350B favors the build-out; the sector to watch is office assistants and agents, where the bottom-up ceiling (\$150B) and the IT-budget cap (\$270B) leave the most room to differ.
    \item \textbf{Impairment and retreat.} If by end-2030 cumulative AI-asset impairments across hyperscalers and neo-clouds under ASC 360-10 are below \$100B \emph{and} aggregate Big-5 capex has not fallen 30\% from its peak, the paper's most likely outcome was wrong.
\end{enumerate}
Other observables will move the distribution as they arrive (Appendix~\ref{app:update} gives the calendar and the re-issue rule), in rough order of leverage: the labs' calendar-2026 bookings and gross margins (the base and $\theta$); Microsoft's next paid-seat disclosure (the assistant ceiling); Data Center Watch's quarterly count (the deferral fraction); the operators' year-end useful-life notes (silicon life); Gartner's 2027 AI-share revision (the feasibility band). Softer expectations that do not decide the case, because their attribution is weak: BLS professional-services employment shows under 1.5\% net AI-attributable displacement; metered frontier API prices fall more than 80\% from 2026 while integrator revenue grows.

If the paper is right, the sequence is: 2027--28, lab growth decelerates below 2$\times$, the 2024--25 startup cohort thins, auditors scrutinize investor-funded compute contracts, first \$50B+ impairment wave; 2028--29, CFOs cut experimental budgets, capex retreats toward the \$450--550B the realized $g$ can carry, neo-cloud restructurings; 2030, 100+ TOPS local silicon is standard, open-weight small models the enterprise default, software margin with domain integrators.

\section*{Acknowledgments}
The first circulated version was reviewed by George Kountouris, Francis Pedraza, Matthew Fitzpatrick and Viraj Mehta, who are thanked as its first reviewers. The last-mile implementer's view --- that AI applied to a poorly run process automates the defects, and that the money moves to whoever fixes the process --- shaped Sections~\ref{sec:slack} and~\ref{sec:lastmile}; the commodities view --- economic versus accounting depreciation, the utilization wrinkle, and price converging to marginal cash cost --- is Section~\ref{sec:commodity}. Errors remain the author's. The paper is released so that a community can improve it as the data arrive: every input is in an open register with its evidence and its range, and anyone with a better-sourced number is invited to replace it.

\newpage
\appendix
\section{Model Specification}
All parameters triangular Triang(min, mode, max); \$B.
\begin{itemize}
    \item Capex history $I_{2023..2026}=161, 260, 448, 780$~\cite{ref:epoch_capex,ref:mlq2026}; forward paths Plateau (780 flat), Arms race (1000 then 1100), Rationalize (550 then 500), each held flat to 2035; and an Adaptive path in which next year's capex scales with last year's ratio of receipts to carry, damped (exponent 0.5) and bounded to \$400B--\$1.1T (Appendix~\ref{app:horizon}). No crystal ball chooses among the first three: the plateau is the base because it is the operators' stated 2026 level held constant; the arms race is their 2027 trajectory continued; the cut is what an impairment test would force. The adaptive path is what a rational operator would do with the same information the model has. Non-AI baseline $N_c=(161-A_{2023})(1+g_N)^{c-2023}$, $A_{2023}\sim$Triang(20, 30, 40), $g_N\sim$Triang(0.02, 0.06, 0.10). Grid energization cap Triang(1150, 1400, 1680) \$B/yr; siting deferral $f_{\text{delay}}\sim$Triang(0.10, 0.20, 0.35) from 2026; deferred capex carries $r$ only. Shell cost uplift $\kappa_{\text{site}}\sim$Triang(1.02, 1.05, 1.09); opex uplift $u\sim$Triang(0, 0.03, 0.06).
    \item Cohorts: $\phi\sim$Triang(0.45, 0.52, 0.60); silicon life $=1/(f_{\text{tr}}/L_{\text{tr}}+(1-f_{\text{tr}})/L_{\text{inf}})$ with $f_{\text{tr}}\sim$Triang(0.35, 0.45, 0.55), $L_{\text{tr}}\sim$Triang(2.5, 3.5, 4.5), $L_{\text{inf}}\sim$Triang(4, 5, 6); $L_{\text{sh}}\sim$Triang(15, 20, 25); $\ell_{\text{si}}\sim$Triang(0.5, 0.7, 0.9), $\ell_{\text{sh}}\sim$Triang(0.2, 0.4, 0.6); straight-line, mid-year.
    \item Carry: $c_{\text{opex}}\sim$Triang(0.059, 0.0665, 0.074), decomposed; $r\sim$Triang(0.090, 0.0945, 0.099), CAPM; $H_t$ per Eq.~\ref{eq:Ht}.
    \item Demand: $D_{2026}\sim$Triang(163, 179, 194) and $\theta\sim$Triang(0.77, 0.815, 0.86), both from the decompositions of Section~\ref{sec:decomp}; $D_t=\theta D_{2026}(1+g)^{t-2026}$, $g$ swept.
    \item Bottom-up demand: fifteen sectors per Table~\ref{tab:sectors}; logistic ramps; common factors on ceilings Triang(0.7, 1.0, 1.4) and ramps Triang(0.7, 1.0, 1.3); circular deduction Triang(8, 12, 18); $\gamma_{\text{sov}}$ applied to sovereign-exposed sectors.
    \item Secondary variables added at v25: economic depreciation rate $\delta\sim$Triang(0.5, 0.7, 0.9); inference utilization Triang(0.40, 0.45, 0.50) with a doubling case; token price index and marginal cash cost per million tokens (Section~\ref{sec:commodity}).
    \item Limitations: top-down parameters independent (understates spread); no demand saturation (favors the bull at high $g$); no SPV debt-service cliffs or grid-stranded cohorts (favors the bull); a single capex path per run rather than a policy response to realized $g$; investor-funded lab compute counted as receipts.
\end{itemize}


\section{Beyond 2030: The Horizon to 2035 and a Rational Operator}\label{app:horizon}
Ending at 2030 invites the objection that the deficit is a timing artifact: the stock is built ahead of demand, and by 2032 the demand has arrived. This appendix extends the same model to 2035. After 2030 the sector ceilings grow with the budgets that set them --- IT spending, advertising and users --- at a rate $g_C\sim$Triang(6, 10, 14)\%/yr, and the logistic ramps continue; capex paths are held at their 2028--30 levels. A fourth, \emph{adaptive}, path lets capex respond: each year's capex is last year's scaled by the ratio of receipts to carry, damped by a square root and bounded to \$400B--\$1.1T --- the behavior of an operator who reads the same numbers this paper does and neither ignores nor panics at them.

\begin{figure}[htbp]\centering\includegraphics[width=\textwidth]{f19_horizon2035.pdf}
\caption{Annual and cumulative surplus or deficit to 2035 by capex path (medians). A plateau reaches annual break-even only around 2035 and never recovers cumulatively within the horizon; the adaptive path is profitable from 2031 and recovers cumulatively in 2034--35.}\label{fig:horizon}\end{figure}

\begin{table}[htbp]
\centering\small
\caption{2031--2035: median annual $\Delta_t$ (\$B) with P(profit), and cumulative from 2027}
\label{tab:horizon}
\begin{tabular}{lccccc c}
\toprule
\textbf{Path} & \textbf{2031} & \textbf{2032} & \textbf{2033} & \textbf{2034} & \textbf{2035} & \textbf{Cum.\ 2027--35} \\
\midrule
Plateau & $-230$ (1\%) & $-190$ (4\%) & $-145$ (13\%) & $-95$ (26\%) & $-35$ (42\%) & $-1{,}550$ \\
Arms race & $-575$ (0\%) & $-575$ (0\%) & $-565$ (0\%) & $-540$ (0\%) & $-510$ (0\%) & $-4{,}260$ \\
Rationalize & $+80$ (83\%) & $+150$ (95\%) & $+220$ (98\%) & $+300$ (100\%) & $+385$ (100\%) & $+880$ \\
Adaptive & $+55$ (80\%) & $+165$ (99\%) & $+245$ (100\%) & $+290$ (100\%) & $+300$ (100\%) & $+505$ \\
\bottomrule
\end{tabular}
\end{table}

\begin{figure}[htbp]\centering\includegraphics[width=0.85\textwidth]{f20_adaptive.pdf}
\caption{The adaptive path. Capex falls to $\sim$\$400B by 2030--31 as receipts lag carry, then re-expands as receipts overtake it; the stock is right-sized to demand with a two-year lag.}\label{fig:adaptive}\end{figure}

Three readings. First, the plateau does not recover by 2035: infrastructure receipts grow from \$480B to \$910B, but the plateau's carry grows too --- replacement of short-lived silicon plus the shell cohorts still depreciating --- from \$725B to \$930B, so the annual gap is still open at the end of the horizon and the cumulative deficit of $\sim$\$1.55T is not repaid. ``By 2035 it all comes back'' is true for the technology and false for the plateau balance sheet. Second, the arms race never closes. Third, a 35\% cut or an adaptive operator turns profitable in 2030--31 and recovers cumulatively by 2034--35, at a stock that has fallen to roughly \$400--500B/yr of capex --- close to what Section~\ref{sec:results} identified as the level realized demand can carry. The adaptive path also answers a question the fixed paths cannot: how does the industry stay on the rails? By letting the operators' own recoverability tests do their work a year at a time, rather than by any single actor's failure. Growth $g$ is in this sense not an exogenous unknown but the outcome of decisions the operators will make with the same numbers, which is why the paper reports distributions and thresholds rather than a forecast.

\section{Sector Ceiling Derivations}\label{app:sectors}
Each ceiling is a Fermi chain to 2030; ranges become the triangular min/max in Table~\ref{tab:sectors}.
\begin{itemize}
    \item \textbf{Office assistants / agents.} 350M knowledge-worker seats $\times$ 35--55\% adoption $\times$ \$25--35/seat/month $\times$ 12 $=$ \$37--83B of seat revenue; agentic and workflow consumption at 1.5--3$\times$ seat revenue $=$ \$55--250B; ceiling capped by an AI share of 3--5\% of $\sim$\$5.5T global IT spend $=$ \$165--275B. Mode \$200B. $\theta=0.6$: routine tasks migrate to on-device NPUs and private small models (Section~\ref{sec:bottomup}).
    \item \textbf{Advertising uplift.} Meta and Google advertising $\sim$\$500B in 2026; AI ranking, creative and measurement lift yield by 10--20\% $=$ \$50--100B; AI-native conversational and generative formats \$30--150B. Mode \$150B. $\theta=1.0$: the infrastructure owner is the advertiser's platform.
    \item \textbf{Consumer subscriptions.} $\sim$1B monthly users $\times$ 5--10\% paid $\times$ \$20--25/month $\times$ 12 $=$ \$12--30B; free-tier advertising at \$30--80/user-year on 300--600M ad-eligible users $=$ \$10--50B; premium tiers and devices remainder. Mode \$80B.
    \item \textbf{Coding tools.} 29M developers~\cite{ref:octoverse} $\times$ 50--70\% paid $\times$ (\$20--200/month seat + \$1--3k/yr agent tokens) $=$ \$35--95B. Mode \$60B. Non-recurring: refactoring impulse $\alpha e^{-\lambda t}$ with maintenance tail $\beta\sim$Triang(0.05, 0.08, 0.12) of build tokens.
    \item \textbf{Search \& internal use.} Google search $\sim$\$250B; AI overviews and modes defend share rather than expand it; incremental monetizable AI search 5--15\% $=$ \$12--38B; hyperscaler internal productivity, fraud and recommendation uplift \$15--40B. Mode \$50B.
    \item \textbf{Financial services.} Banking and insurance IT and operations $\sim$\$1.5T; AI share 2--4\% $=$ \$30--60B, constrained by model-risk governance (SR~11-7, FINRA) and on-premise deployment ($\theta=0.75$). Mode \$40B.
    \item \textbf{Customer service / CX.} \$350B contact-center labor; 10--15\% displaced by 2030 $=$ \$35--55B of savings; vendors and integrators price at 50--80\% of savings $=$ \$25--45B. $\theta=0.5$: speech and RAG run on small models; 10--35:1 token ratio.
    \item \textbf{Public sector / defense.} Government IT $\sim$\$700B; AI share 4--8\% $=$ \$28--56B. Sovereign carve-out $\gamma$ applies.
    \item \textbf{Healthcare / clinical.} 1.1M US physicians $\times$ 60--80\% ambient-scribe adoption $\times$ \$400--600/month $=$ \$3--6B; global multiple 2--3$\times$; prior-authorization, coding, imaging triage \$10--25B. Mode \$35B. Market reports place clinical documentation at \$5B in 2026~\cite{ref:healthcare_rm}.
    \item \textbf{Autonomous fleets.} Autonomous-driving and fleet software forecasts of \$8B (2026) to \$32B (early 2030s)~\cite{ref:av_fbi}; 5--7 year chassis cycles; most inference on-vehicle ($\theta=0.45$).
    \item \textbf{Generative video / gaming / media.} Cloud GPU cost $\sim$\$0.30/user-hour against consumer willingness to pay; live generative content 100--300M users $\times$ \$5--10/month $=$ \$6--36B~\cite{ref:newzoo}.
    \item \textbf{Legal \& compliance.} \$950B legal services; software 2--3\% $=$ \$19--28B; verification rules and hallucination liability cap adoption. Mode \$20B.
    \item \textbf{Science / R\&D.} Pharma R\&D $\sim$\$300B $\times$ 4--8\% $=$ \$12--24B; materials, climate, engineering simulation \$5--20B. Mode \$25B.
    \item \textbf{Factory / robotics software.} Operational robot stock $\sim$4--5M~\cite{ref:ifr} $\times$ \$2--4k/yr software $=$ \$8--20B; bounded by Asian wage floors (\$1.05--1.60/hr) below robot TCO ($\sim$\$2.9/hr); inference at the edge ($\theta=0.45$).
    \item \textbf{Eldercare / assistive.} 160M OECD seniors $\times$ 15--25\% smart-home penetration $\times$ \$100--200/yr $=$ \$2--8B; 85\% on-device for privacy ($\theta=0.15$). 90\% of eldercare cost is physical labor software cannot displace.
    \item \textbf{Edge siphon (cross-cutting).} $350\text{M}\times15\times250\times500=6.6\times10^{14}$ routine tokens/yr; at \$0.10--0.50/M $=$ \$66--330B of potential cloud revenue; 60--85\% executed locally $\Rightarrow$ \$40--280B/yr never billed. This is the physical basis for $\theta<1$ in the office assistant, CX, robotics and eldercare rows.
    \item \textbf{Spatial video tokens (embodied AI).} $4\text{ cameras}\times30\text{ fps}\times3600\text{ s}\times62.5\text{ tokens/frame}=2.7\times10^{7}$ tokens per camera-hour, $\sim$5.4$\times10^{8}$ per robot per 20-hour day; favors \$0.02--0.035/kWh corridors over cheap chips.
\end{itemize}


\section{Error Propagation Across Chains}\label{app:errors}
The paper combines estimates in three ways, and each has its own propagation rule. All three are checked below against the Monte Carlo, which does not assume any of them.

\subsection{Rule 1: within a chain, factors multiply, log-variances add}
A single chain is a product, $Y=\prod_i X_i$. With independent factors,
\begin{equation}
\sigma^{2}_{\ln Y}=\sum_i\sigma^{2}_{\ln X_i},
\end{equation}
so \emph{relative} errors add in quadrature and the loose factors dominate. For a triangular input with min $a$ and max $c$, $\sigma_{\ln X}\approx\ln(c/a)/4.9$. In this paper every input is tighter than a factor of two: the supply-side inputs have $\sigma_{\ln}$ between 0.06 (cost of capital) and 0.17 (energization lag); the sector ceilings have 0.15--0.25. The chain is where individual expertise enters, and where a referee should attack one factor at a time.

\subsection{Rule 2: across chains that are summed, absolute variances add --- plus the common factor}
Sectors are added, not multiplied: $D=\sum_s R_s$. For independent sectors,
\begin{equation}
\sigma^{2}_{D}=\sum_s\sigma^{2}_{R_s}
\qquad\Longleftrightarrow\qquad
\sigma^{2}_{\ln D}\approx\sum_s w_s^{2}\,\sigma^{2}_{\ln R_s},\quad w_s=\frac{R_s}{D},
\label{eq:sumrule}
\end{equation}
so the relative error of the sum is the \emph{share-weighted} root-sum-square of the sector relative errors. Because the shares $w_s$ are squared, a sector with 5\% of the total contributes 1/25 as much variance as one with 25\%, even at equal relative uncertainty; and because there are many sectors, the idiosyncratic error of the sum is far smaller than that of any one sector. This is the averaging that rewards detail.

What does not average is anything the sectors share. If a common factor $M$ (the $M_C$ of Section~\ref{sec:bottomup}) multiplies every ceiling --- macro adoption, a recession, a regulatory shock --- then $D=M\sum_s R_s'$ and
\begin{equation}
\sigma^{2}_{\ln D}\approx\sigma^{2}_{\ln M}+\sum_s w_s^{2}\,\sigma^{2}_{\ln R_s'}.
\label{eq:common}
\end{equation}
The common factor adds its full variance once, unweighted, and cannot be diversified away. In practice it dominates: the residual uncertainty of a well-decomposed demand estimate is almost entirely the shared factor, not the sector details.

\subsection{Rule 3: supply minus demand, absolute variances add}
The result is a difference, $\Delta=\text{Receipts}-H$. Supply and demand share no parameters in this model, so
\begin{equation}
\sigma^{2}_{\Delta}=\sigma^{2}_{\text{Receipts}}+\sigma^{2}_{H}\qquad(\text{plus }-2\,\text{Cov if correlated}),
\end{equation}
and the strength of the conclusion is the $z$-score $\Delta/\sigma_{\Delta}$. Two things follow. First, because the uncertainties add rather than the medians, a gap must exceed roughly $2\sigma_{\Delta}$ before the sign is a Fermi-robust result. Second, if capex and demand were positively correlated in reality (nobody spends into flat demand), the covariance term would \emph{reduce} $\sigma_{\Delta}$ and sharpen the sign; we do not take that credit.

\subsection{Numerical check, 2030, Plateau}
\begin{table}[htbp]
\centering\small
\caption{Propagation rules versus the Monte Carlo (2030, Plateau, $N=100{,}000$)}
\label{tab:propcheck}
\begin{tabular}{lrr}
\toprule
\textbf{Quantity} & \textbf{Predicted by rule} & \textbf{Observed in MC} \\
\midrule
\multicolumn{3}{l}{\textit{Rule 2, sectors without the common factor}} \\
$\sum_s w_s^{2}\sigma^{2}_{\ln R_s}$ (Table below) & 0.0026 & \\
$\sigma_{\ln D}$, idiosyncratic sum & 0.051 & 0.051 \\
\multicolumn{3}{l}{\textit{Rule 2 with common ceiling factor, $\sigma_{\ln M}=0.14$}} \\
$\sigma_{\ln D}=\sqrt{0.0026+0.0195}$ & 0.149 & 0.133$^{\dagger}$ \\
\multicolumn{3}{l}{\textit{Rule 1, supply chain}} \\
$\sigma_{\ln H_{2030}}$ & --- & 0.047 \\
\multicolumn{3}{l}{\textit{Rule 3, difference}} \\
$\sigma_{\text{Receipts}}$, $\sigma_{H}$ (\$B) & 65, 34 & 65, 34 \\
$\sigma_{\Delta}=\sqrt{65^{2}+34^{2}}$ (\$B) & 73 & 73 (Cov $\approx-4$) \\
Median $\Delta_{2030}$ (\$B) & & $-259$ \\
$z=\Delta/\sigma_{\Delta}$ & & $-3.5$ \\
P($\Delta_{2030}<0$) & & 0.99 \\
\bottomrule
\end{tabular}
\\[2pt]{\footnotesize $^{\dagger}$Slightly below prediction because the logistic has not fully reached the ceiling by 2030, so the common ceiling factor passes through at less than unit elasticity.}
\end{table}

\begin{table}[htbp]
\centering\small
\caption{Sector contributions to Eq.~\ref{eq:sumrule}, 2030 (idiosyncratic, common factor removed)}
\label{tab:sectorvar}
\begin{tabular}{lccc}
\toprule
\textbf{Sector} & \textbf{Share $w_s$} & $\sigma_{\ln R_s}$ & $w_s^{2}\sigma^{2}_{\ln R_s}$ \\
\midrule
Office assistants / agents & 0.21 & 0.16 & 0.0011 \\
Advertising uplift & 0.18 & 0.11 & 0.0004 \\
Consumer subscriptions & 0.11 & 0.16 & 0.0003 \\
Coding tools & 0.09 & 0.16 & 0.0002 \\
Search \& internal use & 0.07 & 0.18 & 0.0001 \\
Customer service / CX & 0.05 & 0.17 & 0.0001 \\
Financial services & 0.05 & 0.17 & 0.0001 \\
Remaining eight sectors & 0.01--0.04 & 0.15--0.25 & 0.0002 total \\
\midrule
Sum & 1.00 & & 0.0026 $\Rightarrow\sigma_{\ln D}=0.051$ \\
\bottomrule
\end{tabular}
\end{table}

The lesson of the table is the one the method section promised. Fifteen sectors, each uncertain by 11--25\%, sum to an idiosyncratic uncertainty of 5.1\%; assistants carry 40\% of that now that the advertising and consumer ceilings are decomposed. What remains --- three-quarters of the total demand variance --- is the common adoption factor at 14\%, which sector detail cannot reduce but which Section~\ref{sec:decomp} splits into an IT-budget part (5\%), a price-deflation part (8\%) and a residual capability part (10\%). The supply side, a chain of seventeen tight inputs, is uncertain by 5\%. Subtracting, the 2030 gap of \$259B against $\sigma_{\Delta}=\$73$B is a 3.5$\sigma$ result: robust by the Fermi standard, not a certainty, and its residual doubt lives in the residual capability factor and the assistant ceiling. Those are the two axes of Figure~\ref{fig:probmap}.

\section{Other Derivations}
\subsection{Carry at the modes, 2028 Plateau}
AI cohorts in service: 2023: 0; 2024: 95; 2025: 283; 2026: 615; 2027: 615; 2028: 308 (half energized). Gross $\approx$1{,}915. Silicon 55\%: depreciation $1{,}053/4=263$; shell 45\%: $862/20=43$. Opex $0.08\times1{,}915=153$. Net book $\approx$1{,}280; return $0.105\times1{,}280=134$. $H_{2028}\approx590$; Monte Carlo median 565.
\subsection{Enterprise benefit requirement, 2028}
$0.7\times570\approx400$; $\times(1+1.5)=1{,}000$; $\times1.2\approx1{,}200$. Payroll shares: $1.2/30=4\%$ global; $1.2/9=13\%$ US-alone.
\subsection{Slack dissipation}
$\tau(1-0.65)(1-0.20)=0.28\,\tau$; at $\tau=0.15$: 4.2\% of payroll.
\subsection{Contact-center ratio}
Labor \$350B; displacement 5--10\% $=$ \$17--35B; tokens \$1--2B; ratio 10--35:1.
\subsection{Break-even revenue levels}
Plateau 2030 at $g^{*}=49\%$: $179\times1.49^{4}\approx880$; at $\theta=0.815$, infrastructure receipts $\approx720$, against $H_{2030}\approx710$.

\newpage
\section{Every Equation, Term by Term}\label{app:equations}
Each numbered equation in the paper, restated and explained one symbol at a time, with units and a worked number where one helps. Symbols follow the glossary table (Appendix~\ref{app:glossary}); a few notational conventions come first.

\subsection*{Conventions}
Lower-case $\ell$ is an energization lag (a fraction); upper-case $L$ is a life in years --- they are different quantities and the letters are deliberately different. Lower-case $k_s$ is a sector's logistic ramp rate (per year); $k$ in the Arrhenius and Black equations is Boltzmann's constant --- the only place a letter is reused, and it is flagged where it occurs. $f_{\text{tr}}$ is the training share of the silicon fleet; $f_{\text{delay}}$ is the siting deferral. $M_C$ and $M_k$ are the common factors; $m$ is the enterprise benefit margin. $\omega(e)$ is an edge weight in the circular-flow graph; $W$ is the payroll pool; $w_s$ is a sector share. Triang($a$, $b$, $c$) always means a triangular distribution with minimum $a$, mode $b$, maximum $c$. Money is in billions of US dollars per year unless stated; rates are per year; shares are fractions between 0 and 1.

\subsection*{Equation~\ref{eq:var}: how uncertainties combine in a product}
\[\sigma^{2}_{\ln Y}=\sum_{i=1}^{n}\sigma^{2}_{\ln X_i}.\]
$Y$ is a quantity computed as a product of $n$ factors $X_1\cdots X_n$ (a Fermi chain). $\ln$ is the natural logarithm; taking logs turns the product into a sum, $\ln Y=\sum\ln X_i$. $\sigma_{\ln X_i}$ is the standard deviation of $\ln X_i$, which for modest uncertainties is the same as the \emph{relative} uncertainty of $X_i$ (a $\sigma_{\ln}$ of 0.10 means ``about $\pm$10\%''). For independent factors, variances add: the relative uncertainty of the product is the square root of the sum of the squares of the factors' relative uncertainties. Worked: three factors at 0.10, 0.15 and 0.30 give $\sqrt{0.01+0.0225+0.09}=\sqrt{0.1225}=0.35$, not $0.55$; the loosest factor dominates.

\subsection*{Equation~\ref{eq:Ht}: the annual carry of the installed base}
\[H_t=\sum_{c\le t}\sum_{j\in\{\text{si},\text{sh}\}}\left[\frac{G_{c,j}}{L_j}+c_{\text{opex}}\,G_{c,j}+r\,G_{c,j}\!\left(1-\frac{t-c+\tfrac12}{L_j}\right)\right]\mathbf{1}\!\left[t-c+\tfrac12<L_j\right]+\;r\,\text{Deferred}_t.\]
$H_t$ is the carry in year $t$ (\$B). The outer sum runs over every cohort year $c$ up to and including $t$; the inner sum over the two components $j$ of each cohort: silicon (si) and shell (sh). $G_{c,j}$ is the in-service gross value of that component of that cohort (\$B): AI-attributable capex $A_c$ times the component's share ($\phi$ for silicon, $1-\phi$ for shell), times the energized fraction ($\ell_j$ in the spend year, 1 afterwards), times the site cost uplift $\kappa_{\text{site}}$ for shell. Inside the bracket, three costs: \emph{depreciation} $G_{c,j}/L_j$, straight-line over the component's life $L_j$; \emph{facility opex} $c_{\text{opex}}G_{c,j}$ (times $(1+u)$ where dry cooling applies); and \emph{return on net book} $r$ times what is left undepreciated, where $t-c+\tfrac12$ is the cohort's age at mid-year and $1-\text{age}/L_j$ is the undepreciated fraction. The indicator $\mathbf{1}[\cdot]$ is 1 while the cohort is younger than its life and 0 after it has fully depreciated, so retired cohorts drop out. The final term charges the cost of capital on capital that is deferred (by the grid cap or siting) and not yet earning. Worked, 2028 plateau at the modes: in-service AI gross about \$1.9T, of which silicon \$1.05T $/4.2$ yr $=$ \$250B of depreciation and shell \$0.86T $/20$ yr $=$ \$43B; opex $0.0665\times1.9\text{T}=\$126$B; return $0.0945\times$ net book $\sim$\$1.3T $=\$123$B; total $\approx$ \$540B, matching Table~\ref{tab:carry}.

\subsection*{Effective silicon life (harmonic mean)}
\[L_{\text{si}}=\frac{1}{f_{\text{tr}}/L_{\text{tr}}+(1-f_{\text{tr}})/L_{\text{inf}}}.\]
Because depreciation is $1/L$ per year, a fleet's average depreciation rate is the share-weighted average of $1/L$, and the effective life is its reciprocal. At $f_{\text{tr}}=0.45$, $L_{\text{tr}}=3.5$, $L_{\text{inf}}=5$: $0.45/3.5+0.55/5=0.129+0.110=0.239$, so $L_{\text{si}}=4.2$ years.

\subsection*{The top-down demand path}
\[D_t=\theta\,D_{t_0}(1+g)^{t-t_0},\qquad t_0=2026.\]
$D_{t_0}$ is calendar-2026 end-customer AI revenue (\$B); $g$ the assumed constant annual growth (0.49 means 49\%/yr); $(1+g)^{t-t_0}$ compounds it over $t-t_0$ years; $\theta$ converts end-customer revenue to infrastructure receipts. Worked: $179\times1.49^{4}=880$ end-customer, $\times0.815=717$ receipts in 2030. The break-even $g^{*}$ is the $g$ at which this equals $H_t$; it is found by bisection in the code.

\subsection*{The logistic sector curve}
\[R_s(t)=\frac{C_s}{1+\left(C_s/R_{s,0}-1\right)e^{-k_s(t-t_0)}}.\]
At $t=t_0$ the exponential is 1 and the expression reduces to $C_s/(C_s/R_{s,0})=R_{s,0}$: the curve starts at the 2026 value. As $t$ grows the exponential shrinks toward 0 and $R_s(t)\to C_s$: it approaches the ceiling. The bracket $(C_s/R_{s,0}-1)$ sets how far below the ceiling the sector starts; $k_s$ sets how fast the gap closes --- half the remaining distance is covered in about $\ln 2/k_s$ years. Worked, office assistants: $R_{s,0}=30$, $C_s=150$, $k_s=0.8$: in 2030, $e^{-3.2}=0.041$, bracket $=4$, denominator $=1.16$, $R_s=129$. The sector's growth rate is the slope of this curve divided by its value, which falls from about 60\%/yr in 2027 to about 12\%/yr in 2030. Sums: $D^{\text{BU}}_t=\sum_s R_s(t)$ and $\text{Receipts}_t=\sum_s\theta_sR_s(t)$, with common factors $M_C$ multiplying every $C_s$ and $M_k$ every $k_s$, and sovereign-exposed sectors multiplied by $(1-\gamma_{\text{sov}})(1-g_{\text{EU}}\cdot\text{ramp})$.

\subsection*{The annual surplus or deficit}
\[\Delta_t=\text{Receipts}_t-F_{\text{circ}}-\kappa\,H_t.\]
Receipts less round-trip revenue $F_{\text{circ}}$ (\$B) less the carry, with $\kappa$ the sovereignty capital multiplier (1 in the base case, 1.2 in the overlay). Positive is a surplus. P(profit) is the fraction of Monte Carlo trials in which $\Delta_t>0$; the $z$-score is $\text{median}(\Delta_t)/\text{sd}(\Delta_t)$.

\subsection*{Realized enterprise benefit}
\[B_{\text{realized}}=W\,\tau\,(1-\sigma_{\text{slack}})(1-\nu_{\text{verify}})-C_{\text{severance}}.\]
$W$ is the payroll pool (\$T); $\tau$ the nominal fraction of work time AI saves; $(1-\sigma_{\text{slack}})$ the fraction of that saving not absorbed as slack; $(1-\nu_{\text{verify}})$ the fraction not consumed by verification; $C_{\text{severance}}$ one-time separation costs. Worked at the modes: $\tau\times0.35\times0.80=0.28\,\tau$; with $\tau=0.15$, 4.2\% of payroll --- \$1.26T on the \$30T global pool before severance.

\subsection*{The enterprise benefit requirement}
\[B_{\text{req}}=S_{\text{ent}}(1+m),\qquad S_{\text{ent}}=H_{\text{ent}}(1+\mu_{\text{integ}}),\qquad H_{\text{ent}}=f_{\text{ent}}H_t.\]
$f_{\text{ent}}$ ($\approx$0.7) is the enterprise share of the carry that enterprises must fund; $\mu_{\text{integ}}$ the integration spend per dollar of software and tokens; $m$ the required margin over cost. Worked, 2028: $0.7\times540=378$; $\times2.5=945$; $\times1.2\approx1{,}130$, i.e.\ $\sim$\$1.1--1.2T.

\subsection*{Propagation across summed chains (Appendix~\ref{app:errors})}
\[\sigma^{2}_{\ln D}\approx\sigma^{2}_{\ln M}+\sum_s w_s^{2}\,\sigma^{2}_{\ln R_s}.\]
$D=M\sum_sR_s$ is a sum of sectors with a common multiplier $M$ ($=M_C$ here). $w_s=R_s/D$ is sector $s$'s share. The relative variance of a sum of independent terms is the share-weighted (squared) sum of their relative variances --- small sectors contribute almost nothing --- plus the common factor's variance, which enters once, unweighted, because it moves every term together. Worked: fifteen sectors at 0.11--0.25 with the shares of Table~\ref{tab:sectorvar} sum to 0.0026 ($\sigma=0.051$); adding $\sigma_{\ln M}=0.14$ gives $\sqrt{0.0026+0.0195}=0.149$.

\subsection*{Propagation across the difference}
\[\sigma^{2}_{\Delta}=\sigma^{2}_{\text{Receipts}}+\sigma^{2}_{H}-2\,\text{Cov}(\text{Receipts},H).\]
Absolute (not relative) variances add for a difference; the covariance is near zero here because supply and demand share no inputs. Worked: $\sqrt{65^2+34^2}=73$; $z=-259/73=-3.5$.

\subsection*{The incumbent's revenue ratio}
\[\frac{\text{Rev}_{2030}}{\text{Rev}_{2026}}=(1+g_{\text{seat}})^{4}\,(1-\text{EU})\,\bigl[1-\epsilon\bigl((1+p)^{4}-1\bigr)\bigr]\,(1+p)^{4}.\]
Seats fall with the agent-driven trend $(1+g_{\text{seat}})^4$, the sovereignty exit $(1-\text{EU})$, and the price-driven loss $\epsilon$ times the cumulative price rise; price rises by $(1+p)^4$. Worked at the modes: $0.97^4=0.885$; $\times0.92=0.814$; price rise $1.08^4-1=0.36$, seat loss $0.35\times0.36=0.126$, factor $0.874$; $\times1.36$: ratio $=0.814\times0.874\times1.36=0.97$ --- revenue about flat.

\subsection*{Arrhenius acceleration}
\[AF=\exp\!\left[\frac{E_a}{k}\left(\frac{1}{T_1}-\frac{1}{T_2}\right)\right].\]
$AF$ is the factor by which a temperature-driven failure rate rises from absolute temperature $T_1$ to $T_2$ (kelvin); $E_a$ the activation energy in electron-volts; $k$ Boltzmann's constant, $8.617\times10^{-5}$ eV/K. Worked: $E_a=0.7$ eV, $T_1=333$ K (60$^{\circ}$C), $T_2=343$ K (70$^{\circ}$C): $\frac{0.7}{8.617\times10^{-5}}(1/333-1/343)=8124\times8.76\times10^{-5}=0.71$; $AF=e^{0.71}=2.0$. Failure rate doubles per 10$^{\circ}$C in this range.

\subsection*{Black's equation}
\[\text{MTTF}\propto J^{-n}\,e^{E_a/kT}.\]
Mean time to failure by electromigration falls with current density $J$ raised to a power $n$ (typically about 2) and rises with the same Arrhenius exponential; high current and high temperature together shorten life.

\subsection*{Cost of capital (CAPM and WACC)}
\[r_e=r_f+\beta\,\text{ERP},\qquad r=\frac{E}{V}r_e+\frac{D}{V}r_d(1-T_{\text{tax}}).\]
Cost of equity: risk-free rate plus beta times the equity risk premium. WACC: equity weight times cost of equity plus debt weight times after-tax cost of debt. Worked: $4.3+1.2\times5.0=10.3\%$; $0.82\times10.3+0.18\times6.5\times0.79=8.45+0.92=9.4\%$.

\subsection*{Internal rate of return}
\[\sum_{t=0}^{n}\frac{F_t}{(1+d)^{t}}=0.\]
$F_t$ are the project's cash flows (negative outlay at $t=0$, positive returns after); $d$ is the discount rate. The IRR is the $d$ that makes the discounted flows sum to zero. A hurdle rate is the minimum IRR a firm accepts.

\section{Every Input, Appraised}\label{app:inputs}
Fermi's discipline is to do one's best on each number and to show the work so that anyone with better information can replace it. This appendix does that for every input in the register. For each: what is actually known and from whom; the range chosen and why; how far the source can be trusted, including where a company has a reason to shade the figure; what would improve it; and, where we have a view, what we think is likely to happen. Ranges are triangular (min / mode / max) unless stated. The register (\texttt{parameters.json}) is the place to substitute a different opinion; the rebuild script shows what it does.

A general caution that applies throughout, and a request. Corporate disclosures in this area are made by parties with a stake in the answer: hyperscalers choose accounting lives; labs choose run-rate months and revenue bases; forecasters sell the forecasts they publish; opposition groups count victories. None of that makes a figure false or implies bad faith, and nothing below is a claim about any company's conduct; it is the ordinary observation that a reader should ask ``who benefits from this number being large (or small)?'' and weight accordingly. Where we note an incentive, that is what it is --- an incentive, not an accusation. The request follows from it: the companies concerned could remove most of this appendix by disclosing a few numbers they already have --- paid seats and their trend, AI revenue on a stated basis, compute cost as a share of lab revenue, accelerator lives in service. Better guidance would cost them little and would replace inference with fact in every analysis of this kind, including those of their own lenders and auditors.

\subsection*{Supply side}

\paragraph{Capex history, $I_c$ (161, 260, 448, 780).} Known well. 2023--25 are Epoch AI's reconstruction from XBRL filings (cash PP\&E plus new finance leases), which is the best available basis and reconciles with company statements. 2026 is guidance, and guidance has been revised upward all year; companies have an incentive to guide high (signals commitment to investors and to each other) and then spend it. Improve with: the October actuals and the Q4 10-K cash-flow statements; note whether finance leases are growing faster than cash capex, which moves cost off the capex line. Likely: 2026 lands within \$40B of guidance; 2027 guidance will be the first real signal of whether the plateau or the arms race is the path.

\paragraph{AI capex in 2023, $A_{2023}$ (20 / 30 / 40).} Inferred, not disclosed. Total capex was flat in 2023 while H100 purchases began; Nvidia's data-center revenue that year (\$47B) and the Big-5's share of it bound the figure. The range is honest about the inference. Improve with: any company-level statement about 2023 AI spend. Likely: near the mode; it matters mainly because it pins the baseline.

\paragraph{Non-AI baseline growth, $g_N$ (2 / 6 / 10 \%/yr).} Judgment. Fulfillment capex is cyclical, ordinary cloud refresh tracks non-AI cloud revenue ($\sim$15\%/yr) at a lower capital intensity. Improve with: Amazon's segment capex, Microsoft's and Alphabet's statements about non-AI infrastructure. Likely: 4--8\%; the low end if capital is being diverted to AI.

\paragraph{Silicon share, $\phi$ (0.45 / 0.52 / 0.60).} Anchored on Nvidia's data-center revenue against AI-attributable capex, plus custom silicon, memory and networking. The trap is dividing by total capex rather than AI capex. Companies do not disclose the split; Nvidia's share of Big-5 purchases is estimated. Improve with: any hyperscaler statement of ``servers and equipment'' versus ``land and buildings'' in the PP\&E note (Meta and Microsoft give this). Likely: the mode; rising if custom silicon displaces shell-heavy new builds, falling if power infrastructure costs keep rising.

\paragraph{Training share of fleet, $f$ (0.35 / 0.45 / 0.55).} Judgment. Frontier training clusters are a minority of installed accelerators by count and a larger share by capital. No company discloses it. Improve with: statements about inference versus training capacity (Microsoft has made them qualitatively). Likely: falling over time as inference scales.

\paragraph{Training-cluster life, $L_{\text{tr}}$ (2.5 / 3.5 / 4.5 yr).} Physics and field data, not accounting. Arrhenius acceleration, the 24/7 duty cycle, Meta's reported failure rates, and the economics of the power slot argue for the short end; A100s still rented argue for the long. The companies' accounting lives (5--6 years) are longer than this, and the companies have an incentive to keep them long (each year added to useful life reduces reported depreciation by billions). Amazon's 2025 shortening to five years, ``due to the pace of AI,'' is the one disclosure that cuts the other way and is the most credible datum here. Improve with: the December 10-K useful-life notes; any field reliability data; secondary-market prices for H100s. Likely: the mode; an impairment or life shortening at one operator would move the others.

\paragraph{Inference-fleet life, $L_{\text{inf}}$ (4 / 5 / 6 yr).} Bounded above by the accounting lives (a company cannot credibly claim its assets last longer than it depreciates them) and below by the training life. Same corporate incentive as above. Likely: the mode.

\paragraph{Shell life, $L_{\text{sh}}$ (15 / 20 / 25 yr).} Industry norm; low leverage on the result. Improve only if evidence emerges that AI halls are being rebuilt for new density standards faster than that.

\paragraph{Energization lags, $\ell_{\text{si}}$ (0.5 / 0.7 / 0.9), $\ell_{\text{sh}}$ (0.2 / 0.4 / 0.6).} Judgment from build times: silicon racks within one to three quarters, shell takes twelve to twenty-four months. Low leverage after decomposition. Improve with: any disclosed construction-in-progress balances (the PP\&E note shows ``construction in progress,'' which is exactly un-energized capital).

\paragraph{Facility opex share, $c_{\text{opex}}$ (5.9 / 6.65 / 7.4 \% of gross).} Decomposed from PUE, tariff, utilization, capex per MW, staffing, tax, network. The tariff and PUE are public; staffing and maintenance are judgment. Improve with: any operator's disclosed cost per MW-year; colocation providers' (Equinix, Digital Realty) unit costs are the nearest public proxy. Likely: rising, as tariffs rise and dry cooling spreads. \emph{Second anchor (v27):} Goldman Sachs, citing Epoch AI, uses \$836M per GW-year of non-depreciation operating cost, about 2\% of capital --- the paper's power line alone. The 4.6\% gap (staffing, property tax and insurance, network, other) is the one input on which the two supply-side frameworks differ materially (Section~\ref{sec:gs}, Table~\ref{tab:gs_reconcile}); it is worth \$85--115B/yr of carry in 2028--30. Flagged: adopting the lower anchor would be a material re-issue under Appendix~\ref{app:update} and awaits a sourced operator figure.

\paragraph{Cooling opex uplift, $u$ (0 / 3 / 6 \%).} Engineering estimate of dry versus evaporative cooling. Improve with: the fraction of new capacity in water-constrained jurisdictions. Likely: the mode, rising with siting friction.

\paragraph{WACC, $r$ (9.0 / 9.45 / 9.9 \%).} CAPM with Damodaran inputs. The one input that is honestly well known; changes with rates. Improve with: the January update.

\paragraph{Grid energization cap (1150 / 1400 / 1680 \$B/yr).} 22--40 GW/yr of connectable new AI load, globally, at \$38--55M/MW. Weakly sourced: assembled from utility and queue reports, not a single dataset. Does not bind on a plateau, so its looseness does not matter unless capex exceeds $\sim$\$1.3T. Improve with: a consolidated GW-per-year connection figure from the grid operators. \emph{High anchor (v27):} Goldman Sachs's semiconductor build implies 20, 35 and 57 GW of AI deployments worldwide in 2026, 2027 and 2028~\cite{ref:gs_roic2026}; the 2028 figure is above this range's top and is the observable to test against utility energization data.

\paragraph{Siting deferral, $f_{\text{delay}}$ (10 / 20 / 35 \%).} Calibrated on Data Center Watch's \$130B per quarter of blocked or delayed projects against a pipeline of order \$1T. Data Center Watch tracks opposition outcomes and has an incentive to count broadly; its figures are project values, not annual spend, and most delayed projects eventually proceed elsewhere. The range discounts for both. Improve with: an independent count; the fraction of delayed projects that were ultimately built. Likely: rising through 2027 as moratoria spread, then falling as developers move to jurisdictions that want them.

\paragraph{Site cost uplift, $\kappa_{\text{site}}$ (1.02 / 1.05 / 1.09).} Judgment: community-benefit agreements, water mitigation, private generation, relocation premia. Improve with: disclosed terms of any large community-benefit agreement.

\paragraph{Sovereignty capital multiplier, $\kappa_{\text{sov}}$ (1.1 / 1.2 / 1.3; overlay only).} Engineering judgment about utilization and scale of smaller in-jurisdiction facilities. No regulator publishes it. Kept out of the base case for that reason.

\subsection*{Demand side}

\paragraph{The 2026 base, $D_{2026}$ (163 / 179 / 194).} Seven components. The two largest (Anthropic, OpenAI) rest on run-rates the companies chose to disclose, in the months they chose to disclose them, on a basis (gross, for Anthropic) that flatters the figure; both companies were raising capital at the time. Run-rates overstate calendar bookings. The hyperscaler component depends on what Microsoft counts in its ``AI business,'' a basis it has not fully defined. The advertising component is the least observable; Meta and Google attribute growth to AI on their calls without publishing the decomposition. The overlap correction is our own. The incentive on every component points upward. Improve with: January's 10-Ks (audited), an OpenAI S-1, and any Anthropic audited figure. Likely: calendar-2026 bookings near the mode; the run-rate figures will look generous in retrospect.

\paragraph{Cloud-capture share, $\theta$ (0.77 / 0.815 / 0.86), $\theta_{\text{lab}}$ ($\sim$0.635).} OpenAI's reported margin and Azure spend are leaked, unaudited figures; Anthropic's margin is undisclosed. The labs, raising capital, have the ordinary incentive to present improving margins. Improve with: the S-1; any audited lab financials. Likely: $\theta_{\text{lab}}$ falls over time as the labs' margins improve --- which is good for the labs and bad for the infrastructure owners, and is the paper's central value-capture argument in one number.

\paragraph{Circular flows, $F_{\text{circ}}$ (8 / 12 / 18).} From related-party notes and press. Small; kept for completeness and because readers will look for it. Improve with: the 10-K notes.

\paragraph{Sector bases, $R_{s,0}$.} Calibrated to sum to the observed base; not independent evidence. Individual sector bases are judgment allocations of a known total. Improve with: any vendor's sector revenue disclosure.

\paragraph{Sector ceilings, $C_s$.} The paper's most consequential judgments. Each is a chain (Appendix~\ref{app:sectors}). The two that matter: office assistants (100 / 150 / 230), anchored on Microsoft's disclosed 20M paid seats --- Microsoft discloses seat counts at its discretion and has not committed to a regular series --- and on an IT-budget hard cap; and advertising (105 / 136 / 170), anchored on Meta's disclosed lifts --- disclosed by a company with an interest in AI being credited for ad growth. Both have an upward incentive at the source and the ranges allow for it. The remaining thirteen are smaller and their errors average. Improve with: Copilot seats each quarter; Gartner's application-software category data; Meta's and Google's ad-revenue decomposition if ever given. Likely: assistants at or below the mode unless agents are trusted with whole workflows; advertising near the mode.

\paragraph{Ramp rates, $k_s$.} Judgment about time-to-saturation by sector, calibrated so that 2027 growth matches the labs' guided deceleration. Improve with: 2027 actuals.

\paragraph{Sector cloud shares, $\theta_s$.} Judgment about on-device, on-premise and integrator leakage by sector. The edge-siphon chain and the last-mile evidence support the shape; the levels are estimates. Improve with: any data on the share of enterprise inference running outside hyperscaler clouds (surveys exist; none are authoritative).

\paragraph{Common factors, $M_C$ (0.7 / 1.0 / 1.4), $M_k$ (0.7 / 1.0 / 1.3).} Judgment about how much macro and capability background the sectors share. This is where the residual doubt lives and it is stated as such. Improve only by experience.

\paragraph{Post-2030 ceiling growth, $g_C$ (6 / 10 / 14 \%/yr).} Generic; deliberately generous to the build-out. Should be tied to the torn-up budget pool (seat deflation, services shrinkage) in a later version, which would lower it.

\paragraph{European displacement, $g_{\text{EU}}$ (3 / 5 / 8 \% of sovereign-exposed demand by 2030).} From the sovereignty survey (a vendor of private email, with a stake in the answer), France's mandate, and the German state exits. Probably conservative given a national government has now moved. Improve with: procurement announcements; the 2027 survey. Likely: above the mode.

\paragraph{Sovereignty carve-out, $\gamma_{\text{sov}}$ (15 / 25 / 35 \%; overlay only).} Judgment; kept as an overlay.

\paragraph{Residual adoption factor (swept, not ranged).} The one input the paper refuses to assign a range to, because no 2026 data bear on it. It is whether agents become trustworthy with whole workflows and how fast enterprises reorganize around them. Our view, for what it is worth: the sector chains already assume substantial adoption; a residual factor of 1.5 on top of them --- what the plateau needs --- would mean the largest and fastest enterprise software adoption in history. Possible; not the base case.

\subsection*{Feasibility band}

\paragraph{Software and services pools (1440, 1870 \$B).} Gartner, whose forecasts moved 3.4 points on the same year within five months and who sells the forecasts. The levels are more reliable than the growth rates. Improve with: the October update; cross-check against IDC.

\paragraph{Seat share of software (0.50 / 0.55 / 0.60).} Judgment. Improve with: Gartner's application-versus-infrastructure split.

\paragraph{Seat-SaaS growth, $g_{\text{seat}}$ ($-6$ / 0 / $+5$ \%/yr; secondary).} The market has repriced the sector on the assumption of decline; Atlassian's seat decline is the first hard datum; vendors have an incentive to report consumption metrics in place of seat counts. Improve with: quarterly seat disclosures from the vendors that still give them. Likely: negative; the question is how negative.

\paragraph{Reallocation to AI (12 / 20 / 32 \%; secondary).} Pure judgment: how much of the seat and services money becomes AI-solution spend rather than savings returned to the CFO. The most improvable number in the paper with the least data. Improve with: any vendor's disclosure of AI-consumption revenue replacing seat revenue (Salesforce, ServiceNow, SAP); Gartner's 2027 category data. Likely: the mode, with the money arriving later than the seat losses.

\paragraph{Infrastructure software and services growth (4/8/12; $-2$/3/7 \%/yr).} Judgment; services could shrink faster if AI removes consulting hours quickly.

\paragraph{European pool displacement (10 / 20 / 35 \%).} As for $g_{\text{EU}}$; applied to the pool rather than to demand.

\paragraph{New AI spend (2 / 4 / 7 \%).} Judgment; small.

\paragraph{Advertising and consumer caps (140/170/200; 80/105/130).} From the decomposed ceilings; same sources and biases as above.

\subsection*{Incumbent case}

\paragraph{Microsoft 365 seats and price (420 / 450 / 480 M; \$19 / 21 / 23).} Widely quoted; not verified against a filing; Microsoft quotes seats selectively. The illustration hangs on these two numbers and they are flagged; a single disclosed seat series would settle it. Improve with: the 10-K segment revenue and the last disclosed seat count.

\paragraph{Price growth, $p$ (5 / 8 / 12 \%/yr).} From the announced July 2026 increases. Microsoft's incentive is to raise price to offset seats; it has done so.

\paragraph{Seat elasticity, $\epsilon$ (0.15 / 0.35 / 0.60; secondary).} Judgment; enterprise software has historically been sticky, agents and sovereignty exits make it less so. Improve with: renewal data, if any vendor ever discloses it.

\paragraph{Agent seat change, EU exit, attach.} Judgment ranges consistent with the sector chains; the seat trend is now anchored on Microsoft's disclosed +6--7\% commercial seat growth through FY26, from which the model assumes deceleration. Attach (12 / 22 / 35 \% by 2030, from $\sim$6.5\%) is the number that decides the illustration and the one a regular disclosure would settle.

\subsection*{Water and thermal}
\paragraph{WUE, evaporative share, $E_a$.} Engineering ranges from public sources (ASHRAE, JEDEC, operator sustainability reports --- which have an incentive to report improving WUE). Used in the text, not in the carry.

\subsection*{What we think is likely, in one paragraph}
The supply-side numbers will be settled by filings within two quarters and are unlikely to move the carry by more than 10\%. The demand-side base will come in near its mode once audited figures replace run-rates, and $\theta$ will drift down as the labs' margins improve. The sector ceilings will hold except for office assistants, where the answer depends on agent adoption that no one can see from here, and advertising, where Meta and Google will keep attributing growth to AI whether or not the attribution is clean. The budget pool will be smaller and more reallocated than the base case assumes, hurting seat-priced incumbents more than the aggregate. Siting friction will peak in 2027. The result --- a sunk deficit in 2027--28 under every path, a plateau that pays only with adoption far above what is visible, and a rational operator who cuts to \$400--500B and recovers by 2034 --- is more likely to be moved by a capex decision than by any demand surprise. That decision is the thing to watch on October 28.

\section{Update Protocol: Keeping the Estimate Live}\label{app:update}
Every number in this paper is a function of a register of inputs, each with a source, a source date, the date its source next publishes, and the model parameter it feeds. The register (\texttt{parameters.json}), the model, and a rebuild script are deposited with the paper. Updating is mechanical: change the input, run the rebuild, read the delta.

\subsection*{The register}
Forty-two entries in five groups --- supply (capex history, silicon share, lives, opex, WACC, lags, grid and siting), demand (the 2026 base and its components, cloud-capture share, sector ceilings, European displacement, common factors), feasibility (software and services pools, seat share, seat growth, reallocation), the incumbent worked case (seats, price, attach, elasticity), and the re-issue rule. Entries marked \emph{secondary} are swept in their own figures rather than trusted at a range. Two commitments follow from the register: nothing enters the model except through it, so the provenance trail is complete; and any reader can substitute their own value and rebuild.

\subsection*{The calendar}
\begin{table}[htbp]
\centering\small
\caption{Disclosure calendar and what each event moves}
\label{tab:calendar}
\begin{tabularx}{\textwidth}{L{3.2cm}L{4.2cm}X}
\toprule
\textbf{Date} & \textbf{Event} & \textbf{Register entries touched} \\
\midrule
Oct.~28--30, 2026 & Microsoft, Alphabet, Amazon, Meta Q3 calls & 2026 capex actuals and 2027 guidance (capex history, paths); M365 Copilot paid seats and any seat-count change (assistant ceiling, incumbent case); Advantage+ run-rate (advertising ceiling); Gemini MAU (consumer ceiling); AI run-rates (2026 base) \\
Oct.~2026 & Data Center Watch Q3 report & Siting deferral fraction \\
Oct.~2026 & Gartner Q4 IT-spending update & Software and services pools; AI-share signal (feasibility band) \\
Nov.~19, 2026 & Nvidia Q3 FY27 & Data-center revenue (silicon share; capex cross-check) \\
Monthly & Lab run-rate reports (Epoch AI revenue dataset) & 2026 base; growth versus $g^{*}$ \\
Dec.~2026 & Oracle Q2 FY27 & Big-5 capex; RPO \\
Jan.~2027 & Damodaran cost-of-capital update & WACC \\
Late Jan.~2027 & Big-4 Q4 calls and 10-Ks & Calendar-2026 revenue test (interim marker); useful-life notes (silicon lives); related-party notes (circular flows); capex 2027 guidance \\
Q1 2027 & OpenAI S-1 if it becomes public & Audited lab revenue, margin, compute commitments ($\theta_{\text{lab}}$, base) \\
Quarterly & Seat-metric disclosures (Atlassian, Salesforce, ServiceNow, Adobe) & Seat-SaaS growth (secondary variable) \\
Mar.~2027 & State of Digital Sovereignty 2027 survey & European displacement \\
Jan.~2028 & Calendar-2027 revenue & Interim marker: $\sim$\$330B is the bottom-up path, above \$350B favors the build-out \\
Jan.~2029 & Calendar-2028 revenue & The revenue test: above \$400B overturns the most likely outcome \\
\bottomrule
\end{tabularx}
\end{table}

\subsection*{The re-issue rule}
After each rebuild the script reports the change in three summary statistics against the previous run: the 2030 plateau $z$-score, the 2030 plateau break-even growth $g^{*}$, and P(profit, 2030) under rationalization. A change is \emph{minor} if $|\Delta z|<0.5$, $|\Delta g^{*}|<5$ points and $|\Delta\text{P}|<15$ points: tables are updated, the version number is incremented, and no announcement is made. A change is \emph{material} if any threshold is crossed, or if any registered 2026 observable arrives outside its P10--P90 range: the paper is re-issued with a delta log --- a table of every input that moved, its old and new range, its source, and the effect on the three statistics. Cadence is set by the calendar: a rebuild after each Big-5 earnings week and after each Nvidia report, and a material re-issue whenever the rule fires.

\subsection*{A dry run}
To test the loop before it is needed: suppose the October calls put 2026 capex at \$800B and calendar-2026 AI revenue at \$195B rather than \$179B. Rebuilding gives $g^{*}_{2030}=46\%$ (from 49\%), $z=-3.57$ (from $-3.54$) and P(profit, rationalize, 2030) two points lower: all three inside the minor thresholds. The higher base helps the demand side, the higher capex hurts the supply side, and they nearly cancel. A disclosure that Copilot paid seats had reached 40M, by contrast, would move the assistant ceiling's floor and be reported as material.

\subsection*{Commitment for this release}
This version (v26) is issued on 30 September 2026 as the pre-earnings estimate. A post-earnings update (v27) will be issued within two weeks of Nvidia's 19 November report, after the Big-5 Q3 calls of 28--30 October, whether or not the re-issue rule fires: it will carry a delta log of every register entry that moved, the rebuilt tables and figures, and a one-paragraph statement of what the calls did to the three summary statistics. Readers should regard v26's numbers as the stake in the ground that v27 measures against. Register changes made between 29 and 30 September (the OpenAI and Anthropic September disclosures, the batch share, the effective-price tier) rebuilt as minor under the rule: $z$ $+0.07$, $g^{*}_{2030}$ $-1$ point, P(profit, rationalize, 2030) $+2$ points.

\subsection*{Formats}
The paper is released as LaTeX source and PDF, deposited on Zenodo with the register, model and figure scripts; the PDF is the reading copy and the source is the object of record. Word and other office formats are produced only for reviewers' comments and are not released.

\subsection*{What is not automated}
Judgment ranges --- the common factors, the residual adoption factor, the reallocation share, elasticities --- do not update from a feed and are re-examined only when a disclosure bears on them. The re-issue rule guards against the two failure modes an updated model is prone to: over-reacting to one quarter, and quietly drifting through many small changes that add up to a different paper.

\section{Glossary}\label{app:glossary}
Terms are grouped by discipline; a table lists every model variable and parameter with its symbol, register name, units, base range and meaning; and a final section explains the phrases and idioms the paper relies on. Where a term is used with a paper-specific meaning, that meaning is given.

\subsection*{Paper-specific terms and model symbols}
\begin{description}\setlength{\itemsep}{1pt}
\item[Carry ($H_t$)] The annual cost of owning the installed base in year $t$: depreciation plus facility opex plus a return on net book value. The revenue the base must earn just to stand still.
\item[Cohort] Capital spent in one calendar year, tracked as a unit through its own depreciation life. A \emph{vintage model} sums the cohorts alive in each year.
\item[AI-attributable capex ($A_c$)] Big-5 capex in year $c$ minus a non-AI baseline (fulfillment, ordinary cloud refresh). Only this part enters the carry.
\item[Big-5] Microsoft, Alphabet, Amazon, Meta, Oracle --- the five whose capex Epoch AI tracks and whose infrastructure hosts most frontier AI.
\item[Silicon share ($\phi$)] Fraction of AI-attributable capex that is accelerators, memory, networking and servers, as opposed to shell, power and cooling.
\item[Shell] Buildings, power distribution, cooling plant --- the long-lived (15--25 yr) part of a data center.
\item[Energization] Connecting built capacity to grid power so it can run. The \emph{lag} ($\ell$) is the share of a cohort energized in its spend year; the rest follows a year later.
\item[Grid energization cap] Global capacity to connect new AI load each year, in GW/yr; capex above it is deferred.
\item[Siting / consent deferral ($f_{\text{delay}}$)] Share of a cohort delayed a year by permitting, litigation, moratoria or relocation.
\item[Site cost uplift ($\kappa_{\text{site}}$)] Multiplier on shell cost for community-benefit agreements, water mitigation, dry cooling, private generation, relocation.
\item[Sovereignty overlay ($\kappa_{\text{sov}}$, $\gamma_{\text{sov}}$)] Capital duplication multiplier and demand carve-out from data-residency and procurement rules.
\item[End-customer revenue ($D_t$; $R_s(t)$)] What buyers --- enterprises, consumers, advertisers --- pay in year $t$ for AI products and services, in \$B. $D_t$ is the total; $R_s(t)$ is one sector's part of it, and $D_t=\sum_s R_s(t)$. It is counted \emph{before} the frontier labs take their margin and before the part that never reaches a cloud (leakage) is removed. See Figure~\ref{fig:flow}.
\item[Infrastructure receipts (also ``receipts'', ``cloud receipts'')] The part of end-customer revenue that reaches the owners of the data centers, in \$B: $\text{Receipts}_t=\sum_s\theta_s R_s(t)$. This, not end-customer revenue, is what is compared with the carry.
\item[Cloud-capture share ($\theta_s$; $\theta$)] Dimensionless, 0--1. The fraction of sector $s$'s end-customer revenue that becomes infrastructure receipts, after the lab's margin, on-device execution, on-premise deployment and integrator capture. $\theta$ without a subscript is the revenue-weighted average; $\theta_{\text{lab}}$ is the same quantity for the frontier-lab channel alone, i.e.\ the labs' compute cost as a share of their revenue.
\item[Ceiling ($C_s$)] In \$B. A sector's maximum 2030 end-customer revenue, set by a physical or budgetary constraint (seats, IT budgets, ad yields, users). Given as a triangular range; the logistic never quite reaches it.
\item[Ramp ($k_s$)] In units of per year. The rate parameter of the logistic curve; a sector covers half the remaining distance to its ceiling in about $\ln2/k_s$ years (0.9 years at $k=0.8$). It is not a growth rate and not ``$1+k$''.
\item[Logistic] The S-shaped saturation curve $R(t)=C/[1+(C/R_0-1)e^{-k(t-t_0)}]$ with $t_0=2026$ and $R_0$ the 2026 value: growth that decays as the ceiling is approached.
\item[Break-even growth ($g^{*}$)] In percent per year (0.49 $=$ 49\%/yr). The constant annual growth rate of end-customer revenue from 2026 at which infrastructure receipts equal carry in a given year. Held constant by construction so that one number can be tested; the bottom-up growth is not constant (it decays).
\item[Residual adoption/capability factor] Dimensionless multiplier on 2030 receipts (1.0 $=$ the sector chains as estimated). It stands for what the chains cannot see from 2026: the pilots and seat counts visible today let one see roughly four years ahead, but not whether a capability step makes agents trustworthy with whole workflows. Because it cannot be estimated, it is swept as a what-if axis rather than assigned a range.
\item[Feasibility cutoff / band] In \$B. The most that could be spent on AI in 2030 out of the enterprise software and services pool --- itself decomposed into seat-priced software, infrastructure software, and services, each with its own growth, less the share displaced from US vendors --- plus advertising and consumer caps. The cutoff is uncertain, so it is drawn as a band from its P10 to its P90 (\$920B--\$1.31T); its P50 is $\sim$\$1.10T.
\item[Common factor] A multiplier drawn once per Monte Carlo trial and applied to every sector's ceiling (and a second to every ramp), representing the macroeconomic and technological background all sectors share --- recession, rates, regulation, or a capability step. Without it, fifteen independent sectors would average to a falsely narrow total.
\item[Plateau / Arms race / Rationalize / Adaptive] The forward capex paths: flat at the 2026 level of \$780B; the 2027 trajectory continued at \$1.0--1.1T (an \emph{arms race}: each operator building because the others are); a 35\% cut to \$500--550B; and an adaptive path in which capex tracks realized receipts (Appendix~\ref{app:horizon}).
\item[Top-down vs bottom-up] Two independent demand estimates: a base grown at $g$; and fifteen sector chains summed. They share only the 2026 base.
\item[Fermi chain] A quantity estimated as a product of several separately-bounded factors, in the manner of Fermi's piano-tuner estimate. In this paper each factor is bounded to tens of percent, not to an order of magnitude; the Fermi label refers to the method of decomposition and propagation, not to the precision.
\item[Slack trap] The paper's term for productivity gains absorbed as organizational slack (verification, output inflation, fractional time savings) rather than realized as cash flow.
\item[Last-mile operator] An external firm that delivers AI outcomes to an enterprise on gain-share pricing, typically on small open-weight models.
\end{description}

\subsection*{Model variables and parameters}
Every quantity the model uses, with its symbol as it appears in the text, its name in the register (\texttt{parameters.json}), its units, its base-case range (triangular min / mode / max unless stated) and what it means. Ranges are those of the current version; the register is authoritative.

\begin{center}\footnotesize
\begin{tabularx}{\textwidth}{L{1.7cm}L{2.6cm}L{1.3cm}L{2.2cm}X}
\toprule
\textbf{Symbol} & \textbf{Register name} & \textbf{Units} & \textbf{Base range} & \textbf{Meaning} \\
\midrule
\multicolumn{5}{l}{\textit{Supply side --- the installed base and its carry}} \\
$I_c$ & capex\_history & \$B/yr & 161, 260, 448, 780 (2023--26); paths after & Big-5 capital expenditure in cohort year $c$ (cash PP\&E plus new finance leases) \\
$A_{2023}$ & ai\_capex\_2023 & \$B & 20 / 30 / 40 & AI purchases in 2023; with the flat 2023 total, pins the non-AI baseline $N_{2023}=161-A_{2023}$ \\
$g_N$ & nonai\_growth & 1/yr & 0.02 / 0.06 / 0.10 & Growth of the non-AI baseline (fulfillment, ordinary cloud refresh) \\
$N_c$ & --- (derived) & \$B/yr & $\sim$156 in 2026 & Non-AI baseline in year $c$: $N_{2023}(1+g_N)^{c-2023}$ \\
$A_c$ & --- (derived) & \$B/yr & --- & AI-attributable capex: $\max(I_c-N_c,0)$ \\
$\phi$ & silicon\_share & 0--1 & 0.45 / 0.52 / 0.60 & Share of $A_c$ that is accelerators, memory, networking, servers (rest is shell) \\
$f_{\text{tr}}$ & training\_share & 0--1 & 0.35 / 0.45 / 0.55 & Share of the silicon fleet in training clusters (rest inference) \\
$L_{\text{tr}}$ & life\_training & years & 2.5 / 3.5 / 4.5 & Economic life of training-cluster silicon (power-slot obsolescence) \\
$L_{\text{inf}}$ & life\_inference & years & 4 / 5 / 6 & Economic life of inference-fleet silicon \\
$L_{\text{si}}$ & --- (derived) & years & $\sim$4.2 & Effective silicon life, harmonic mean: $1/(f_{\text{tr}}/L_{\text{tr}}+(1-f_{\text{tr}})/L_{\text{inf}})$ \\
$L_{\text{sh}}$ & life\_shell & years & 15 / 20 / 25 & Life of buildings, power distribution, cooling plant \\
$\ell_{\text{si}}$, $\ell_{\text{sh}}$ & lag\_silicon, lag\_shell & 0--1 & 0.5/0.7/0.9; 0.2/0.4/0.6 & Fraction of a cohort's silicon (shell) energized in its spend year; remainder the year after \\
$c_{\text{opex}}$ & opex\_share & 1/yr & 0.059 / 0.0665 / 0.074 & Facility opex per year as a fraction of in-service gross stock (power, staff, tax, network, other) \\
$u$ & cooling\_opex\_uplift & 0--1 & 0 / 0.03 / 0.06 & Opex uplift where evaporative cooling is refused (dry cooling) \\
$r$ & wacc & 1/yr & 0.090 / 0.0945 / 0.099 & Weighted average cost of capital (CAPM); return charged on net book value \\
$G_{c,j}$ & --- (derived) & \$B & --- & In-service gross of cohort $c$, component $j\in\{\text{si},\text{sh}\}$ \\
$H_t$ & --- (output) & \$B/yr & $\sim$725 in 2030 (plateau) & Carry: depreciation $+$ opex $+$ return on net book $+$ return on deferred capital, Eq.~\ref{eq:Ht} \\
grid cap & grid\_cap & \$B/yr & 1150 / 1400 / 1680 & Energizable AI capex per year (22--40 GW at \$38--55M/MW); excess deferred a year \\
$f_{\text{delay}}$ & siting\_deferral & 0--1 & 0.10 / 0.20 / 0.35 & Share of each cohort from 2026 deferred a year by consent, siting, water \\
$\kappa_{\text{site}}$ & site\_cost\_uplift & multiplier & 1.02 / 1.05 / 1.09 & Shell cost uplift for community benefits, water mitigation, relocation \\
$\kappa_{\text{sov}}$ & (overlay) & multiplier & 1.1 / 1.2 / 1.3 & Capital duplication from data-residency rules; sovereignty scenario only \\
\bottomrule
\end{tabularx}
\vspace{4pt}
\begin{tabularx}{\textwidth}{L{1.7cm}L{2.6cm}L{1.3cm}L{2.2cm}X}
\toprule
\textbf{Symbol} & \textbf{Register name} & \textbf{Units} & \textbf{Base range} & \textbf{Meaning} \\
\midrule
\multicolumn{5}{l}{\textit{Demand side --- end-customer revenue and infrastructure receipts}} \\
$D_{2026}$, $D_{t_0}$ & base\_2026 & \$B & 163 / 179 / 194 & Calendar-2026 end-customer AI revenue, all channels (seven components) \\
$g$ & (swept) & 1/yr & 0.2--0.8 & Sustained growth of end-customer revenue in the top-down sweep \\
$g^{*}$ & --- (output) & 1/yr & 0.49 (plateau, 2030) & Break-even sustained growth: receipts $=$ carry in the target year \\
$\theta$ & theta & 0--1 & 0.77 / 0.815 / 0.86 & Revenue-weighted cloud-capture share, top-down \\
$\theta_s$ & sectors (table) & 0--1 & per sector & Cloud-capture share of sector $s$: fraction of its end-customer revenue reaching infrastructure owners \\
$\theta_{\text{lab}}$ & --- (component of $\theta$) & 0--1 & $\sim$0.635 & Labs' compute cost as a share of their revenue, by lab \\
$R_{s,0}$, $R_s(t)$ & sectors (table) & \$B & per sector & Sector $s$'s 2026 revenue; its revenue in year $t$ on the logistic \\
$C_s$ & sectors (table) & \$B & per sector & Sector $s$'s 2030 ceiling, from a physical or budgetary constraint \\
$k_s$ & sectors (table) & 1/yr & per sector & Sector $s$'s logistic ramp rate; half the remaining distance in $\ln2/k_s$ years \\
$D^{\text{BU}}_t$ & --- (output) & \$B/yr & $\sim$626 in 2030 & Bottom-up end-customer revenue, $\sum_s R_s(t)$ \\
Receipts$_t$ & --- (output) & \$B/yr & $\sim$479 in 2030 & Infrastructure receipts, $\sum_s\theta_s R_s(t)$ \\
$M_C$, $M_k$ & common\_factor\_ceiling, \_ramp & multiplier & 0.7/1.0/1.4; 0.7/1.0/1.3 & Common factors on all ceilings (ramps), drawn once per trial \\
$g_C$ & post2030\_ceiling\_growth & 1/yr & 0.06 / 0.10 / 0.14 & Growth of sector ceilings after 2030 (Appendix~\ref{app:horizon}) \\
$g_{\text{EU}}$ & eu\_displacement & 0--1 by 2030 & 0.03 / 0.05 / 0.08 & European displacement of US vendors, ramped 2027--30 on sovereign-exposed sectors \\
$\gamma_{\text{sov}}$ & (overlay) & 0--1 & 0.15 / 0.25 / 0.35 & Demand carve-out for domestic champions and enclaves; sovereignty scenario only \\
$F_{\text{circ}}$ & circular & \$B/yr & 8 / 12 / 18 & Round-trip revenue (equity returning as cloud spend), subtracted from receipts \\
$\Delta_t$ & --- (output) & \$B/yr & $\sim-259$ in 2030 & Receipts$_t-F_{\text{circ}}-\kappa H_t$: annual surplus (deficit) \\
$z$ & --- (output) & --- & $\sim-3.5$ & $\text{median}(\Delta_{2030})/\sigma_{\Delta}$ \\
residual factor & (swept) & multiplier & 0.6--1.7 & Multiplier on 2030 receipts for adoption the sector chains cannot see (secondary variable) \\
\bottomrule
\end{tabularx}
\vspace{4pt}
\begin{tabularx}{\textwidth}{L{1.7cm}L{2.6cm}L{1.3cm}L{2.2cm}X}
\toprule
\textbf{Symbol} & \textbf{Register name} & \textbf{Units} & \textbf{Base range} & \textbf{Meaning} \\
\midrule
\multicolumn{5}{l}{\textit{Feasibility band --- what enterprise, advertising and consumer budgets can finance}} \\
SW$_{26}$, SV$_{26}$ & software\_2026, services\_2026 & \$B & 1380/1440/1500; 1800/1870/1940 & Gartner 2026 enterprise software and IT services pools \\
seat share & seat\_share\_of\_software & 0--1 & 0.50 / 0.55 / 0.60 & Share of software that is seat-priced application SaaS \\
$g_{\text{seat}}$ & seat\_saas\_growth & 1/yr & $-0.06$ / 0 / 0.05 & Growth of seat-priced software 2026--30 (secondary variable) \\
$g_{\text{infra}}$, $g_{\text{sv}}$ & (fixed ranges) & 1/yr & 0.04/0.08/0.12; $-0.02$/0.03/0.07 & Growth of infrastructure software; of IT services \\
disp & eu\_pool\_displacement & 0--1 & 0.10 / 0.20 / 0.35 & Share of Europe's pool ($\sim$23\% of global) moved away from US vendors by 2030 \\
realloc & reallocation\_to\_ai & 0--1 & 0.12 / 0.20 / 0.32 & Share of the accessible pool reallocated from seats and services to AI solutions (secondary variable) \\
ainew & (fixed range) & 0--1 & 0.02 / 0.04 / 0.07 & Genuinely new AI spend as a share of the pool \\
ads cap, cons cap & (fixed ranges) & \$B & 140/170/200; 80/105/130 & Advertising and consumer ceilings entering the band \\
cutoff & --- (output) & \$B & 920--1310 (P10--P90) & Total 2030 AI revenue financeable: pool $\times$ (realloc $+$ ainew) $+$ ads $+$ cons \\
\bottomrule
\end{tabularx}
\vspace{4pt}
\begin{tabularx}{\textwidth}{L{1.7cm}L{2.6cm}L{1.3cm}L{2.2cm}X}
\toprule
\textbf{Symbol} & \textbf{Register name} & \textbf{Units} & \textbf{Base range} & \textbf{Meaning} \\
\midrule
\multicolumn{5}{l}{\textit{Incumbent worked case (Microsoft 365 commercial)}} \\
seats$_{26}$, ARPU$_{26}$ & seats\_2026\_M, arpu\_2026 & M; \$/seat/mo & 420/450/480; 19/21/23 & Paid commercial seats and blended price, 2026\verify \\
$p$ & price\_growth & 1/yr & 0.05 / 0.08 / 0.12 & List-price growth including the July-2026 increase \\
$\epsilon$ & seat\_elasticity & --- & 0.15 / 0.35 / 0.60 & Seats lost per unit of cumulative price increase (secondary variable) \\
agent seat change & (fixed range) & 1/yr & $-0.06$ / $-0.03$ / 0 & Seat change from agents replacing human seats \\
EU exit & (fixed range) & 0--1 by 2030 & 0.03 / 0.08 / 0.15 & Cumulative sovereignty seat loss \\
attach & copilot\_paid\_seats\_M (base) & 0--1 & 0.12 / 0.22 / 0.35 by 2030 & Paid Copilot attach on remaining seats ($\sim$6.5\% today) \\
\bottomrule
\end{tabularx}
\vspace{4pt}
\begin{tabularx}{\textwidth}{L{1.7cm}L{2.6cm}L{1.3cm}L{2.2cm}X}
\toprule
\textbf{Symbol} & \textbf{Register name} & \textbf{Units} & \textbf{Base range} & \textbf{Meaning} \\
\midrule
\multicolumn{5}{l}{\textit{Water and thermal (used in text, not in the carry)}} \\
WUE & --- & L/kWh & 0.3 / 1.0 / 1.8 & Water consumed per kWh of IT energy \\
evaporative share & --- & 0--1 & 0.4 / 0.6 / 0.8 & Share of the fleet on evaporative cooling \\
$E_a$, $k$, $AF$ & --- & eV; eV/K; --- & $E_a\approx0.7$ & Activation energy, Boltzmann constant, Arrhenius acceleration factor \\
\bottomrule
\end{tabularx}
\vspace{4pt}
\end{center}

\subsection*{Symbols used in the derivations (not model inputs)}
\begin{description}\setlength{\itemsep}{1pt}
\item[$Y$, $X_i$, $n$] In the method section: a Fermi chain $Y=\prod_{i=1}^{n}X_i$ of $n$ factors $X_i$.
\item[$\sigma_{\ln X}$, $\sigma_{\ln Y}$, $\sigma_{\ln D}$, $\sigma_{\ln R_s}$] Relative (log) standard deviations of a factor, a chain, the demand total and a sector; see \emph{$\sigma_{\ln}$} above.
\item[$w_s$] Sector $s$'s share of total end-customer revenue, $R_s/D$; the weight (squared) in the sum rule of Appendix~\ref{app:errors}.
\item[$\ell$ vs $L$] Lower-case $\ell$ (with subscript si or sh) is an energization lag, a fraction between 0 and 1; upper-case $L$ is a life in years. Different letters for different quantities.
\item[$\omega(e)$] Dollars flowing along edge $e$ of the circular-flow graph; the revenue a cycle can round-trip is the smallest $\omega$ on it.
\item[$j$, $L_j$, $G_{c,j}$] Component index in the carry equation: $j=\text{si}$ (silicon) or $\text{sh}$ (shell); the life and in-service gross of that component of cohort $c$.
\item[{$\mathbf{1}[\cdot]$}] Indicator: 1 if the bracketed condition holds (the cohort is still alive), 0 otherwise.
\item[$\kappa$ (unsubscripted)] In $\Delta_t=\text{Receipts}_t-F_{\text{circ}}-\kappa H_t$: the sovereignty multiplier applied to carry, 1 in the base case.
\item[$\tau$] Nominal task-level time saving from AI as a fraction of payroll (Triang 0.10 / 0.15 / 0.25), Section~\ref{sec:slack}.
\item[$\sigma_{\text{slack}}$] Share of nominal saving absorbed as organizational slack ($\approx$0.65): fractional time, output inflation, the manager's dilemma.
\item[$\nu_{\text{verify}}$] Share of the remainder consumed by human verification and liability review ($\approx$0.20).
\item[$W$] The payroll pool the saving applies to: $\sim$\$30T global knowledge-worker compensation, $\sim$\$9T US.
\item[$B_{\text{realized}}$, $C_{\text{severance}}$] Cash benefit a firm can actually book, $W\tau(1-\sigma_{\text{slack}})(1-\nu_{\text{verify}})-C_{\text{severance}}$; one-time separation costs.
\item[$B_{\text{req}}$, $S_{\text{ent}}$, $H_{\text{ent}}$] Benefit enterprises must realize ($\sim$\$1.2T in 2028); their total AI outlay including integration ($\sim$\$1.0T); the enterprise share of the carry ($\sim$\$400B).
\item[$\mu_{\text{integ}}$, $m$] Integration spend per dollar of software and tokens (Triang 1.0 / 1.5 / 2.5); the required benefit margin over cost ($\approx$0.2).
\item[$AF$, $E_a$, $k$, $T_1$, $T_2$] Arrhenius acceleration factor between absolute temperatures $T_1$ and $T_2$; activation energy ($\approx$0.7 eV); Boltzmann's constant ($8.617\times10^{-5}$ eV/K). Note $k$ here is not the logistic ramp $k_s$.
\item[$J$, $n$ (Black's equation)] Current density and its exponent (typically $\sim$2) in electromigration lifetime, MTTF $\propto J^{-n}e^{E_a/kT}$.
\item[$\epsilon$, $p$ (incumbent case)] Seat-loss elasticity and annual price growth; revenue ratio $(1+p)^4[1-\epsilon((1+p)^4-1)]$ times the seat trend.
\item[$g$ subscripts] $g_N$ baseline growth; $g_C$ post-2030 ceiling growth; $g_{\text{EU}}$ European displacement; $g_{\text{seat}}$, $g_{\text{infra}}$, $g_{\text{sv}}$ growth of the three budget pools; $g^{*}$ break-even growth; bare $g$ the top-down sweep. All per year.
\item[$t_0$, $t$, $c$] The base year (2026); the year being evaluated; a cohort (capex) year.
\item[P10, P50, P90] Percentiles: the value below which 10\%, 50\% (median), 90\% of Monte Carlo outcomes fall.
\end{description}

\subsection*{Statistics and estimation}
\begin{description}\setlength{\itemsep}{1pt}
\item[Triang($a$, $b$, $c$)] A triangular probability distribution with minimum $a$, most likely value $b$, maximum $c$: zero probability outside $[a,c]$, peaking at $b$. Used for every input because it needs only three numbers an expert can defend. Its relative uncertainty, $\sigma_{\ln}$, is approximately $\ln(c/a)/4.9$: the full log-range divided by about 4.9, because a triangle's standard deviation is roughly a fifth of its base (exactly $1/\sqrt{24}$ of the base for a symmetric triangle).
\item[P10 / P90] The 10th and 90th percentiles of a distribution: the range inside which 80\% of Monte Carlo outcomes fall.
\item[Monte Carlo (MC)] Repeating a calculation $N=100{,}000$ times with inputs drawn from their distributions, to obtain the distribution of the result.
\item[$\sigma_{\ln}$] Standard deviation of the natural log of a quantity; approximately its relative (percentage) uncertainty.
\item[Quadrature / root-sum-square (RSS)] Combining independent uncertainties as $\sqrt{\sum\sigma_i^2}$ rather than $\sum\sigma_i$. Two independent errors of 10 combine to 14, not 20. In the variance budget the same rule is applied to the half-swings, which play the role of the $\sigma_i$.
\item[Half-swing] Take one input; set every other input to its mode; move the chosen input from its minimum to its maximum and record how much the result changes. Half of that change is the half-swing. Min-to-max (not mode-to-max) is used because the triangular ranges are not symmetric, and halving makes it comparable to a standard deviation. Summed in quadrature, the half-swings reproduce the Monte Carlo spread (Appendix~\ref{app:errors}).
\item[Variance budget] Apportionment of total uncertainty among inputs by squared half-swing; shows which inputs matter.
\item[$z$-score] The median gap between receipts and carry ($\Delta$), divided by the standard deviation of that gap: $z=\Delta/\sigma_\Delta$. A $z$ of $-3.5$ means the median deficit is 3.5 standard deviations below zero. $|z|\ge2$ is the paper's threshold for a robust sign; it is a statement about the modeled distribution, not a guarantee.
\item[Tornado chart] Horizontal bars of half-swings, largest at top, showing sensitivity to each input.
\item[Harmonic mean] The right average for asset lives, because depreciation is $1/L$ per year: a fleet that is 45\% training silicon ($f_{\text{tr}}$) at 3.5 years and 55\% inference silicon at 5 years depreciates at $0.45/3.5+0.55/5=0.239$ per year, i.e.\ an effective life of $1/0.239=4.2$ years --- not the arithmetic 4.3. The shorter-lived part dominates.
\item[CAGR] Compound annual growth rate.
\item[Run-rate vs booked] Run-rate (ARR) is the latest month or quarter annualized; booked is what was recognized over the calendar year. In a fast-growing business the two differ: a company at \$20B run-rate in January and \$40B in December booked roughly the average, \$30B, if growth was linear --- less if it accelerated late. Reported run-rates are converted to calendar bookings this way throughout.
\item[Gross vs net basis] Whether money a reseller collects and forwards to a supplier is counted as the reseller's revenue (gross) or only the reseller's margin (net). It applies to any reseller relationship; the case that matters here is that Anthropic books end-customer spend on Amazon Bedrock and Google Vertex as its own gross revenue, so the same dollar can appear in both a lab's and a hyperscaler's AI figures. The paper subtracts an overlap for this.
\end{description}

\subsection*{Finance and accounting}
\begin{description}\setlength{\itemsep}{1pt}
\item[Capex] Capital expenditure: cash spent on long-lived assets, here including new finance leases.
\item[Opex] Operating expenditure. Here \emph{facility opex}: power, cooling, on-site staff and maintenance, property tax, insurance, network --- everything it costs to run the building each year other than depreciation and financing.
\item[Depreciation] The annual write-down of an asset's cost over its useful life. \emph{Accounting life} is the company's book assumption (5--6 yr for servers); \emph{economic life} is how long the asset actually earns its keep.
\item[Net book value] Cost less accumulated depreciation; the base on which a return is charged.
\item[WACC] Weighted average cost of capital: the blended annual return a company must earn on invested capital to satisfy its shareholders and lenders. The paper's $r$. It is a property of the company, not of the project, so the same $r$ applies to AI and non-AI assets; the extra risk of AI is handled through the hurdle rate on the customer side, not the WACC.
\item[CAPM] Capital asset pricing model. Cost of equity $=r_f+\beta\times\text{ERP}$: the risk-free rate plus the stock's beta times the equity risk premium. With $r_f=4.3\%$, $\beta=1.2$ and ERP $=5\%$, cost of equity is $4.3+1.2\times5=10.3\%$; blended with cheaper debt it gives the WACC.
\item[Beta] A stock's sensitivity to the market; $>$1 means more volatile than the market.
\item[Equity risk premium] The excess return investors demand over the risk-free rate for holding equities. (Not to be confused with ERP, enterprise resource planning software, which also appears.)
\item[IRR] Internal rate of return. Future cash is worth less than present cash; a \emph{discount rate} $d$ converts a cash flow $F_t$ in year $t$ to its present value $F_t/(1+d)^t$. The IRR is the $d$ at which a project's discounted inflows exactly repay its outlay, $\sum_t F_t/(1+d)^t=0$: the annual return the project earns on the money tied up in it.
\item[Hurdle rate] The minimum IRR a company requires before approving an investment. Every business applies one to all spending; it is set above the WACC, and higher again for projects whose payback is uncertain --- so AI projects face a higher hurdle than routine capital, not a lower one.
\item[ROI / ROIC] Return on investment / on invested capital.
\item[NOPAT] Net operating profit after tax: operating profit (EBIT) less tax at the statutory rate; the numerator of Goldman Sachs's ROIC definition in Section~\ref{sec:gs}.
\item[EBIT] Earnings before interest and tax: operating profit.
\item[TCO] Total cost of ownership over an asset's life.
\item[COGS] Cost of goods sold; direct costs of delivering a product.
\item[Gross margin] Revenue less COGS, as a share of revenue.
\item[TAM] Total addressable market.
\item[ARPU] Average revenue per user (or per seat), in \$ per month or per year; for Microsoft 365 commercial, the blended price across E3, E5 and Business plans ($\sim$\$21/seat/month, unverified).
\item[M365] Microsoft 365, the commercial Office suite subscription; the largest seat-priced software base.
\item[DINUM] France's interministerial digital directorate, which mandated the government Windows-to-Linux migration in April 2026.
\item[IDC] International Data Corporation, the other large IT-market forecaster; a cross-check on Gartner.
\item[MIT NANDA] The MIT initiative whose 2025 report found $\sim$95\% of enterprise generative-AI pilots without measurable P\&L impact.
\item[DOI] Digital object identifier: the permanent citable identifier Zenodo issues for a deposit.
\item[Apple, Qualcomm, AMD, Intel (edge silicon)] Makers of the NPUs in phones and PCs (Apple Silicon, Snapdragon X, Ryzen AI, Core Ultra) that run routine inference locally.
\item[SK Hynix, Samsung, Micron] The three HBM suppliers; the bullwhip risk sits with their 2027--28 capacity.
\item[JPMorgan; Bank of America] Investment banks whose 2026 software-sector notes are cited on the SaaS repricing.
\item[MAU] Monthly active users: anyone who used the product at least once in the month. The distribution of use behind that count is extremely skewed --- a small fraction of heavy daily users, a long tail of occasional ones --- which is why the consumer ceiling multiplies MAU by a paid-conversion rate (4--8\%) rather than treating users as alike.
\item[RPO] Remaining performance obligations: contracted revenue not yet delivered (a backlog).
\item[DSCR] Debt-service coverage ratio: cash flow available divided by debt payments due; lenders set tripwires on it.
\item[SPV] Special-purpose vehicle: a separate legal entity used to hold assets and debt off a parent's balance sheet. SPVs change who carries the carry and where it appears, not its size; they matter to the paper only because their debt-service tripwires can turn a slow revenue shortfall into a fast liquidation.
\item[VC] Venture capital.
\item[ASC 360-10] The US accounting standard (FASB Codification) governing impairment of long-lived assets: if projected undiscounted cash flows fall below carrying value, the asset is written down to fair value.
\item[FASB] Financial Accounting Standards Board, the US accounting standard-setter.
\item[SEC; 10-K, 10-Q, 8-K, S-1] US Securities and Exchange Commission; its annual report, quarterly report, current-event report, and registration statement for an IPO.
\item[XBRL] The machine-readable tagging of SEC filings, from which Epoch AI's capex series is built.
\item[PP\&E] Property, plant and equipment.
\item[ROU asset] Right-of-use asset: a leased asset recognized on the balance sheet.
\item[YTD; FY; Q1--Q4; YoY; QoQ] Year to date; fiscal year; quarters; year-on-year; quarter-on-quarter. Nvidia's and Microsoft's fiscal years do not align with the calendar.
\item[ARR] Annual recurring revenue (see run-rate).
\end{description}

\subsection*{Computing and hardware}
\begin{description}\setlength{\itemsep}{1pt}
\item[GPU] Graphics processing unit; the general term for the accelerators (Nvidia H100, B200) that train and run AI models.
\item[TPU; Trainium; MTIA] Google's, Amazon's and Meta's in-house AI accelerators.
\item[NPU] Neural processing unit: an AI accelerator built into consumer devices (phones, PCs).
\item[TOPS] Tera-operations per second; the throughput rating of an NPU.
\item[HBM (HBM3, HBM3e, HBM4)] High-bandwidth memory: DRAM stacked vertically next to the accelerator die; successive generations.
\item[TSV] Through-silicon via: the vertical copper interconnect through stacked memory dies.
\item[FP4 / FP8 / FP16] Floating-point number formats of 4, 8 and 16 bits; lower precision means more throughput per watt.
\item[Microscaling (MX formats)] A block floating-point scheme, standardized by the Open Compute Project in 2023, in which a group of 32 values shares one scale factor so that each value needs only 4 bits. It is what makes FP4 accurate enough for inference; expected to spread fastest where accelerators are scarce, including China.
\item[Ternary weights] Model weights restricted to $\{-1,0,1\}$, eliminating multiplications.
\item[LLM; SLM] Large / small language model.
\item[GPT-4] OpenAI's March 2023 model; used here as the date from which AI capex is measured.
\item[Open-weight model] A model whose parameters are published for anyone to run (Llama, Mistral, DeepSeek).
\item[Distilled / quantized model] A smaller model trained to mimic a larger one; a model stored at lower numeric precision. Both cut compute cost.
\item[Token] The unit of text a model reads or writes (roughly 3/4 of a word); the unit in which inference is priced.
\item[Inference vs training] Running a trained model to answer requests, versus the compute-intensive process of creating it.
\item[RAG] Retrieval-augmented generation: a model answering with reference to documents fetched at query time.
\item[Agentic; agents] AI systems that execute multi-step tasks and take actions, rather than single responses.
\item[Office assistant / agent] The paper's generic term for AI assistants embedded in productivity suites and enterprise workflows (Microsoft 365 Copilot, Google Gemini for Workspace, Claude and ChatGPT enterprise seats, and the agentic workloads built on them). Lower-case ``copilot'' has become industry shorthand for this category; the paper avoids it to keep the category distinct from the product.
\item[Copilot (Microsoft 365 Copilot; GitHub Copilot)] Microsoft's branded assistants for Office and for coding. Their disclosed paid seat counts anchor the office-assistant sector's 2026 base.
\item[API] Application programming interface; here, metered access to a model over the network.
\item[Bedrock; Vertex; Azure AI Foundry] The AI platform services of AWS, Google Cloud and Microsoft Azure through which third-party models are resold.
\item[AWS; GCP; OCI; Azure] Amazon Web Services, Google Cloud Platform, Oracle Cloud Infrastructure, Microsoft's cloud.
\item[Neo-cloud] A GPU-rental specialist without a broad software business (CoreWeave and peers).
\item[Hyperscaler] A cloud operator building at multi-gigawatt scale (the Big-5 and their global peers).
\item[Frontier lab] A developer of leading models (OpenAI, Anthropic, xAI, Mistral, DeepSeek).
\item[CPU] Central processing unit; conventional server silicon.
\item[MLOps] Machine-learning operations: the engineering discipline of deploying and maintaining models in production.
\item[Edge] Computing on the user's own device --- a phone or PC with an NPU. \emph{Not} the same as on-premise.
\item[On-premise] Computing on servers the customer owns and operates in its own facilities, typically running open-weight models. Edge and on-premise both bypass the cloud meter; the paper's $\theta_s$ nets out both.
\item[InfiniBand; RoCE] High-speed cluster networking fabrics.
\item[A100; H100; B200] Successive Nvidia data-center accelerator generations (2020, 2022, 2024--25).
\end{description}

\subsection*{Data centers, power and reliability}
\begin{description}\setlength{\itemsep}{1pt}
\item[MW; GW; TWh; kWh] Megawatt, gigawatt (power); terawatt-hour, kilowatt-hour (energy). Data-center scale is quoted as IT load in MW or GW.
\item[PUE] Power usage effectiveness: total facility power divided by IT power; 1.2 means 20\% overhead for cooling and distribution.
\item[WUE] Water usage effectiveness: litres of water consumed per kWh of IT energy.
\item[Evaporative vs dry cooling] Cooling by evaporating water (cheap, water-intensive) versus by air-cooled heat exchangers (no water, more power).
\item[Interconnection queue] The waiting list to connect a new load or generator to the transmission grid. Its length differs by geography: 5--7 years in PJM and similar US regions, comparable or longer in Ireland, the Netherlands and Germany, and effectively a planning allocation rather than a queue in China.
\item[PJM; ERCOT; CAISO] The grid operators for the mid-Atlantic, Texas and California.
\item[EirGrid; CRU; TenneT; Liander] Ireland's transmission operator and regulator; the Dutch transmission and distribution operators.
\item[FLAPD] Frankfurt, London, Amsterdam, Paris, Dublin: Europe's five main data-center hubs.
\item[EIA] US Energy Information Administration.
\item[ETS] The EU Emissions Trading System, which prices carbon at $\sim$EUR~70/t. The US has no federal carbon price; California's cap-and-trade and the northeastern RGGI cover a minority of data-center load.
\item[Moratorium] A formal halt on new connections or approvals.
\item[Arrhenius relation] Failure rate rising exponentially with temperature: $AF=\exp[(E_a/k)(1/T_1-1/T_2)]$. With activation energy $E_a\approx0.7$\,eV this gives roughly a doubling per 10$^{\circ}$C.
\item[Activation energy ($E_a$); Boltzmann constant ($k$)] The energy barrier of a failure mechanism; the constant relating temperature to energy.
\item[Black's equation] Electromigration lifetime $\propto J^{-n}e^{E_a/kT}$: metal lines fail faster at high current density and temperature.
\item[Coffin--Manson] Fatigue life under thermal cycling; governs solder and micro-bump failure in stacked memory.
\item[MTTF] Mean time to failure.
\item[JEDEC] The semiconductor standards body; JEP122 is its failure-mechanism reference.
\item[ASHRAE TC 9.9] The technical committee that sets recommended data-center temperature and humidity envelopes.
\item[Salvage / second life] Reusing retired accelerators for less demanding work; the paper argues the power slot, not the silicon, is the binding constraint.
\end{description}

\subsection*{Economics, policy and institutions}
\begin{description}\setlength{\itemsep}{1pt}
\item[TFP] Total factor productivity: output growth not explained by more labor or capital.
\item[Productivity J-curve] The pattern (Brynjolfsson et al.) in which a general-purpose technology first depresses measured productivity, because of intangible complementary investment, then raises it.
\item[Solow paradox] ``You can see the computer age everywhere but in the productivity statistics.''
\item[Jevons paradox; Parkinson's law] Efficiency gains increase total consumption; work expands to fill the time available. Combined here as \emph{Parkinson--Jevons}: cheaper output inflates the demand for output.
\item[Surplus capture (Nordhaus)] Innovators historically kept about 2\% of the social surplus from their innovations.
\item[Gain-share] Pricing as a share of the customer's realized savings.
\item[Seat] A per-user software license.
\item[IT budget / IT spend] Enterprise spending on software, services, devices and data-center systems, as tracked by Gartner.
\item[GDPR; EU AI Act; BSI C5; DPDP] EU data-protection regulation and its Chapter V transfer rules; the EU's 2024 AI regulation; Germany's cloud security catalogue; India's 2023 Digital Personal Data Protection Act.
\item[Digital sovereignty] Policies requiring data and compute to remain under national jurisdiction.
\item[DPI] Digital public infrastructure (India's state-run digital platforms).
\item[East Data, West Computing] China's program to site compute in its western, low-cost-power provinces.
\item[GCC; G20] Gulf Cooperation Council states; the Group of Twenty economies.
\item[NHTSA; HIPAA; FINRA; SEC; SR 11-7] US vehicle-safety regulator; health-privacy law; securities self-regulator; securities regulator; the Federal Reserve's model-risk-management guidance.
\item[OECD; IFR; ILO; BLS; NEA] The OECD; International Federation of Robotics; International Labour Organization; US Bureau of Labor Statistics; China's National Energy Administration.
\item[Gartner] The largest IT-market research and forecasting firm; source of the IT-spending and AI-share figures.
\item[Epoch AI] A nonprofit research group that maintains open datasets on AI compute, spending and revenue; its capex series, built from SEC XBRL filings, is the paper's capex history, and its revenue-reports dataset anchors the lab figures.
\item[Data Center Watch] A research firm that tracks data-center projects blocked or delayed by local opposition; source of the siting-deferral calibration.
\item[CFO; CX; POC] Chief financial officer; customer experience (contact centers); proof of concept.
\item[Knowledge worker] Office, professional, technical and managerial labor; the paper's $\sim$1.05B global and 62M US pools.
\item[BLS OES / ECEC] The Bureau of Labor Statistics' occupational employment and wage survey, and its employer-cost-of-compensation series.
\end{description}

\subsection*{Phrases and idioms used in the paper}
\begin{description}\setlength{\itemsep}{1pt}
\item[The build-out] The 2023--2030 wave of AI data-center construction by the Big-5 and their peers.
\item[Earn its carry] For the installed base to generate receipts at least equal to its annual carry (depreciation, opex, return). The paper's test of solvency.
\item[The scissors] Figure~\ref{fig:scissors}: carry rising on one blade, receipts on the other; where and whether they cross.
\item[Sunk deficit] The 2027--28 shortfall that no growth rate can close, because the cohorts causing it are already on the books.
\item[Bull case / bear case] The optimistic and pessimistic readings of the same evidence; the paper tries to state each in its strongest form.
\item[Referee] An anonymous academic reviewer; ``what a referee would attack first'' means the weakest point a hostile expert reader would find.
\item[Coin-toss neighborhood] The region of parameter space where P(profit) is near 50\%: the sign of the result is not determined.
\item[Fermi-robust] A conclusion whose gap exceeds the propagated uncertainty ($|z|\ge2$), so it survives ordinary disputes about individual inputs.
\item[One-at-a-time] Moving a single input across its range with all others held at their modes; the basis of the tornado charts and variance budget.
\item[Move the estimate vs tighten it] Decomposition can shift the central value (a move) as well as narrow its range (a tightening); the paper records both.
\item[Toward the bull / toward the bear] The direction a revision moved the result: in favor of the build-out paying, or against.
\item[Hard cap / cutoff] A ceiling set by something that cannot be exceeded (a seat count, an IT budget share, physics), as opposed to a soft estimate.
\item[Round trip; circular flow; closed loop] Revenue that a company's own investment returns to it, e.g.\ equity in a lab that comes back as cloud spend.
\item[Pass-through; gross basis] Revenue a reseller collects and forwards to the supplier; counting it as the supplier's own revenue (gross) creates double counts across companies.
\item[Investor-funded compute] Cloud purchases by frontier labs paid for with venture and strategic capital rather than customer revenue; real receipts to hyperscalers, but not self-sustaining.
\item[Bridging] Temporary funding that covers a gap until revenue arrives; here, the investor-funded share of lab compute.
\item[Value-capture leakage; the 10--50:1 disconnect] Savings realized by AI buyers that do not flow to the cloud meter --- kept by the buyer, the integrator, or executed on-device.
\item[Cloud meter] Metered, pay-per-token or per-hour cloud billing; the channel through which infrastructure earns.
\item[Base-load inference] Routine, high-volume model usage (summaries, drafting, classification) that fills capacity between peaks.
\item[Local-silicon siphon; edge siphon] Routine inference migrating to NPUs in phones and PCs, where it is not billed to any cloud.
\item[Frontier vs routine] Bleeding-edge reasoning (a small share of volume) versus everyday tasks (the large majority) that small models handle.
\item[Economic depreciation ($\delta$)] The fall in what an installed accelerator can earn, driven by the newest chip's lower cost per token; modeled as a declining balance at 50--90\%/yr, versus straight-line accounting at 20--30\%/yr. A secondary variable.
\item[Unrecognized impairment] The gap between an asset's straight-line net book value and its economic value; what a fair-value write-down would recognize.
\item[Marginal cash cost of inference] Cash cost of producing one more token: power and facility opex per accelerator-hour divided by tokens per hour; excludes depreciation and return. $\sim$\$0.03--0.05 per million tokens.
\item[Commodity floor] The price to which a commodity converges when supply is ample: its marginal cash cost. Receipts at the floor cover opex and nothing else.
\item[Batch share ($b$)] Share of inference tokens processed asynchronously in batch windows (24 h) at 50--75\% discounts, as opposed to interactively. Raises utilization, lowers effective price; a secondary variable.
\item[Allocation vs throughput utilization] Allocation: the share of a fleet rented out (Oracle reports 97.9\%). Throughput: the share of a rented accelerator's capacity doing paid work (40--50\%). Different quantities; the batch argument concerns the second.
\item[Cost curve / supply stack] Producers ranked by full cost per unit, with capacity on the horizontal axis --- the standard commodity picture; the price settles where demand meets the curve, and producers above it lose their carry.
\item[GPU rental price] Market price to rent an accelerator by the hour (Silicon Data index); H100 \$2.63/hr in Sept.~2026. Rising rentals signal scarcity; the ratio of rental to marginal cash cost (\$0.20--0.35/hr) is the GPU-level version of the price-to-floor ratio.
\item[Hashprice] Bitcoin mining's revenue per unit of computation; converged to the marginal cost of the cheapest power within two years of each ASIC generation --- the precedent for inference pricing.
\item[Inference utilization] Share of nominal accelerator throughput actually used for paid inference (40--50\% today; latency and peak buffers take the rest).
\item[Amortization overhang; impairment wave] Depreciation obligations from past cohorts that current revenue cannot cover; the resulting write-downs when ASC 360-10 tests fail.
\item[Recoverability test] The ASC 360-10 step comparing an asset's carrying value with the undiscounted cash flows management projects for it.
\item[Wattage chokepoint; the power slot] The idea that a grid connection, not the silicon, is the scarce resource; an old accelerator is retired when a new one earns more per watt in the same slot.
\item[De-racked] Physically removed from the data center.
\item[Competing with one's own ghost] Selling retired hardware into a secondary market that then undercuts the seller's new hardware.
\item[Salvage fallacy; second-life fallacy] The assumption that retired accelerators keep meaningful value; the paper argues they do not, for power, pricing and thermal reasons.
\item[Dual-use fallacy] The assumption that AI halls can be repurposed for conventional hosting; their density and cooling make that uneconomic.
\item[Arrhenius budget] The total high-temperature operating time a part can accumulate before wear-out; consumed faster at higher temperature.
\item[Replacement vs growth capex] Spending that merely replaces retired capacity versus spending that adds capacity; at short silicon lives, a flat budget is increasingly replacement.
\item[Bullwhip] Over-ordering up a supply chain in response to demand signals, followed by oversupply; applied here to HBM.
\item[Deflationary wedge] Open-weight models driving API prices toward the marginal cost of electricity.
\item[Electricity floor] The lowest sustainable price of inference: the power to run it.
\item[Slack trap; slack absorption] Productivity gains absorbed into organizational slack rather than cash flow (see Section~\ref{sec:slack}).
\item[Indivisibility (of labor)] Fractional time savings across many people cannot be aggregated into a whole position eliminated.
\item[Verification burden] The human review that liability requires of AI output, which offsets generation savings.
\item[Manager's dilemma] Managers' incentive to report ``bandwidth expansion'' rather than surrender headcount.
\item[Avoided hiring; attrition] Cash-flow benefit realized by not replacing departures, rather than by dismissals; slower but real.
\item[Timing chasm] The mismatch between 8--12 year enterprise adoption cycles and 3--5 year silicon lives.
\item[Phase 2 (Goldman Sachs)] The 2026--27 AI-compute capex cohort of six US hyperscalers isolated in Goldman Sachs's September 2026 ROIC framework (Section~\ref{sec:gs}); the paper's cohorts $c=2026, 2027$.
\item[Take-or-pay] A power contract under which the buyer pays for the contracted energy whether or not it is used; converts power from a variable operating cost into a fixed charge (Section~\ref{sec:firstround}).
\item[Operational ramp] The share of built capacity that is in service in a given year (Goldman Sachs: 60/80/85\% in 2028--30); the paper's energization lag plays the same role. Distinct from inference utilization, which is how hard in-service capacity is run.
\item[Talent trap] The cost and failure rate of building AI capability in-house, which pushes enterprises to external operators.
\item[Secondary variable] An input too loose to pin from current data and too consequential to bury: held at a central range in the base case and swept explicitly in its own figure, so the reader sees what the result does across it. In this paper: the residual adoption factor, the assistant ceiling, the reallocation share, and seat-software growth.
\item[Active share (of paid seats)] The fraction of paid assistant seats actually used weekly; independent estimates 20--30\% in some large deployments, Microsoft claims parity with Outlook and Teams. Triang(0.25, 0.45, 0.70), a secondary variable; separates a bundle from a product.
\item[Attach rate] The fraction of a vendor's existing seats that also buy a paid add-on (here, Copilot on Microsoft 365). About 6.5\% in mid-2026 (30M paid Copilot seats on $\sim$460M commercial seats).
\item[Price elasticity of seats ($\epsilon$)] Seats lost per unit of cumulative price increase; 0.35 means a 20\% cumulative rise loses 7\% of seats.
\item[AI tax] Trade-press term for bundling AI features into existing subscriptions at a higher price whether or not the customer wanted them.
\item[SaaSpocalypse] Wall Street's name for the January--April 2026 repricing of enterprise software, $\sim$\$2T of market value, on the recognition that AI agents replace the human seats per-seat software is priced on.
\item[CLOUD Act] The 2018 US law under which US providers must produce data held anywhere in the world on a US order; the legal basis of Europe's sovereignty exit from US cloud and software.
\item[EuroStack] The European initiative for a sovereign digital stack --- cloud, operating systems, productivity software --- independent of US vendors.
\item[Outcome-priced; seat-priced; billable-unit collapse] Pricing on results; pricing per user; the erosion of per-seat and per-hour revenue models as outcomes become purchasable.
\item[Sovereign carve-out; domestic champion] Demand reserved by governments for national providers (Mistral, Falcon, Bhashini) or air-gapped enclaves.
\item[Fortress West] The scenario label for a Western bloc with strong sovereignty rules and restricted flows.
\item[Tri-bloc] The US, Europe and China as three diverging compute regimes.
\item[Cash drag; stranded capital] Capital committed to sites that cannot be energized or built, carrying cost without revenue.
\item[Energizable capex] The dollar value of capacity the grid can connect in a year.
\item[Consent bottleneck] Local political permission as the binding constraint on siting (Forbes' phrase).
\item[Community-benefit agreement] Payments or commitments a developer makes to a locality to obtain approval.
\item[Ratepayer backlash] Public opposition driven by electricity bills rising with data-center load.
\item[Winners and losers; exposure hierarchy] The paper's tables of who bears the deficit and who captures the surplus.
\item[What would change the distribution] The specific observable outcomes that would move the paper's result decisively (Section~\ref{sec:falsify}); stated as tests on a distribution, not as a yes/no refutation.
\item[Observables] Public, verifiable numbers the model must reconcile with (Table~\ref{tab:observe}).
\item[Interim marker] A 2027 checkpoint on the path to the 2028 revenue test.
\item[Provisional; [v] / [VERIFY]] A figure marked for checking against its primary source before publication.
\end{description}


\newpage
\begin{thebibliography}{99}

\bibitem{ref:epoch_capex}
Juniewicz, I. (2026). ``Hyperscaler capex has quadrupled since GPT-4's release.'' Epoch AI Data Insight, Feb.~26, 2026. Quarterly capex (cash PP\&E plus new finance-lease ROU assets) for Alphabet, Amazon, Meta, Microsoft and Oracle from SEC XBRL, Q1 2022--Q4 2025. CC-BY. \url{https://epoch.ai/data-insights/hyperscaler-capex-trend}; data: \url{https://epoch.ai/data/charts/hyperscaler-capex-trend/stacked_hyperscalers_trend.csv}. Accessed Sept.~8, 2026.

\bibitem{ref:mlq2026}
MLQ News (2026). ``Big Tech's 2026 Capex Range Reaches \$720 Billion to \$745 Billion,'' Aug.~2, 2026. Cites Amazon cash capex \$48.1B / \$77.7B / \$128.3B (2023--25), Alphabet \$32.3B / \$52.5B / \$91.4B, Meta \$27.0B / \$37.3B / \$69.7B, and 2026 guidance. \url{https://mlq.ai/news/big-techs-2026-capex-range-reaches-720-billion-to-745-billion/}

\bibitem{ref:statista2026}
Statista (2026). ``Big Tech's AI Spending to Reach \$760 Billion in 2026,'' chart, Jul.~31, 2026, from Q2 2026 earnings guidance; 2025 comparison \$413B. \url{https://www.statista.com/chart/35046/capital-expenditure-of-meta-alphabet-amazon-and-microsoft/}

\bibitem{ref:futurum2026}
Patience, N. (2026). ``AI Capex 2026: The \$690B Infrastructure Sprint,'' Futurum Group, Feb.~12, 2026. \url{https://futurumgroup.com/insights/ai-capex-2026-the-690b-infrastructure-sprint/}

\bibitem{ref:yahoo_goldman2026}
Yahoo Finance (2026). ``Meta, Microsoft, Amazon, and Alphabet are about to spend a shocking amount of money to dominate the AI era,'' Jun.~3, 2026 (Goldman Sachs \$5.3T 2025--30 capex estimate). \url{https://finance.yahoo.com/sectors/technology/article/meta-microsoft-amazon-and-alphabet-are-about-to-spend-a-shocking-amount-of-money-to-dominate-the-ai-era-115359575.html}

\bibitem{ref:yahoo2026}
Yahoo Finance (2026). ``Hyperscalers Hit \$700 Billion in 2026 AI Spending Plans,'' Q1 2026 earnings coverage (Alphabet Q1 capex \$35.67B; Google Cloud backlog $>$\$460B; Microsoft AI run-rate $>$\$37B). \url{https://finance.yahoo.com/sectors/technology/articles/hyperscalers-hit-700-billion-2026-111243744.html}\verify

\bibitem{ref:traderinsights2026}
Trader Insights (2026). ``Two Weeks That Rewired \$200 Billion in AI Capital,'' Apr.~2026 (AWS AI run-rate $>$\$15B; OpenAI \$122B round). \url{https://traderinsightsmedia.beehiiv.com/p/two-weeks-that-rewired-200-billion-in-ai-capital}\verify

\bibitem{ref:nvda_fy26}
NVIDIA Corporation (2026). ``NVIDIA Announces Financial Results for Fourth Quarter and Fiscal 2026,'' Form 8-K exhibit, Feb.~25, 2026 (FY2026 revenue \$215.9B; data-center \$193.7B). \url{https://www.sec.gov/Archives/edgar/data/1045810/000104581026000019/q4fy26pr.htm}

\bibitem{ref:nvda_q1fy27}
NVIDIA Corporation (2026). CFO Commentary on First Quarter Fiscal 2027 Results, Form 8-K exhibit (data-center revenue \$75.2B; hyperscale $\sim$50\%). \url{https://www.sec.gov/Archives/edgar/data/0001045810/000104581026000051/q1fy27cfocommentary.htm}

\bibitem{ref:anthropic_seriesg}
Anthropic (2026). ``Anthropic raises \$30 billion Series G at \$380 billion post-money valuation,'' Feb.~12, 2026 (run-rate revenue \$14B). \url{https://www.anthropic.com/news/anthropic-raises-30-billion-series-g-funding-380-billion-post-money-valuation}

\bibitem{ref:techcrunch_anthropic}
TechCrunch (2026). ``Anthropic's annualized revenue surges to \$65B,'' Aug.~17, 2026, citing Bloomberg and the Financial Times (Anthropic \$65B run-rate end-July from \$9B end-2025; \$65B round at \$965B; OpenAI \$40B). \url{https://techcrunch.com/2026/08/17/anthropics-annualized-revenue-surges-to-65b/}

\bibitem{ref:axios_anthropic}
Axios (2026). ``Anthropic's revenue run rate reportedly surpasses \$65 billion pre-IPO,'' Aug.~17, 2026 (Q2 2026 revenue $>$\$11.5B; Q1 \$4.73B). \url{https://www.axios.com/2026/08/17/anthropic-revenue-run-rate-ipo-openai}

\bibitem{ref:epoch_revenue}
Epoch AI (2026). \textit{AI companies revenue reports} dataset (run-rate and booked revenue with source quotes: OpenAI \$21.4B end-2025, ``more than \$40B'' Aug.~2026; Anthropic \$14B Feb.~2026; others). \url{https://epoch.ai/data/ai_companies_revenue_reports.csv}

\bibitem{ref:epoch_revenue_x}
Epoch AI (2026). ``Anthropic could surpass OpenAI in annualized revenue by mid-2026,'' Data Insight, Feb.~19, 2026 (growth rates; company-guided 2026 growth of $\sim$2.2$\times$ OpenAI, $\le$4$\times$ Anthropic per The Information). \url{https://epoch.ai/data-insights/anthropic-openai-revenue}

\bibitem{ref:epoch_spending}
Epoch AI (2026). ``Compute accounts for the majority of expenses of AI companies,'' Data Insight, Feb.~4, 2026. \url{https://epoch.ai/data-insights/company-spending-breakdown}

\bibitem{ref:valueadd_openai}
ValueAdd VC (2026). ``OpenAI Revenue 2026: \$40B ARR as Anthropic Hits \$65B,'' Sept.~2026 (OpenAI \$122B round at \$852B, Mar.~31, 2026; reported 2025 booked revenue \$13.07B and operating loss \$20.9B; confidential S-1 Jun.~8, 2026). Secondary source; figures attributed to leaked financials. \url{https://valueaddvc.com/blog/openai-revenue-2026-20b-arr-4b-month-path-to-profitability}\verify

\bibitem{ref:meta_q2_2026}
Adweek (2026). ``Meta's Tepid Revenue Outlook Undercuts its AI Spending Spree,'' Jul.~29, 2026 (Q2 2026 call: Advantage+ end-to-end solutions above a \$75B annual run-rate; AI ad matching yielding $>$8\% more clicks and 15.7\% more Facebook conversions). \url{https://www.adweek.com/media/metas-tepid-revenue-outlook-undercuts-its-ai-spending-spree/}

\bibitem{ref:meta_q2_2026b}
Meta Platforms, Inc. (2026). Second Quarter 2026 Results, Jul.~29, 2026 (revenue \$60.8B, +28\%; ad impressions +14\%, average price per ad +12\%); Q1 2026 results (revenue \$56.3B, +33\%). \url{https://investor.atmeta.com/}; summarized in \url{https://www.refontelearning.com/blog/meta-advantage-plus-explained} and \url{https://www.contentgrip.com/meta-ai-scaled-ad-engine-q1-2026/}

\bibitem{ref:sacra_anthropic}
Sacra (2026). ``Anthropic revenue, valuation \& funding,'' updated Jul.~8, 2026 (gross-basis recognition of reseller revenue; $\sim$\$45B run-rate May 2026 vs OpenAI $\sim$\$33B; OpenAI Q1 2026 loss $\sim$\$7B; Microsoft revenue share $\sim$\$6B/yr). \url{https://sacra.com/c/anthropic/}

\bibitem{ref:gartner_it}
Gartner (2026). \textit{Forecast: Worldwide IT Spending}, 2026 updates (see \cite{ref:gartner_it_feb2026,ref:gartner_it_jul2026}).

\bibitem{ref:gartner_it_feb2026}
Gartner (2026). ``Gartner Forecasts Worldwide IT Spending to Grow 10.8\% in 2026, Totaling \$6.15 Trillion,'' press release, Feb.~3, 2026 (software above \$1.4T; GenAI model spending growth 80.8\%; GenAI share of software up 1.8 points). \url{https://www.gartner.com/en/newsroom/press-releases/2026-02-03-gartner-forecasts-worldwide-it-spending-to-grow-10-point-8-percent-in-2026-totaling-6-point-15-trillion-dollars}

\bibitem{ref:gartner_it_jul2026}
Gartner (2026). ``Gartner Forecasts Worldwide IT Spending to Grow 14.2\% in 2026, Totaling \$6.37 Trillion,'' press release, Jul.~27, 2026; segment detail (software \$1.44T, IT services \$1.87T, data center systems \$788B) from the April update as reported by eeNews Europe. \url{https://www.gartner.com/en/newsroom/press-releases/2026-07-27-gartner-forecasts-worldwide-it-spending-to-grow-14-point-2-percent-in-2026-totaling-6-point-37-trillion}; \url{https://www.eenewseurope.com/en/gartner-projects-it-spending-at-6-31-trillion-in-2026-on-ai-surge/}

\bibitem{ref:gartner_ai_platforms}
Gartner (2026). ``Gartner Forecasts Worldwide AI Platforms and Models Market to Grow 63\% in 2026,'' press release, Jul.~20, 2026 (end-user spending \$64B in 2026 from \$39B in 2025). \url{https://www.gartner.com/en/newsroom/press-releases/2026-07-20-gartner-forecasts-worldwide-ai-platforms-and-models-market-to-grow-63-percent-in-2026}

\bibitem{ref:gartner_ai_spend}
Gartner (2026). ``Gartner Forecasts Worldwide AI Spending to Grow 47\% in 2026'' (\$2.59T, of which more than 45\% is vendor-built infrastructure), May~19, 2026. \url{https://www.gartner.com/en/newsroom/press-releases/2026-05-19-gartner-forecasts-worldwide-ai-spending-to-grow-47-percent-in-2026}

\bibitem{ref:consumer_mau}
Alphabet Q2 2026 earnings (Gemini app above 950M monthly active users) and Meta Q1 2026 earnings (Meta AI 1.2B monthly active users), as compiled in \url{https://www.instantpress.co/ai-statistics}\verify

\bibitem{ref:msft_copilot_seats}
Microsoft Corporation (2026). FY2026 Q3 earnings call, Apr.~29, 2026 (S.~Nadella: Microsoft 365 Microsoft 365 Copilot at 20 million paid enterprise seats, from 15 million at FY2026 Q2; customers with more than 50{,}000 seats quadrupled year over year; Accenture agreement for over 740{,}000 seats). Coverage: \url{https://letsdatascience.com/news/microsoft-office assistant-reaches-20-million-paid-enterprise-seats-eed5aec2}; transcript at \url{https://www.microsoft.com/en-us/investor/}\verify

\bibitem{ref:dcw_q1_2026}
Data Center Watch (2026). \textit{Q1 2026 Report: Blocked and Delayed Data Center Projects}, covering January--March 2026 (at least 75 projects, $\sim$\$130B; opposition groups more than doubled to 833 across 49 states; 300+ state bills; statewide moratoria proposed in 14 states). \url{https://www.datacenterwatch.org/q1-2026}; 2024--25 baseline report \url{https://www.datacenterwatch.org/report}

\bibitem{ref:tomshw_dcw}
Tom's Hardware (2026). ``More than 75 data center build-outs worth \$130 billion have been successfully blocked in the first three months of 2026,'' Jun.~13, 2026 (Ipsos: local opposition $\sim$70\%). \url{https://www.tomshardware.com/tech-industry/artificial-intelligence/more-than-75-data-center-build-outs-worth-usd130-billion-have-been-successfully-blocked-in-the-first-four-months-of-2026-bipartisan-opposition-mounts-nationwide-over-fears-of-soaring-power-and-water-costs}

\bibitem{ref:forbes_consent}
Szczerba, R. (2026). ``\$130 Billion In AI Data Centers Stalled. The Bottleneck Is Consent,'' \textit{Forbes}, Jul.~22, 2026. \url{https://www.forbes.com/sites/robertszczerba/2026/07/22/130-billion-in-ai-data-centers-stalled-the-bottleneck-is-consent/}

\bibitem{ref:kucoin_dcw}
Data Center Watch and Carbon Direct tracking, as reported Sept.~2026 (46 AI data-center projects, \$137--172B, stalled, blocked or withdrawn Jan.~2024--May 2026; Tract's \$14B Arizona withdrawal; QTS/Compass \$24.7B Virginia litigation; New York halt on large facilities, Jul.~2026; cancellations 6 in 2024, 25 in 2025, 20+ in early 2026). \url{https://cryptobriefing.com/us-data-center-projects-150b-stalled-cancelled/}\verify

\bibitem{ref:mit_nanda}
MIT NANDA (2025). \textit{The GenAI Divide: State of AI in Business 2025}, Jul.~2025 (about 95\% of enterprise generative-AI pilots with no measurable P\&L impact; gains concentrated in workflow-integrated deployments). \url{https://nanda.media.mit.edu/}\verify

\bibitem{ref:mckinsey_ai2025}
McKinsey \& Company (2025). \textit{The State of AI in 2025: Agents, Innovation, and Transformation}, Nov.~2025 (majority of firms using AI in at least one function; $\sim$1\% describing deployments as mature). \url{https://www.mckinsey.com/capabilities/quantumblack/our-insights/the-state-of-ai}\verify

\bibitem{ref:brynjolfsson_callcenter}
Brynjolfsson, E., Li, D., \& Raymond, L. (2025). ``Generative AI at Work,'' \textit{Quarterly Journal of Economics}, 140(2), 889--942 (14\% average productivity gain for customer-support agents; 34\% for the least experienced). \url{https://doi.org/10.1093/qje/qjae044}

\bibitem{ref:cui_copilot}
Cui, Z., Demirer, M., Jaffe, S., Musolff, L., Peng, S., \& Salz, T. (2024). ``The Effects of Generative AI on High-Skilled Work: Evidence from Three Field Experiments with Software Developers,'' SSRN 4945566 (26\% more completed tasks). \url{https://ssrn.com/abstract=4945566}

\bibitem{ref:dellacqua_bcg}
Dell'Acqua, F., McFowland, E., Mollick, E., et al. (2023). ``Navigating the Jagged Technological Frontier,'' Harvard Business School Working Paper 24-013 (consultants with GPT-4: 12.2\% more tasks, 25.1\% faster, 40\% higher quality inside the frontier). \url{https://www.hbs.edu/faculty/Pages/item.aspx?num=64700}

\bibitem{ref:saaspocalypse_ti}
Tech Insider (2026). ``AI Agents Just Erased \$2T in SaaS Value --- Who Survives,'' Aug.~2026 (JPMorgan: largest non-recessionary 12-month software drawdown in 30+ years; IGV down 20--30\% Feb.--Apr.; \$285B in 48 hours on the Claude Cowork launch). \url{https://tech-insider.org/saas-stock-crash-ai-agents-2-trillion-2026/}

\bibitem{ref:saaspocalypse_nx}
NxCode (2026). ``SaaSpocalypse 2026: Why AI Just Wiped \$285B from Software,'' Feb.~5, 2026 (Bank of America's ``internally inconsistent'' objection: investors cannot believe both that AI capex fails and that AI disrupts all software). \url{https://www.nxcode.io/resources/news/saaspocalypse-2026-software-stock-crash}

\bibitem{ref:saaspocalypse_fc}
FinancialContent / MarketMinute (2026). ``The SaaSpocalypse of 2026: How Generative AI Broke the Software Growth Engine,'' Mar.~30, 2026 (Atlassian $-35$\% on first systemic seat-count decline; Adobe $-36$\%; Salesforce $-33$\%). \url{https://markets.financialcontent.com/wral/article/marketminute-2026-3-30-the-saaspocalypse-of-2026-how-generative-ai-broke-the-software-growth-engine}

\bibitem{ref:saaspocalypse_arbi}
Arbisoft (2026). ``The SaaSpocalypse and What Survives It,'' May 2026 (Gartner: 40\% of enterprise applications with AI agents by end-2026, from under 5\% in 2025). \url{https://arbisoft.com/blogs/the-saa-spocalypse-and-what-survives-it}

\bibitem{ref:eu_sovereignty}
Emmermann, V. (2026). ``France Is Ditching Windows for Linux --- And Every EU Business Should Pay Attention,'' Apr.~2026 (DINUM announcement of Apr.~8, 2026; Schleswig-Holstein's cabinet decision of Apr.~2024; Denmark; CNAM's 80,000-agent migration). \url{https://blog.cierra.ai/blog/post/france-ditching-windows-linux-eu-digital-sovereignty?lang=en}

\bibitem{ref:tuta_sovereignty}
Tuta (2026). ``Digital Sovereignty Study: EU IT leaders are quitting U.S. Big Tech,'' summarizing \textit{State of Digital Sovereignty 2026} (1,818 IT decision-makers in Germany, France and the Nordics; 40\% of German IT leaders actively implementing European software, from 20\%; 87.5\% favoring European products in critical infrastructure; the ICC prosecutor account incident). \url{https://tuta.com/blog/digital-sovereignty-study}\verify

\bibitem{ref:m365_price_wl}
Windows Latest (2026). ``Microsoft 365 just got a price hike over continuous innovation, but Copilot is the AI tax on businesses,'' Jul.~5, 2026 (increases up to 43\% on some SKUs effective Jul.~1, 2026, announced Dec.~4, 2025; Security Copilot at 400 SCUs per 1,000 licenses bundled into E5). \url{https://www.windowslatest.com/2026/07/05/microsoft-365-just-got-a-price-hike-over-continuous-innovation-but-copilot-is-the-ai-tax-on-businesses/}

\bibitem{ref:m365_price_uscloud}
US Cloud (2026). ``Microsoft 365 Price Increase July 1, 2026: What Enterprise Buyers Need to Know,'' Jul.~2026 (E5/E7 packaging, Security Copilot capacity limits, renewal impact). \url{https://www.uscloud.com/blog/microsoft-365-price-increase-july-1-2026-what-enterprise-buyers-need-to-know-part-1/}

\bibitem{ref:m365_price_ow}
Office Watch (2026). ``Microsoft 365 Business Price Increase: Who Pays More in 2026,'' Jul.~8, 2026 (Windows Enterprise increase; ``a chunk of that money is funding Copilot for a customer base that has consistently shown little interest''). \url{https://office-watch.com/2026/microsoft-365-business-price-increase-2026/}

\bibitem{ref:mehta_2026}
Mehta, V. (2026). Reviewer correspondence on version 24 of this paper, 29 September 2026: economic versus accounting depreciation of accelerators (80--90\%/yr vs 20--30\%/yr), inference utilization modeled at 40--50\% and its possible doubling, and convergence of token prices to marginal cash cost; market context on neo-cloud short interest. Quoted with permission\verify.

\bibitem{ref:orcl_q1fy27}
Oracle Corporation (2026). Q1 fiscal 2027 results and earnings call, Sept.~10, 2026 (revenue \$19.3B, +30\%; OCI \$7.4B, +121\%; RPO \$664B, +\$26B q/q; GPU utilization 97.9\%; 850 MW delivered in the quarter; FY26 capex \$55.7B, free cash flow $-$\$23.7B; $\sim$\$40B financing planned for FY27). Transcript: \url{https://www.fool.com/earnings/call-transcripts/2026/09/11/oracle-orcl-q1-2027-earnings-call-transcript/}; summaries: \url{https://finance.yahoo.com/markets/stocks/articles/oracle-corp-orcl-q1-2027-090101356.html}

\bibitem{ref:orcl_q1fy27_mw}
MarketWise (2026). ``Oracle Earnings: Watch Out for Massive OpenAI Risk Despite Blowout First Quarter,'' Sept.~2026 (OpenAI about half of RPO; CFO: new contracts ``via prepay or bring your own hardware''; RPO not reaching revenue or capex until fiscal 2028 or beyond). \url{https://marketwise.com/investing/oracle-earnings-openai-risk/}

\bibitem{ref:sacra_openai}
Sacra (2026). ``OpenAI revenue, valuation \& funding,'' updated Jul.~29, 2026 and after (Q1 2026 revenue \$5.7B, Q2 \$6.7B; run-rate +35\% quarter-to-date in mid-August, enterprise +50\%; ads \$1B annualized; Claude 56M MAU in Q2 2026; OpenRouter monthly token volumes 25.6T vs 8.8T). \url{https://sacra.com/c/openai/}

\bibitem{ref:aibw_openai}
AI Business Weekly (2026). ``OpenAI Statistics 2026,'' Jul.~17, 2026 (ChatGPT 900M WAU Feb.~2026, 1B MAU Jun.~2026 per Reuters/Sensor Tower; 50M+ consumer subscribers; 9M+ paying business users; S-1 filed confidentially Jun.~8). \url{https://aibusinessweekly.net/p/openai-statistics}\verify

\bibitem{ref:batch_pricing}
Batch and caching price tiers: OpenAI Batch API and prompt caching; Anthropic Message Batches and prompt caching; Google Vertex AI batch prediction --- lab pricing pages, Sept.~2026 (batch 50--75\% off interactive rates, 24-hour windows; caching up to 90\% off cached prefixes). Aggregator summaries consulted are listed in the repository's REFERENCES.md and are not cited here\verify.

\bibitem{ref:alphabet_q2_2026}
Alphabet Inc. (2026). Second Quarter 2026 Results, Jul.~22, 2026 (revenue \$119.8B, +24\%; Google Cloud \$24.8B, +82\%, operating income \$8.8B; cloud backlog \$514B; capex \$44.9B; 2026 capex guidance raised to \$195--205B; free cash flow $-$\$5.9B). Coverage: \url{https://www.cnbc.com/2026/07/22/google-earnings-q2-goog-live-updates.html}; \url{https://www.metirai.com/blog/alphabet-q2-2026-earnings-google-cloud-ai-capex}

\bibitem{ref:q2_scorecard}
Digital Applied (2026). ``AI Capex Scorecard: What Earnings Week Actually Showed,'' Aug.~2026 (Microsoft FY26 capex \$115.9B, $\sim$\$175B calendar-2026, $>$\$50B Q1 FY27, Azure +43\%, Copilot paid seats past 30M; Meta FY26 capex \$130--145B; Amazon Q2 net sales \$200.6B, no full-year capex guidance). \url{https://www.digitalapplied.com/blog/ai-capex-scorecard-earnings-week-july-2026}

\bibitem{ref:futurex_q2}
FutureX Capital (2026). ``Q2 2026 Earnings: Big Tech's AI Capex Race Hits \$730 Billion,'' Aug.~3, 2026 (quarterly capex GOOG \$44.9B, MSFT \$41B, AMZN \$54.2B, META \$31.08B; Amazon 2026 cash capex $\sim$\$220B on memory prices; Meta \$135--145B; AWS +37\%). \url{https://futurex.capital/en/ai-lab/reports/big-tech-ai-capex-2026q2}

\bibitem{ref:directions_msft}
Directions on Microsoft (2026). ``Microsoft: Expect Capacity Constraints, CapEx Acceleration to Continue,'' Jul.~2026 (FY26 Q4 revenue \$90B, Azure +43\%, Azure past \$100B annual, Q4 capex \$41B, Q1 FY27 Azure guide up to 45\%). \url{https://www.directionsonmicrosoft.com/microsoft-expect-capacity-constraints-capex-acceleration-to-continue/}

\bibitem{ref:msft_8k_fy27}
Microsoft Corporation (2026). Form 8-K, Exhibit 99.1, FY2027 investor-metrics restatement (Microsoft 365 commercial cloud revenue growth 14--20\%; Microsoft 365 commercial seat growth 6--7\% year over year each quarter of FY2026; Azure growth 39--43\%). \url{https://www.sec.gov/Archives/edgar/data/0000789019/000119312526380280/d291965dex991.htm}

\bibitem{ref:ai2_copilot}
AI2.Work (2026). ``Microsoft 365 Copilot Hits 30M Paid Seats, Backlog Tops \$678B,'' Jul.~2026 (E7 bundling, July price increases and usage-based billing behind the seat jump; 6.5\% penetration of $\sim$464M commercial seats; independent estimates of 20--30\% weekly active use of assigned seats in some large deployments). \url{https://ai2.work/blog/microsoft-365-copilot-hits-30m-paid-seats-backlog-tops-678b}\verify

\bibitem{ref:windowsforum_copilot}
WindowsForum (2026). ``Microsoft 365 Copilot's 60 Million Seat Forecast: E7 Bundles, Pricing and Renewal Risks,'' Sept.~30, 2026 (seat series 15M/20M/30M; E7 general availability May 1, 2026; usage-based billing added in July; commercial seats +6\%, revenue +17\%). \url{https://windowsforum.com/news/microsoft-365-copilots-60-million-seat-forecast-e7-bundles-pricing-and-renewal-risks.446453/}

\bibitem{ref:fool_copilot}
The Motley Fool (2026). ``Prediction: Microsoft Will Double Its Paid Copilot Seats in the Next 12 Months,'' Sept.~29, 2026 (30M seats on $>$450M commercial base growing $\sim$6\%/yr; FY26 M365 commercial cloud revenue +17\%, seats +6\%; cost of revenue in the segment +12\% on Copilot infrastructure). \url{https://www.fool.com/investing/2026/09/29/prediction-microsoft-will-double-its-paid-copilot-seats-in-the-next-12-months/}

\bibitem{ref:china_q2_capex}
NextFin News via TMTPost (2026). ``China's Tech Giants Pour Billions in AI Infrastructure in Q2,'' Sept.~2026 (Alibaba Q2 revenue RMB~269B, AI cloud and compute revenue RMB~48.4B +45\%, adjusted profit $-$38\%; Tencent capex RMB~52.8B, free cash flow $-$RMB~13.8B on RMB~51.4B compute prepayments; Baidu profit $-$68\%; Chinese cloud ROIC 13--20\% vs US 25--50\%). \url{https://www.tmtpost.com/8119172.html}

\bibitem{ref:36kr_china}
36Kr (2026). ``Alibaba and Tencent capex tops RMB 120 billion in one quarter,'' Sept.~2026 (Alibaba capex RMB~67.7B +75\%, a quarter of revenue; Tencent RMB~52.8B +176\%; Baidu capex 36\% of revenue, free cash flow $-$RMB~7.95B). \url{https://eu.36kr.com/en/p/3948069327591304}

\bibitem{ref:trendforce_csp9}
TrendForce (2026). Top-nine cloud service provider capex forecast, revised May 2026 to $\sim$\$830B (+79\%); reported by Evertiq and iConnect007. \url{https://evertiq.com/news/2026-05-06-ai-boom-pushes-hyperscaler-capex-towards-usd-830-billion-in-2026}; \url{https://iconnect007.com/article/149904/north-american-ai-data-center-expansion-drives-2026-capex-of-top-nine-csps-to-us830-billion/149901/aep}

\bibitem{ref:taiwannews_csp9}
Taiwan News (2026). ``Taiwan IT makers to benefit as cloud giants' capex jumps 90\%,'' Sept.~2026 (TrendForce: nine CSPs $>$\$886.7B in 2026; Chinese four AI capex +80\%; ByteDance leading). \url{https://www.taiwannews.com.tw/en/news/6414153}

\bibitem{ref:david1990}
David, P. A. (1990). ``The Dynamo and the Computer: An Historical Perspective on the Modern Productivity Paradox,'' \textit{American Economic Review}, 80(2), 355--361. \url{https://www.jstor.org/stable/2006600}

\bibitem{ref:jorgenson2008}
Jorgenson, D. W., Ho, M. S., \& Stiroh, K. J. (2008). ``A Retrospective Look at the U.S. Productivity Growth Resurgence,'' \textit{Journal of Economic Perspectives}, 22(1), 3--24. \url{https://doi.org/10.1257/jep.22.1.3}

\bibitem{ref:brynjolfsson_hitt2000}
Brynjolfsson, E., \& Hitt, L. M. (2000). ``Beyond Computation: Information Technology, Organizational Transformation and Business Performance,'' \textit{Journal of Economic Perspectives}, 14(4), 23--48. \url{https://doi.org/10.1257/jep.14.4.23}

\bibitem{ref:bls_productivity}
U.S. Bureau of Labor Statistics. \textit{Labor Productivity and Costs}, nonfarm business sector, historical series (annual growth $\sim$2.8\% 1948--73, 1.4\% 1973--95, 2.9\% 1995--2004, 1.4\% 2005--19). \url{https://www.bls.gov/productivity/}

\bibitem{ref:census_btos}
U.S. Census Bureau. \textit{Business Trends and Outlook Survey}, AI supplement, 2025 (share of firms using AI to produce goods or services). \url{https://www.census.gov/data/experimental-data-products/business-trends-and-outlook-survey.html}\verify

\bibitem{ref:a16z_som}
Andreessen Horowitz (2026). \textit{State of Markets}, September 2026 (90 slides: hyperscaler capex consensus from S\&P Capital IQ as of Sept.~18, 2026; J.P.~Morgan capex and FCF series; BofA cycle comparison and ROIC; cloud backlogs; Silicon Data GPU rentals; OpenRouter token and spend shares; enterprise AI budget surveys; software fundamentals; Watten--Bistline--Blanford electricity study). Supplied to the author; \url{https://a16z.com/}

\bibitem{ref:gs_roic2026}
Sheridan, E., Borges, G., Schneider, J., Vegliante, A., Fein-Ashley, J., Huang, E., Valenti, C., Hager, M., Zhang, S., and Makkar, A. (2026). ``Americas Technology: Sizing the AI Economy Needed to Justify ROIC on Hyperscaler AI Capex,'' Goldman Sachs Global Investment Research, Equity Research, Sept.~24, 2026, 29 pp.\ (Phase-2 2026--27 AI capex \$1.726T; \$42.4B/GW; 70/30 compute/shell at 5/15-yr lives; \$836M/GW-yr opex; 15\% ROIC; \$1.42T 2028--30 revenue, \$11.6B/GW-yr; backlog \$1.69T; semiconductor build 20/35/57 GW). Client research, not publicly linkable. Primary sell-side source.

\bibitem{ref:ppa_press}
Constellation Energy (2024). ``Constellation to Launch Crane Clean Energy Center, Restoring Jobs and Carbon-Free Power to The Grid,'' press release, Sept.~20, 2024 (20-year power purchase agreement with Microsoft); Talen Energy (2024--25), Susquehanna--Amazon Web Services data-center power arrangements; Constellation Energy (2025). ``Constellation and Meta Sign 20-Year Agreement for Clinton Clean Energy Center,'' press release, Jun.~3, 2025. Company press releases; contract terms (take-or-pay structure, volumes) not fully disclosed\verify.

\bibitem{ref:jpmam_gtm}
J.P.~Morgan Asset Management (2026). \textit{Guide to the Markets}, August 2026 (hyperscaler capex 2019--2028E: 71, 97, 131, 159, 154, 241, 416, 791, 1,066, 1,179 \$B; hyperscaler free cash flow as \% of sales), as reproduced in~\cite{ref:a16z_som}.

\bibitem{ref:silicon_data}
Silicon Data (2026). GPU rental price index and inferred residual values, Sept.~7, 2026 (H100 \$2.63, H200 \$3.29, B200 \$5.69, A100 \$1.59 per GPU-hour), as reproduced in~\cite{ref:a16z_som}. \url{https://www.silicondata.com/}\verify

\bibitem{ref:openrouter}
OpenRouter (2026). Token and spend share by model provider, Sept.~28, 2026 (Anthropic, Google and OpenAI $\sim$69\% of spend on $\sim$35\% of tokens; tokens up $\sim$600$\times$ over two years), as reproduced in~\cite{ref:a16z_som}. \url{https://openrouter.ai/rankings}\verify

\bibitem{ref:watten2026}
Watten, S., Bistline, J., \& Blanford, G. (2026). ``Have Data Centers Raised Your Electric Bill?,'' Aug.~24, 2026 (each 10\% increase in state data-center capacity associated with a $\sim$0.4\% fall in average residential rates; 2019--24 capacity growth of 160\% associated with rates $\sim$6\% lower). As quoted in~\cite{ref:a16z_som}\verify.

\bibitem{ref:amzn_10k_life}
Amazon.com, Inc. (2026). Form 10-K for fiscal 2025, Note 1: ``Effective January 1, 2025 we changed our estimate of the useful lives of a subset of our servers and networking equipment from six years to five years. The shorter useful lives are due to the increased pace of technology development, particularly in the area of artificial intelligence and machine learning.'' Summarized with quantified effects in Hudson Labs, \url{https://hudson-labs.com/research/amazon-server-depreciation-amzn}; see also Usvyatsky, O., Deep Quarry, Feb.~11, 2025, \url{https://deepquarry.substack.com/p/amazon-revises-server-lifespan-amid}

\bibitem{ref:hudson_meta}
Hudson Labs (2025). ``Server useful life: AMZN, META'' (Meta extension to 5.5 years; FY2025 depreciation reduced \$2.92B). \url{https://www.hudson-labs.com/research/server-useful-life-amzn-meta}

\bibitem{ref:hudson_msft}
Hudson Labs (2026). ``MSFT server useful life'' (4 to 6 years effective FY2023; 2--6 year range disclosed thereafter). \url{https://www.hudson-labs.com/research/msft-server-useful-life}

\bibitem{ref:hudson_policy}
Hudson Labs. ``Accounting policy changes boost tech earnings'' (Alphabet, Microsoft, Amazon, Oracle useful-life history 2020--2023). \url{https://hudson-labs.com/blog/accounting-policy-changes-boost-tech-earnings}

\bibitem{ref:jedec_jep122}
JEDEC Solid State Technology Association. \textit{JEP122H: Failure Mechanisms and Models for Semiconductor Devices}, 2016 (Arrhenius model; default activation energies by mechanism). \url{https://www.jedec.org/standards-documents/docs/jep-122h}

\bibitem{ref:arrhenius_rule}
Industry rule of thumb of failure-rate doubling per 10$^{\circ}$C, derived from the Arrhenius relation at $E_a\approx0.7$\,eV over 50--100$^{\circ}$C; discussed in ASHRAE TC~9.9, \textit{Thermal Guidelines for Data Processing Environments}, 5th ed.\ (2021), and in vendor fresh-air-cooling studies (Dell, 2012)\verify.

\bibitem{ref:black1969}
Black, J. R. (1969). ``Electromigration --- A Brief Survey and Some Recent Results,'' \textit{IEEE Transactions on Electron Devices}, 16(4), 338--347. \url{https://doi.org/10.1109/T-ED.1969.16754}

\bibitem{ref:llama3}
Grattafiori, A., et al. (Meta, 2024). ``The Llama 3 Herd of Models,'' Section 3.3.4 (Reliability and operational challenges: 419 unexpected interruptions over a 54-day snapshot on 16K H100s; 58.7\% attributed to GPU issues). arXiv:2407.21783. \url{https://arxiv.org/abs/2407.21783}

\bibitem{ref:fasb_asc360}
Financial Accounting Standards Board. \textit{ASC 360-10, Property, Plant, and Equipment --- Overall}, Subsections 35-17 to 35-35 (recoverability test on undiscounted cash flows; impairment measured at fair value). \url{https://asc.fasb.org/}

\bibitem{ref:damodaran2026}
Damodaran, A. (2026). ``Cost of Capital by Sector (US),'' January 2026 data update, Stern School of Business, NYU. \url{https://pages.stern.nyu.edu/~adamodar/New_Home_Page/datafile/wacc.html}

\bibitem{ref:koller2020}
Koller, T., Goedhart, M., \& Wessels, D. (2020). \textit{Valuation: Measuring and Managing the Value of Companies}, 7th ed. Wiley. ISBN 978-1-119-61088-5.

\bibitem{ref:brealey2023}
Brealey, R. A., Myers, S. C., \& Allen, F. (2023). \textit{Principles of Corporate Finance}, 14th ed. McGraw-Hill. ISBN 978-1-265-07415-9.

\bibitem{ref:bls_oes}
U.S. Bureau of Labor Statistics. \textit{Occupational Employment and Wage Statistics}, May 2025. \url{https://www.bls.gov/oes/}\verify

\bibitem{ref:bls_ecec}
U.S. Bureau of Labor Statistics. \textit{Employer Costs for Employee Compensation}, 2025. \url{https://www.bls.gov/ecec/}\verify

\bibitem{ref:eia_epm}
U.S. Energy Information Administration. \textit{Electric Power Monthly}, Table 5.6.A, average industrial price by state, 2024--2026. \url{https://www.eia.gov/electricity/monthly/}

\bibitem{ref:pjm_queue}
PJM Interconnection. \textit{Interconnection Queue} and \textit{2025 Long-Term Load Forecast Report}. \url{https://www.pjm.com/planning/service-requests/interconnection-queues}\verify

\bibitem{ref:eirgrid_cru}
Commission for Regulation of Utilities (Ireland). \textit{CRU/21/124, Direction to the System Operators related to Data Centre grid connection processing}, Nov.~2021; EirGrid, \textit{Shaping Our Electricity Future}. \url{https://www.cru.ie/}\verify

\bibitem{ref:amsterdam}
Gemeente Amsterdam. \textit{Vestigingsbeleid datacenters} (data-centre establishment policy, 2020) and subsequent grid-congestion notices from TenneT and Liander. \url{https://www.amsterdam.nl/}\verify

\bibitem{ref:gdpr}
Regulation (EU) 2016/679 (General Data Protection Regulation), Chapter V, Articles 44--50. \url{https://eur-lex.europa.eu/eli/reg/2016/679/oj}

\bibitem{ref:aiact}
Regulation (EU) 2024/1689 (Artificial Intelligence Act). \url{https://eur-lex.europa.eu/eli/reg/2024/1689/oj}

\bibitem{ref:dpdp}
Government of India. \textit{Digital Personal Data Protection Act, 2023}, No.~22 of 2023. \url{https://www.meity.gov.in/}

\bibitem{ref:coreweave_s1}
CoreWeave, Inc. (2025). Form S-1 Registration Statement and subsequent Forms 10-Q (2025--2026), debt facilities and NVIDIA capacity arrangements. \url{https://www.sec.gov/edgar/browse/?CIK=1769628}

\bibitem{ref:oracle_10q}
Oracle Corporation. Form 10-Q, Q1 FY2026, remaining performance obligations. \url{https://www.sec.gov/edgar/browse/?CIK=1341439}

\bibitem{ref:cbinsights}
CB Insights and PitchBook (2025--2026). AI venture funding and startup outcome datasets; cohort failure tracking for 2024--25 funded generative-AI companies. \url{https://www.cbinsights.com/research/}\verify

\bibitem{ref:octoverse}
GitHub (2025). \textit{Octoverse 2025: The state of open source and AI}. \url{https://github.blog/news-insights/octoverse/}

\bibitem{ref:ifr}
International Federation of Robotics (2025). \textit{World Robotics 2025 --- Industrial Robots}. \url{https://ifr.org/worldrobotics/}

\bibitem{ref:av_fbi}
Fortune Business Insights (2026). \textit{Autonomous Driving Software Market Size, Share \& Industry Analysis, 2026--2034}. \url{https://www.fortunebusinessinsights.com/}\verify

\bibitem{ref:healthcare_rm}
Research and Markets (2026). \textit{AI-Powered Clinical Documentation Market: Global Report}. \url{https://www.researchandmarkets.com/}\verify

\bibitem{ref:newzoo}
Newzoo (2026). \textit{Global Games Market Report 2026}. \url{https://newzoo.com/resources/trend-reports}\verify

\bibitem{ref:cahn2024}
Cahn, D. (2024). ``AI's \$600B Question,'' Sequoia Capital, Jun.~20, 2024. \url{https://www.sequoiacap.com/article/ais-600b-question/}

\bibitem{ref:covello2024}
Goldman Sachs Global Investment Research (2024). ``Gen AI: Too Much Spend, Too Little Benefit?,'' \textit{Top of Mind}, Issue 129, Jun.~25, 2024 (interviews with J.~Covello and D.~Acemoglu). \url{https://www.goldmansachs.com/insights/top-of-mind/gen-ai-too-much-spend-too-little-benefit}

\bibitem{ref:acemoglu2024}
Acemoglu, D. (2024). ``The Simple Macroeconomics of AI,'' NBER Working Paper 32487. \url{https://www.nber.org/papers/w32487}

\bibitem{ref:brynjolfsson2021}
Brynjolfsson, E., Rock, D., \& Syverson, C. (2021). ``The Productivity J-Curve: How Intangibles Complement General Purpose Technologies,'' \textit{American Economic Journal: Macroeconomics}, 13(1), 333--372. \url{https://doi.org/10.1257/mac.20180386}

\bibitem{ref:nordhaus2004}
Nordhaus, W. D. (2004). ``Schumpeterian Profits in the American Economy: Theory and Measurement,'' NBER Working Paper 10433. \url{https://www.nber.org/papers/w10433}

\bibitem{ref:solow1987}
Solow, R. M. (1987). ``We'd Better Watch Out,'' review of \textit{Manufacturing Matters}, \textit{New York Times Book Review}, Jul.~12, 1987, p.~36.

\bibitem{ref:epoch_data2024}
Villalobos, P., Ho, A., Sevilla, J., Besiroglu, T., Heim, L., \& Hobbhahn, M. (2024). ``Will We Run Out of Data? Limits of LLM Scaling Based on Human-Generated Data,'' Epoch AI; arXiv:2211.04325v2. \url{https://arxiv.org/abs/2211.04325}

\bibitem{ref:fermi_trinity}
Fermi, E. (1945). ``My Observations During the Explosion at Trinity on July 16, 1945,'' Los Alamos National Laboratory archive. Reproduced at \url{https://www.atomicarchive.com/resources/documents/trinity/fermi.html}

\bibitem{ref:katz2021}
Katz, J. I. (2021). ``Fermi at Trinity,'' \textit{Nuclear Technology}, 207(sup1), S326--S329; arXiv:2103.05784. \url{https://doi.org/10.1080/00295450.2021.1906623}

\bibitem{ref:morrison1963}
Morrison, P. (1963). ``Fermi Questions,'' \textit{American Journal of Physics}, 31(8), 626--627. \url{https://doi.org/10.1119/1.1969701}

\bibitem{ref:vonbaeyer1993}
von Baeyer, H. C. (1993). \textit{The Fermi Solution: Essays on Science}. Random House. ISBN 978-0-679-42020-3.

\end{thebibliography}
\end{document}
